% ----------------------
\subsection{1 Nephi}
% ----------------------
	\label{1NEPHI}
	
	The Book of Nephi,
	
	His Reign and Ministry
	
	An account of Lehi and his wife Sariah and his four sons,
	being called, beginning at the eldest,
	Laman, Lemuel, Sam, and Nephi.
	The Lord warns Lehi to depart out of the land of Jerusalem
	because he prophesieth unto the people concerning their iniquity
	and they seek to destroy his life.
	He taketh three days' journey into the wilderness with his family.
	Nephi taketh his brethren and returns to the land of Jerusalem
	after the record of the Jews.
	The account of their sufferings.
	They take the daughters of Ishmael to wife.
	They take their families and depart into the wilderness.
	Their sufferings and afflictions in the wilderness.
	The course of their travels.
	They come to the large waters.
	Nephi's brethren rebelleth against him.
	He confoundeth them and buildeth a ship.
	They call the name of the place Bountiful.
	They cross the large waters into the promised land etc.
	This is according to the account of Nephi,
	or in other words, I Nephi wrote this record.

	% -----
	\subsubsection{1 Nephi 1} % *1NEPHI-1
	% -----
		\label{1NEPHI-1}
	
		% --
		\paragraph{1 Nephi 1:1} % *1NEPHI-1-1 
		% ---
			\label{1NEPHI-1-1}

			I Nephi having been born of
				\footnote{%
					% \label{}%
					JS uses this phrase to describe his own parents in his \hyperlink{EMS-1832-JS-CIR-SUM-1832-c}{1832 history}.
					}%
			goodly parents,%
				\footnote{%
					% \label{1NEPHI-1-1-NOTE-goodly}%
					See the definitions in the 1828 \hyperref[GOODLY-1828]{here}, and the OED \hyperref[GOODLY-OED]{here}. The 1828 isn't much help here, and actually might be a bit misleading, because it would make you think that Nephi was referring to how attractive or enjoyable to be around his parents were. But the OED does more---the sense employed here is \#4a, ``Of a good quality or high standard; splendid, excellent, fine.'' The word is used again in \hyperref[MOSIAH-18-7]{Mosiah 18:7}, there referring to a quantity of people, which is OED definition \#5.
					}
			therefore I was taught \hyperref[SOMEWHAT]{somewhat} in all the learning of my father.
			And having seen many afflictions in the course of my days,
			nevertheless having been highly favored%
			\footnote{%
				% \label{}%
				Compare \hyperref[1NEPHI-17-35]{1 Nephi 17:35}. The phrase ``highly favored'' appears eight times in the scriptures. Five of those uses describe a group of people; three describe individual people, and the three are Nephi (here), Mary the mother of Christ (\hyperref[LUKE-1-28]{Luke 1:28}), and the brother of Jared (\hyperref[ETHER-1-34]{Ether 1:34}).
				}
			of the Lord in all my days,
			yea, having had a great knowledge of the goodness and the mysteries of God,
			therefore I make a record of my proceedings in my days.
			
		% ---
		\paragraph{1 Nephi 1:2} % *1NEPHI-1-2 
		% ---
			\label{1NEPHI-1-2}

			Yea, I make a record in the language%
				\footnote{%
					% \label{}%
					Nephi is using a broad conception of the term ``language'' here, since it includes not only a form of written communication but also a system of ``learning.'' Dewey has something like this understanding of language in mind, for example, when he writes, in LTI, ``In this further discussion, language is taken in its widest sense, a sense wider than oral and written speech. It includes the latter. But it includes also not only gestures but rites, ceremonies, monuments and the products of industrial and fine arts. A tool or machine, for example, is not simply a simple or complex physical object having its own physical properties and effects, but is also a mode of language. For it says something, to those who understand it, about operations of use and their consequences. To the members of a primitive community a loom operated by steam or electricity says nothing. It is composed in a foreign language, and so with most of the mechanical devices of modern civilization. In the present cultural setting, these objects are so intimately bound up with interests, occupations and purposes that they have an eloquent voice'' (5868--5869).
					}
			of my father,
			which consists of the learning%
				\footnote{%
					% \label{}%
					See \hyperref[MORMON-1-2]{Mormon 1:2}, 
					}
			of the Jews and the language of the Egyptians.%
				\footnote{%
					% \label{1NEPHI-1-2-NOTE-language}%
					An extremely important verse in thinking about revelation. Here Nephi tells us that, for him, his ``language'' comprises both a system of learning (or a conceptual scheme, to put it in philosophical terms) and an actual language. This is fundamental to the theory of revelation, because later Nephi will also say that ``the Lord God giveth light unto the understanding; for he speaketh unto men according to their language, unto their understanding'' (\hyperref[2NEPHI-31-3]{2 Nephi 31:3}). If God ``speaks according to'' the language that people use, and in that term ``language'' Nephi includes their conceptual scheme, then God is going to be restricted by our concepts in communication with us. See also \hyperref[1NEPHI-3-19]{1 Nephi 3:19--20}, \hyperref[JACOB-7-4]{Jacob 7:4}, \hyperref[ENOS-1-1]{Enos 1:1}, \hyperref[OMNI-1-17]{Omni 1:17--18}, \hyperref[OMNI-1-22]{Omni 1:22}, \hyperref[MOSIAH-1-2]{Mosiah 1:2} and \hyperref[MOSIAH-1-4]{1:4}, \hyperref[MOSIAH-9-1]{Mosiah 9:1},  \hyperref[MOSIAH-28-17]{Mosiah 28:17}, \hyperref[ETHER-12-39]{Ether 12:39}, \hyperref[DC90-11]{DC 90:11, 15} for a distinction between ``language'' and ``tongue''; 
					}

		% ---
		\paragraph{1 Nephi 1:3} % *1NEPHI-1-3 
		% ---
			\label{1NEPHI-1-3}

			And I know that the record which I make to be true.
			And I make it with mine own hand,
				\footnote{%
					% \label{}%
					He's openly telling us that this record is based on his ``knowledge,'' or what he recollects and what he has learned from what happened.
					}%
			and I make it according to my knowledge.

		% ---
		\paragraph{1 Nephi 1:4} % *1NEPHI-1-4 
		% ---
			\label{1NEPHI-1-4}

			For it came to pass in the commencement of the first year
			of the reign of Zedekiah,%
			\footnote{%
				% \label{}%
				For Zedekiah, see AB 9636 (person \#4), 
				}
			king of Judah
			---my father Lehi having dwelt at Jerusalem in all his days---
			and in that same year there came many prophets prophesying unto the people
			that they must repent or the great city Jerusalem must be destroyed.

		% ---
		\paragraph{1 Nephi 1:5} % *1NEPHI-1-5 
		% ---
			\label{1NEPHI-1-5}

			Wherefore it came to pass that
			my father Lehi, as he went forth, prayed unto the Lord,
			yea, even with all his heart,%
				\footnote{%
					% \label{}%
					Compare \hyperref[MORONI-7-48]{Moroni 7:48}, where Mormon instructs us to ``pray unto the Father with all the energy of heart that ye may be filled with this love which he hath bestowed upon all who are true followers of his Son Jesus Christ.'' In \hyperref[1NEPHI-1-15]{verse 15} the promise is fulfilled.
					}
			in behalf of his people.

		% ---
		\paragraph{1 Nephi 1:6} % *1NEPHI-1-6 
		% ---
			\label{1NEPHI-1-6}

			And it came to pass as he prayed unto the Lord,
			there came a pillar of fire and dwelt upon a rock before him,%
				\footnote{%
					% \label{}%
					The first spiritual experience Lehi has after praying ``in behalf of his people'' is fire coming and dwelling upon a rock. Either from out of the fire, or from a different source, he hears and sees many things, and they scare him. This experience is similar to the one Moses has in \hyperref[EXODUS-3]{Exodus 3}. Moses is watching the flocks and takes them a ways away when an angel appears in a bush. The bush burns, but isn't consumed. I assume that Lehi was familiar with this story, so he would probably have recognized the fire as a manifestation of God or the Spirit, and he would have known to pay attention. We know that he was also away from home, perhaps working, because he prays ``as he went forth'' and after the experience he must ``return to his own house at Jerusalem.'' The point here is that God spoke to Lehi in circumstances which were similar to a story he would have been familiar with, maybe to make it easier for Lehi to realize he was having a spiritual experience and easier for him to pay attention. The Lord worked with circumstances and symbols that he was already familiar with. He couldn't grow unless he recognized that he was experiencing and opportunity for growth.
					}
			and he saw and heard much.
			And because of the things which he saw and heard,
			he did quake and tremble%
			\footnote{%
				% \label{}%
				The words ``quake'' and ``tremble'' appear together as a pair five times in the Book of Mormon---here, then again at \hyperref[1NEPHI-22-23]{1 Nephi 22:23}, \hyperref[MOSIAH-27-31]{Mosiah 27:31}, \hyperref[MOSIAH-28-3]{Mosiah 28:3}, and \hyperref[HELAMAN-12-9]{Helaman 12:9}. The words are used of ``the hills and the mountains'' in Helaman, which ``tremble and quake'' in response to God's voice; the other three uses all involve punishment for unrighteousness. I think it's unlikely that Lehi had a vision hear that convicted him of his own sins and made him mourn for the state of his own soul. This would follow a pattern we see with other people in the Book of Mormon, including Enos and Alma. It would also help explain why he leaves the area around the rock to go home to Jerusalem and then has another vision there; he could have been dwelling on his sins along the way, or could have been thinking about mercy and forgiveness in response to his guilt.
				}
			exceedingly.

		% ---
		\paragraph{1 Nephi 1:7} % *1NEPHI-1-7 
		% ---
			\label{1NEPHI-1-7}

			And it came to pass that he returned to his own house at Jerusalem.
			And he cast himself upon his bed,
			being overcome with the Spirit and the things which he had seen.

		% ---
		\paragraph{1 Nephi 1:8} % *1NEPHI-1-8 
		% ---
			\label{1NEPHI-1-8}

			And being thus overcome with the Spirit,
			he was carried away in a vision,
			even that he saw the heavens open
			and he thought he saw God sitting upon his throne,
			surrounded with numberless concourses%
				\footnote{%
					% \label{}%
					See \hyperref[1NEPHI-8-21]{1 Nephi 8:21}, ``numberless concourses of people.''
					}
			of angels
			in the attitude of singing and praising their God.

		% ---
		\paragraph{1 Nephi 1:9} % *1NEPHI-1-9 
		% ---
			\label{1NEPHI-1-9}

			And it came to pass that he saw one%
			\footnote{%
				\label{1NEPHI-1-9-NOTE-one}%
				In the current edition of the BOM this word ``one'' is capitalized. That makes it seem like the ``One'' who comes down is extra special, but you don't get the capital letter in the original text. Based on verses 9 and 10 the ``one'' could be Christ, but when you don't have the capital letter, it makes it seem much more like the ``one'' is just one of the angels mentioned at the end of verse 8.
				}
			descending out of the midst of heaven,
			and he beheld that his luster%
				\footnote{%
					% \label{1NEPHI-1-9-NOTE-luster}%
					This word occurs twice in the scriptures---both in the BOM, here and at \hyperref[MOSIAH-13-5]{Mosiah 13:5}. See \hyperref[LUSTER]{here} for the definitions; it means ``brightness'' or ``splendor.''
					}
			was above that of the sun at noonday.

		% ---
		\paragraph{1 Nephi 1:10} % *1NEPHI-1-10 
		% ---
			\label{1NEPHI-1-10}

			And he also saw twelve others following him,
			and their brightness did exceed that of the stars in the firmament.%
				\footnote{%
					% \label{1NEPHI-1-10-NOTE-stars}%
					Clearly there's the imagery of the kindgoms of glory here, with the sun being mentioned in verse 9, and here the stars. It doesn't quite align with what JS later taught about the kingdoms though; could just refer to the amount of glory and power possessed by the individuals.
					}

		% ---
		\paragraph{1 Nephi 1:11} % *1NEPHI-1-11 
		% ---
			\label{1NEPHI-1-11}

			And they came down and went forth upon the face of the earth.
			And the first%
				\footnote{%
					% \label{}%
					This person appears to be the first \textit{of the twelve}, and not the ``one'' that was shining brightly. But it could also be that ``one.''
					}
			came and stood before my father
			and gave unto him a book and bade him that he should read.

		% ---
		\paragraph{1 Nephi 1:12} % *1NEPHI-1-12 
		% ---
			\label{1NEPHI-1-12}

			And it came to pass that as he read,
			he was filled with the Spirit of the Lord.

		% ---
		\paragraph{1 Nephi 1:13} % *1NEPHI-1-13 
		% ---
			\label{1NEPHI-1-13}

			And he read, saying: Woe woe unto Jerusalem,
			for I have seen thine abominations.
			Yea, and many things did my father read concerning Jerusalem,
			that it should be destroyed and the inhabitants thereof;
			many should perish by the sword
			and many should be carried away captive into Babylon.

		% ---
		\paragraph{1 Nephi 1:14} % *1NEPHI-1-14 
		% ---
			\label{1NEPHI-1-14}

			And it came to pass that when my father had read
			and saw many great and marvelous things,
			he did exclaim many things unto the Lord, such as:
			Great and marvelous are thy works, O Lord God Almighty.
			Thy throne is high in the heavens,
			and thy power and goodness and mercy is over all the inhabitants of the earth.
			And because thou art merciful,
			thou wilt not suffer those who come unto thee that they shall perish.

		% ---
		\paragraph{1 Nephi 1:15} % *1NEPHI-1-15 
		% ---
			\label{1NEPHI-1-15}

			And after this manner was the language of my father in the praising of his God,
			for his soul did rejoice and his whole heart was filled
			because of the things which he had seen,
			yea, which the Lord had shewn unto him.

		% ---
		\paragraph{1 Nephi 1:16} % *1NEPHI-1-16 
		% ---
			\label{1NEPHI-1-16}

			And now I Nephi do not make a full account of the things
			which my father hath written,
			for he hath written many things which he saw in visions and in dreams.
			And he also hath written many things
			which he prophesied and spake unto his children,
			of which I shall not make a full account.

		% ---
		\paragraph{1 Nephi 1:17} % *1NEPHI-1-17 
		% ---
			\label{1NEPHI-1-17}

			But I shall make an account of my proceedings in my days.
			Behold, I make an abridgment of the record of my father
			upon plates which I have made with mine own hands.
			Wherefore after that I have abridged the record of my father,
			then will I make an account of mine own life.

		% ---
		\paragraph{1 Nephi 1:18} % *1NEPHI-1-18 
		% ---
			\label{1NEPHI-1-18}

			Therefore I would that ye should know that
			after the Lord had shewn so many marvelous things unto my father Lehi,
			yea, concerning the destruction of Jerusalem,
			behold, he went forth among the people and began to prophesy
			and to declare unto them concerning the things
			which he had both seen and heard.

		% ---
		\paragraph{1 Nephi 1:19} % *1NEPHI-1-19 
		% ---
			\label{1NEPHI-1-19}

			And it came to pass that the Jews did mock him
			because of the things which he testified of them,
			for he truly testified of their wickedness and their abominations.
			And he testified that the things which he saw and heard,
			and also the things which he read in the book,
			manifested plainly of the coming of a Messiah
			and also the redemption of the world.

		% ---
		\paragraph{1 Nephi 1:20} % *1NEPHI-1-20 
		% ---
			\label{1NEPHI-1-20}

			And when the Jews heard these things, they were angry with him,
			yea, even as with the prophets of old,
			whom they had cast out and stoned and slain.
			And they also sought his life that they might take it away.
			But behold, I Nephi will shew unto you
			that the tender mercies%
				\footnote{%
					% \label{}%
					See \hyperref[LUKE-1-78]{Luke 1:78};
					}
			of the Lord is over all them
			whom he hath chosen because of their faith
			to make them mighty, even unto the power of deliverance.
			
	% -----
	\subsubsection{1 Nephi 2} % *1NEPHI-2
	% -----
		\label{1NEPHI-2}
		
		% ---
		\paragraph{1 Nephi 2:1} % *1NEPHI-2-1 
		% ---
			\label{1NEPHI-2-1}

			For behold, it came to pass that the Lord spake unto my father,
			yea, even in a dream,%
				\footnote{%
					% \label{}%
					Notice the great variety of ways that the Lord communicates information to Lehi during the events recorded in the book.
					}
			and saith unto him:
			Blessed art thou Lehi because of the things which thou hast done.
			And because thou hast been faithful
			and declared unto this people the things which I commanded thee,
			behold, they seek to take away thy life.

		% ---
		\paragraph{1 Nephi 2:2} % *1NEPHI-2-2 
		% ---
			\label{1NEPHI-2-2}

			And it came to pass that the Lord commanded my father, even in a dream,
			that he should take his family and depart into the wilderness.

		% ---
		\paragraph{1 Nephi 2:3} % *1NEPHI-2-3 
		% ---
			\label{1NEPHI-2-3}

			And it came to pass that he was obedient unto the word of the Lord;
			wherefore he did as the Lord commanded him.

		% ---
		\paragraph{1 Nephi 2:4} % *1NEPHI-2-4 
		% ---
			\label{1NEPHI-2-4}

			And it came to pass that he departed into the wilderness.
			And he left his house and the land of his inheritance
			and his gold and his silver and his precious things
			and took nothing with him save it were his family and provisions and tents,
			and he departed into the wilderness.

		% ---
		\paragraph{1 Nephi 2:5} % *1NEPHI-2-5 
		% ---
			\label{1NEPHI-2-5}

			And he came down by the borders near the shore of the Red Sea,
			and he traveled in the wilderness in the borders which was nearer the Red Sea.
			And he did travel in the wilderness with his family,
			which consisted of my mother Sariah and my elder brethren,
			which were Laman, Lemuel, and Sam.

		% ---
		\paragraph{1 Nephi 2:6} % *1NEPHI-2-6 
		% ---
			\label{1NEPHI-2-6}

			And it came to pass that when he had traveled three days in the wilderness,
			he pitched his tent in a valley beside a river of water.

		% ---
		\paragraph{1 Nephi 2:7} % *1NEPHI-2-7 
		% ---
			\label{1NEPHI-2-7}

			And it came to pass that he built an altar of stones,
			and he made an offering unto the Lord and gave thanks unto the Lord our God.

		% ---
		\paragraph{1 Nephi 2:8} % *1NEPHI-2-8 
		% ---
			\label{1NEPHI-2-8}

			And it came to pass that he called the name of the river Laman,
			and it emptied into the Red Sea.
			And the valley was in the borders near the mouth thereof.

		% ---
		\paragraph{1 Nephi 2:9} % *1NEPHI-2-9 
		% ---
			\label{1NEPHI-2-9}

			And when my father saw that the waters of the river
			emptied into the \hyperref[FOUNTAIN]{fountain} of the Red Sea,
			he spake unto Laman, saying:
			O that thou mightest be like unto this river,
			continually running into the \hyperref[FOUNTAIN]{fountain} of all righteousness.%
				\footnote{%
					\label{1NEPHI-2-9-NOTE-fountain}%
					The desire Lehi expresses here is a little strange. The river feeds into the Red Sea, and Lehi says he wants Laman to be like the river, ``running into the fountain of all righteousness.'' See the definitions of ``fountain'' \hyperref[FOUNTAIN-1828]{here}; the most applicable definition seems to be \#1, ``a spring, or source of water.'' If you take God or Christ as the ``fountain of all righteousness'' then it seems like Lehi is saying he hopes Laman will contribute to God or Christ in some way; but another way of reading it is to take the ``fountain'' bit as referring to the collection of all righteous actions. If you see it that way, the Lehi is saying he hopes that Laman will contribute to those righteous actions ``continually''---that Laman will lead a righteous life and combine his righteous efforts with those of other people. Still others would then be able to draw and benefit from his righteous actions.
					}

		% ---
		\paragraph{1 Nephi 2:10} % *1NEPHI-2-10 
		% ---
			\label{1NEPHI-2-10}

			And he also spake unto Lemuel, saying:
			O that thou mightest be like unto this valley,
			firm and steadfast and immovable in keeping the commandments of the Lord.

		% ---
		\paragraph{1 Nephi 2:11} % *1NEPHI-2-11 
		% ---
			\label{1NEPHI-2-11}

			Now this he spake because of the stiffneckedness of Laman and Lemuel.
			For behold, they did murmur in many things against their father
			because that he was a visionary man
			and that he had led them out of the land of Jerusalem,
			to leave the land of their inheritance
			and their gold and their silver and their precious things,
			and to perish in the wilderness.
			And this they said that he had done
			because of the foolish imaginations of his
				\footnote{%
					\label{1NEPHI-2-11-foolish}%
					See \hyperref[1NEPHI-17-20]{1 Nephi 17:20}, and \hyperref[HEART-TDOT]{this discussion} of \he{leb} in the OT.
					}%
			heart.

		% ---
		\paragraph{1 Nephi 2:12} % *1NEPHI-2-12 
		% ---
			\label{1NEPHI-2-12}

			And thus Laman and Lemuel, being the eldest,
			did murmur against their father.
			And they did murmur because they knew not
			the dealings of that God who had created them.

		% ---
		\paragraph{1 Nephi 2:13} % *1NEPHI-2-13 
		% ---
			\label{1NEPHI-2-13}

			Neither did they believe that Jerusalem, that great city,
			could be destroyed according to the words of the prophets.
			And they were like unto the Jews which were at Jerusalem,
			which sought to take away the life of my father.

		% ---
		\paragraph{1 Nephi 2:14} % *1NEPHI-2-14 
		% ---
			\label{1NEPHI-2-14}

			And it came to pass that my father did speak unto them
			in the valley of Lemuel with power, being filled with the Spirit,
			until their frames did shake before him.
			And he did confound them that they durst not utter against him;
			wherefore they did do as he commanded them.

		% ---
		\paragraph{1 Nephi 2:15} % *1NEPHI-2-15 
		% ---
			\label{1NEPHI-2-15}

			And my father dwelt in a tent.

		% ---
		\paragraph{1 Nephi 2:16} % *1NEPHI-2-16 
		% ---
			\label{1NEPHI-2-16}

			And it came to pass that I Nephi being exceeding young,
			nevertheless being large in stature,
			and also having great desires to know of the \hyperref[MYSTERY]{mysteries} of God,
			wherefore I cried unto the Lord.
			And behold, he did visit me and did soften my heart
			that I did believe all the words which had been spoken by my father;
			wherefore I did not rebel against him like unto my brothers.

		% ---
		\paragraph{1 Nephi 2:17} % *1NEPHI-2-17 
		% ---
			\label{1NEPHI-2-17}

			And I spake unto Sam, making known unto him the things
			which the Lord had manifested unto me by his Holy Spirit.
			And it came to pass that he believed in my words.

		% ---
		\paragraph{1 Nephi 2:18} % *1NEPHI-2-18 
		% ---
			\label{1NEPHI-2-18}

			But behold, Laman and Lemuel would not hearken unto my words.
			And being grieved%
				\footnote{%
					% \label{}%
					The 1828 dictionary gives this as a passive participle and says it means, ``Pained; afflicted; suffering sorrow.'' According to the OED, the only relevant meaning in use at the time of the BOM is meaning \#4, ``Affected with grief; vexed, afflicted, troubled or distressed in mind.'' See \hyperref[1NEPHI-7-8]{1 Nephi 7:8} and \hyperref[1NEPHI-15-4]{1 Nephi 15:4}, where Nephi again feels ``grieved'' for the same reason.
					}
			because of the hardness of their hearts,
			I cried unto the Lord for them.

		% ---
		\paragraph{1 Nephi 2:19} % *1NEPHI-2-19 
		% ---
			\label{1NEPHI-2-19}

			And it came to pass that the Lord spake unto me, saying:
			Blessed art thou Nephi because of thy faith,
			for thou hast sought me diligently with \hyperref[LOWLINESS-1828]{lowliness} of heart.%
				\footnote{%
					\label{1NEPHI-2-19-NOTE-lowliness}%
					Very important how Nephi does it---he seeks instruction from the Lord with ``lowliness of heart,'' or without pride. See \hyperref[LOWLINESS-1828]{here} for the 1828 definitions of ``lowliness.''
					}

		% ---
		\paragraph{1 Nephi 2:20} % *1NEPHI-2-20 
		% ---
			\label{1NEPHI-2-20}

			And inasmuch as ye shall keep my commandments,
			ye shall prosper and shall be led to a land of promise,
			yea, even a land which I have prepared for you,%
				\footnote{%
					% \label{}%
					I'm interested in how their understanding and perception of this idea changed over time. Just to note a couple ideas briefly---Jacob later says that there is a happiness ``which is prepared'' for the saints (\hyperref[2NEPHI-9-43]{2 Nephi 9:43}). But later the same Jacob talks about how lonely and unhappy everybody was, at least during certain times of their lives; \hyperref[JACOB-7-26]{Jacob 7:26} is the best example of this, especially at the end, where he says that the people ``did mourn out [their] days.'' So I'm wondering if, in particular, Jacob's view of the ``land of promise'' changed over time: he thought it was actually the land they started to live on after traveling in the boat, but as the people became less happy, Jacob may have come to think of that ``promised land'' which was prepared as more of a metaphor for the kind of blessed state that the Lord would give to the righteous and the saints. I haven't checked in other places for this yet, but I'm struck by \hyperref[JACOB-7-26]{Jacob 7:26} in particular. There are more things going on, though, as \hyperref[ JACOB-7-23]{Jacob 7:23} makes things sound a little better.
					}
			a land which is choice above all other lands.

		% ---
		\paragraph{1 Nephi 2:21} % *1NEPHI-2-21 
		% ---
			\label{1NEPHI-2-21}

			And inasmuch as thy brethren shall rebel against thee,
			they shall be cut off from the presence of the Lord.

		% ---
		\paragraph{1 Nephi 2:22} % *1NEPHI-2-22 
		% ---
			\label{1NEPHI-2-22}

			And inasmuch as thou shalt keep my commandments,
			thou shalt be made a ruler and a teacher over thy brethren.

		% ---
		\paragraph{1 Nephi 2:23} % *1NEPHI-2-23 
		% ---
			\label{1NEPHI-2-23}

			For behold, in that day that they shall rebel against me,
			I will curse them even with a sore curse,
			and they shall have no power over thy seed
			except they shall rebel against me also.

		% ---
		\paragraph{1 Nephi 2:24} % *1NEPHI-2-24 
		% ---
			\label{1NEPHI-2-24}

			And if it so be that they rebel against me,
			they shall be a scourge unto thy seed
			to stir them up in the ways of remembrance.
			
	% -----
	\subsubsection{1 Nephi 3} % *1NEPHI-3
	% -----
		\label{1NEPHI-3}
		
		% ---
		\paragraph{1 Nephi 3:1} % *1NEPHI-3-1 
		% ---
			\label{1NEPHI-3-1}

			And it came to pass that I Nephi returned
			from speaking with the Lord to the tent of my father.

		% ---
		\paragraph{1 Nephi 3:2} % *1NEPHI-3-2 
		% ---
			\label{1NEPHI-3-2}

			And it came to pass that he spake unto me, saying:
			Behold, I have dreamed a dream
			in the which the Lord hath commanded me
			that thou and thy brethren shall return to Jerusalem.

		% ---
		\paragraph{1 Nephi 3:3} % *1NEPHI-3-3 
		% ---
			\label{1NEPHI-3-3}

			For behold, Laban hath the record of the Jews
			and also a genealogy of my forefathers,
			and they are engraven upon plates of brass.

		% ---
		\paragraph{1 Nephi 3:4} % *1NEPHI-3-4 
		% ---
			\label{1NEPHI-3-4}

			Wherefore the Lord hath commanded me
			that thou and thy brothers should go unto the house of Laban
			and seek the records and bring them down hither into the wilderness.

		% ---
		\paragraph{1 Nephi 3:5} % *1NEPHI-3-5 
		% ---
			\label{1NEPHI-3-5}

			And now behold, thy brothers murmur,
			saying it is a hard thing which I have required of them.
			But behold, I have not required it of them,
			but it is a commandment of the Lord.

		% ---
		\paragraph{1 Nephi 3:6} % *1NEPHI-3-6 
		% ---
			\label{1NEPHI-3-6}

			Therefore go, my son,
			and thou shalt be favored of the Lord
			because thou hast not murmured.

		% ---
		\paragraph{1 Nephi 3:7} % *1NEPHI-3-7 
		% ---
			\label{1NEPHI-3-7}

			And it came to pass that I Nephi said unto my father:
			I will go and do the things which the Lord hath commanded,
			for I know that the Lord giveth no commandments unto the children of men
			save he shall prepare a way for them%
				\footnote{%
					% \label{}%
					Compare \hyperref[ROMANS-4-19]{Romans 4:19--22}.
					}
			that they may accomplish the thing which he commandeth them.%
				\footnote{%
					\label{1NEPHI-3-7-NOTE-better}%
					\hyperref[1NEPHI-17-3]{1 Nephi 17:3} gives what I believe is a better and more thorough version of this principle. This verse gets most of the attention but the other is more helpful. Compare also \hyperref[1NEPHI-5-8]{1 Nephi 5:8}, \hyperref[DC5-34]{DC 5:34}, 
					}
					
		% ---
		\paragraph{1 Nephi 3:8} % *1NEPHI-3-8 
		% ---
			\label{1NEPHI-3-8}

			And it came to pass that when my father had heard these words,
			he was exceeding glad,
			for he knew that I had been blessed of the Lord.

		% ---
		\paragraph{1 Nephi 3:9} % *1NEPHI-3-9 
		% ---
			\label{1NEPHI-3-9}

			And I Nephi and my brethren took our journey in the wilderness with our tents
			to go up to the land of Jerusalem.

		% ---
		\paragraph{1 Nephi 3:10} % *1NEPHI-3-10 
		% ---
			\label{1NEPHI-3-10}

			And it came to pass that when we had gone up to the land of Jerusalem,
			I and my brethren did consult one with another.

		% ---
		\paragraph{1 Nephi 3:11} % *1NEPHI-3-11 
		% ---
			\label{1NEPHI-3-11}

			And we cast lots which of us should go in unto the house of Laban.
			And it came to pass that the lot fell upon Laman.
			And Laman went in unto the house of Laban,
			and he talked with him as he sat in his house.

		% ---
		\paragraph{1 Nephi 3:12} % *1NEPHI-3-12 
		% ---
			\label{1NEPHI-3-12}

			And he desired of Laban the records
			which were engraven upon the plates of brass,
			which contained the genealogy of my father.

		% ---
		\paragraph{1 Nephi 3:13} % *1NEPHI-3-13 
		% ---
			\label{1NEPHI-3-13}

			And behold, it came to pass that Laban was angry
			and thrust him out from his presence,
			and he would not that he should have the records.
			Wherefore he said unto him:
			Behold, thou art a robber and I will slay thee.

		% ---
		\paragraph{1 Nephi 3:14} % *1NEPHI-3-14 
		% ---
			\label{1NEPHI-3-14}

			But Laman fled out of his presence
			and told the things which Laban had done
			unto us.
			And we began to be exceeding sorrowful,
			and my brethren were about to return unto my father in the wilderness.

		% ---
		\paragraph{1 Nephi 3:15} % *1NEPHI-3-15 
		% ---
			\label{1NEPHI-3-15}

			But behold, I said unto them that
			as the Lord liveth and as we live,
			we will not go down unto our father in the wilderness
			until we have accomplished the thing which the Lord hath commanded us.

		% ---
		\paragraph{1 Nephi 3:16} % *1NEPHI-3-16 
		% ---
			\label{1NEPHI-3-16}

			Wherefore let us be faithful in keeping the commandments of the Lord.
			Therefore let us go down to the land of our father’s inheritance.
			For behold, he left gold and silver and all manner of riches,
			and all this he hath done because of the commandment.

		% ---
		\paragraph{1 Nephi 3:17} % *1NEPHI-3-17 
		% ---
			\label{1NEPHI-3-17}

			For he knowing that Jerusalem must be destroyed
			because of the wickedness of the people

		% ---
		\paragraph{1 Nephi 3:18} % *1NEPHI-3-18 
		% ---
			\label{1NEPHI-3-18}

			---for behold, they have rejected the words of the prophets---
			wherefore if my father should dwell in the land
			after that he hath been commanded to flee out of the land,
			behold, he would also perish;
			wherefore it must needs be that he flee out of the land.

		% ---
		\paragraph{1 Nephi 3:19} % *1NEPHI-3-19 
		% ---
			\label{1NEPHI-3-19}

			And behold, it is wisdom in God that we should obtain these records
			that we might preserve unto our children the language of our fathers,

		% ---
		\paragraph{1 Nephi 3:20} % *1NEPHI-3-20 
		% ---
			\label{1NEPHI-3-20}

			and also that we may preserve unto them the words
			which have been spoken by the mouth of all the holy prophets,
			which have been delivered unto them by the Spirit and power of God
			since the world began, even down unto this present time.

		% ---
		\paragraph{1 Nephi 3:21} % *1NEPHI-3-21 
		% ---
			\label{1NEPHI-3-21}

			And it came to pass that after this manner of language
			did I persuade my brethren that they might be faithful
			in keeping the commandments of God.

		% ---
		\paragraph{1 Nephi 3:22} % *1NEPHI-3-22 
		% ---
			\label{1NEPHI-3-22}

			And it came to pass that we went down to the land of our inheritance,
			and we did gather together our gold and our silver and our precious things.

		% ---
		\paragraph{1 Nephi 3:23} % *1NEPHI-3-23 
		% ---
			\label{1NEPHI-3-23}

			And after that we had gathered these things together,
			we went up again unto the house of Laban.

		% ---
		\paragraph{1 Nephi 3:24} % *1NEPHI-3-24 
		% ---
			\label{1NEPHI-3-24}

			And it came to pass that we went in unto Laban and desired him
			that he would give unto us the records
			which were engraven upon the plates of brass,
			for which we would give unto him
			our gold and our silver and all our precious things.

		% ---
		\paragraph{1 Nephi 3:25} % *1NEPHI-3-25 
		% ---
			\label{1NEPHI-3-25}

			And it came to pass that
			when Laban saw our property---that it was exceeding great---
			he did lust after it,
			insomuch that he thrust us out and sent his servants to slay us
			that he might obtain our property.

		% ---
		\paragraph{1 Nephi 3:26} % *1NEPHI-3-26 
		% ---
			\label{1NEPHI-3-26}

			And it came to pass that we did flee before the servants of Laban,
			and we were obliged to leave behind our property,
			and it fell into the hands of Laban.

		% ---
		\paragraph{1 Nephi 3:27} % *1NEPHI-3-27 
		% ---
			\label{1NEPHI-3-27}

			And it came to pass that we fled into the wilderness,
			and the servants of Laban did not overtake us,
			and we hid ourselves in the \hyperref[CAVITY]{cavity} of a rock.

		% ---
		\paragraph{1 Nephi 3:28} % *1NEPHI-3-28 
		% ---
			\label{1NEPHI-3-28}

			And it came to pass that Laman was angry with me and also with my father
			---and also was Lemuel, for he hearkened unto the words of Laman---
			wherefore Laman and Lemuel did speak many hard words
			unto us their younger brothers,
			and they did smite us,%
				\footnote{%
					% \label{}%
					I wonder how severe this beating actually was---Nephi says twice (\hyperref[1NEPHI-2-16]{2:16}, \hyperref[1NEPHI-4-31]{4:31}) that he is ``large in stature,'' and he's able to restrain Zoram at the end of chapter 4 without a problem, even though Zoram had to have been older than him. So was Nephi just not fighting back? I'm not sure how much actual damage Laman and Lemuel could have been doing here. They must have fought lots of times before.
					}
			even with a rod.

		% ---
		\paragraph{1 Nephi 3:29} % *1NEPHI-3-29 
		% ---
			\label{1NEPHI-3-29}

			And it came to pass as they smote us with a rod,
			behold, an angel of the Lord came and stood before them,
			and he spake unto them, saying:
			Why do ye smite your younger brother with a rod?
			Know ye not that the Lord hath chosen him to be a ruler over you?
			---and this because of your iniquities.
			Behold, thou shalt go up to Jerusalem again,
			and the Lord will deliver Laban into your hands.

		% ---
		\paragraph{1 Nephi 3:30} % *1NEPHI-3-30 
		% ---
			\label{1NEPHI-3-30}

			And after that the angel had spake unto us, he departed.

		% ---
		\paragraph{1 Nephi 3:31} % *1NEPHI-3-31 
		% ---
			\label{1NEPHI-3-31}

			And after that the angel had departed,
			Laman and Lemuel again began to murmur, saying:
			How is it possible that the Lord will deliver Laban into our hands?
			Behold, he is a mighty man and he can command fifty.
			Yea, even he can slay fifty, then why not us?%
				\footnote{%
					\label{1NEPHI-3-31-NOTE-fifty}%
					The smallness of Laman and Lemuel's thinking has always struck me. They needed to think bigger; they needed to realize that bigger things could happen in their lives than something involving fifty people. The same is true for us. Nephi has it right; in the first verse of the next chapter he throws out five-digit numbers, but it just as well might be bigger.
					}
					
	% -----
	\subsubsection{1 Nephi 4} % *1NEPHI-4
	% -----
		\label{1NEPHI-4}
		
		% ---
		\paragraph{1 Nephi 4:1} % *1NEPHI-4-1 
		% ---
			\label{1NEPHI-4-1}

			And it came to pass that I spake unto my brethren, saying:
			Let us go up again unto Jerusalem,
			and let us be faithful in keeping the commandments of the Lord.
			For behold, he is mightier than all the earth.
			Then why not mightier than Laban and his fifty?
			Yea, or even than his tens of thousands?

		% ---
		\paragraph{1 Nephi 4:2} % *1NEPHI-4-2 
		% ---
			\label{1NEPHI-4-2}

			Therefore let us go up.
			Let us be strong like unto Moses,
			for he truly spake unto the waters of the Red Sea
			and they divided hither and thither,
			and our fathers came through out of captivity on dry ground,
			and the armies of Pharaoh did follow
			and were drownded in the waters of the Red Sea.

		% ---
		\paragraph{1 Nephi 4:3} % *1NEPHI-4-3 
		% ---
			\label{1NEPHI-4-3}

			Now behold, ye know that this is true.
			And ye also know that an angel hath spoken unto you.
			Wherefore can ye doubt?
			Let us go up.
			The Lord is able to deliver us, even as our fathers,
			and to destroy Laban, even as the Egyptians.

		% ---
		\paragraph{1 Nephi 4:4} % *1NEPHI-4-4 
		% ---
			\label{1NEPHI-4-4}

			Now when I had spoken these words,
			they was yet wroth and did still continue to murmur.
			Nevertheless they did follow me up
			until we came without the walls of Jerusalem.

		% ---
		\paragraph{1 Nephi 4:5} % *1NEPHI-4-5 
		% ---
			\label{1NEPHI-4-5}

			And it was by night,
			and I caused that they should hide themselves without the wall.
			And after that they had hid themselves,
			I Nephi crept into the city and went forth towards the house of Laban.

		% ---
		\paragraph{1 Nephi 4:6} % *1NEPHI-4-6 
		% ---
			\label{1NEPHI-4-6}

			And I was led by the Spirit,
			not knowing beforehand the things which I should do.%
				\footnote{%
					\label{1NEPHI-4-6-NOTE-spirit}%
					What an amazing, powerful verse. In \textit{The Fountainhead} Rand wrote, ``Throughout the centuries there were men who took first steps down new roads armed with nothing but their own vision'' (678). She got it almost right---we can rely on the Spirit to guide us as we take those steps down new roads in our own life. If we're going to do what Lehi did, and leave everything behind to become something new and get a new life, then we're going to be living like this a lot---living uncomfortably---and we have to get used to it, and used to seeking the guidance of the Spirit.
					}

		% ---
		\paragraph{1 Nephi 4:7} % *1NEPHI-4-7 
		% ---
			\label{1NEPHI-4-7}

			Nevertheless I went forth.
			And as I came near unto the house of Laban, I beheld a man,
			and he had fallen to the earth before me,
			for he was drunken with wine.

		% ---
		\paragraph{1 Nephi 4:8} % *1NEPHI-4-8 
		% ---
			\label{1NEPHI-4-8}

			And when I came to him, I found that it was Laban.

		% ---
		\paragraph{1 Nephi 4:9} % *1NEPHI-4-9 
		% ---
			\label{1NEPHI-4-9}

			And I beheld his sword, and I drew it forth from the sheath thereof.
			And the hilt thereof was of pure gold,
			and the workmanship thereof was exceeding fine.
			And I saw that the blade thereof was of the most precious steel.

		% ---
		\paragraph{1 Nephi 4:10} % *1NEPHI-4-10 
		% ---
			\label{1NEPHI-4-10}

			And it came to pass that
			I was \hyperref[CONSTRAIN]{constrained} by the Spirit that I should kill Laban.
			But I said in my heart:
			Never at any time have I shed the blood of man.%
				\footnote{%
					% \label{1NEPHI-4-10-NOTE-shed}%
					Nephi believed that the Lord would make a way for him to fulfill the commandment which he had been given. But that way turned out to require him to do something that he did not want to do, and was against his personal standards of righteous living. There is so much to learn from this experience. What sticks out to me most is the Lord's lesson to Nephi about what was going to be required of him on this journey. Nephi is searching for the promised land, something better than what he has. But he learns very early on, almost at the very beginning, that arriving at the promised land will force him to do things he has never done before, to live in ways he has never lived before, and to do new things and things that make him feel uncomfortable. Things have changed, and the old rules no longer apply. He cannot get to where he wants to go by living in the same ways as before. That isn't going to cut it. He has to enter the discomfort zone.
					}
			And I
				\footnote{%
					% \label{}%
					Note the similarity of the language here, ``shrunk, ''would...might not,'' to that in \hyperref[DC19-18]{DC 19:18}, where Christ describes his suffering during the atonement.
					}%
			\hyperref[SHRINK]{shrunk} and would that I might not slay him.

		% ---
		\paragraph{1 Nephi 4:11} % *1NEPHI-4-11 
		% ---
			\label{1NEPHI-4-11}

			And the Spirit saith unto me again:
			Behold, the Lord hath delivered him into thy hands.
			Yea, and I also knew that he had sought to take away mine own life.
			Yea, and he would not hearken unto the commandments of the Lord.
			And he also had taken away our property.

		% ---
		\paragraph{1 Nephi 4:12} % *1NEPHI-4-12 
		% ---
			\label{1NEPHI-4-12}

			And it came to pass that the Spirit said unto me again:
			Slay him, for the Lord hath delivered him into thy hands.

		% ---
		\paragraph{1 Nephi 4:13} % *1NEPHI-4-13 
		% ---
			\label{1NEPHI-4-13}

			Behold, the Lord slayeth the wicked to bring forth his righteous purposes.%
				\footnote{%
					% \label{}%
					Compare the remark from Caiaphas in \hyperref[JOHN-11-50]{John 11:50}.
					}
			It is better that one man should perish
			than that a nation should dwindle and perish in unbelief.%
				\footnote{%
					\label{1NEPHI-4-13-NOTE-spirit}%
					Note the differences in what the Spirit says and what Nephi says:
					
						\begin{tabular}{p{3in}p{3in}}
						\textbf{Spirit} & \textbf{Nephi} \\
						Kill Laban (``I was constrained by the Spirit that I should kill Laban'') & Never at any time have I shed the blood of man \\
						Behold, the Lord delivered him into your hands & Laban had tried to kill me; Laban wouldn't listen to the Lord's commandments; Laban took away our property \\
						Slay him, for the Lord delivered him into your hands &  \\
						The Lord slays wicked people to bring forth his righteous purposes &  \\
						It is better that one man should perish than that a nation should dwindle and perish in unbelief &  \\
						 & I remembered the words of the Lord which he spake unto me in the wilderness, saying that inasmuch as thy seed shall keep my commandments, they shall prosper in the land of promise \\
						 &  They could not keep the commandments of the Lord according to the law of Moses save they should have the law \\
						 & I also knew that the law was engraven upon the plates of brass \\
						 & I knew that the Lord had delivered Laban into my hands for this cause that I might obtain the records according to his commandments \\
						\end{tabular}
						
					The Spirit focuses on delivering the command; the rationalization is mostly left to Nephi, although the Spirit does give a reason at the end. But I think that reason is secondary to the command, and probably won't always come when we are commanded to do the new things that will lead us where the Lord wants us to go. The first three things the Spirit says are just ``Kill him, the Lord delivered him to you'' and that's it. Just like Nephi, we have to be willing to do new things if we want new outcomes. We cannot get to where the Lord wants us to go by doing the same old things.
					}

		% ---
		\paragraph{1 Nephi 4:14} % *1NEPHI-4-14 
		% ---
			\label{1NEPHI-4-14}

			And now when I Nephi had heard these words,
			I remembered the words of the Lord
			which he spake unto me in the wilderness,
			saying that inasmuch as thy seed shall keep my commandments,
			they shall prosper in the land of promise.

		% ---
		\paragraph{1 Nephi 4:15} % *1NEPHI-4-15 
		% ---
			\label{1NEPHI-4-15}

			Yea, and I also thought that they could not keep the commandments of the Lord
			according to the law of Moses save they should have the law.

		% ---
		\paragraph{1 Nephi 4:16} % *1NEPHI-4-16 
		% ---
			\label{1NEPHI-4-16}

			And I also knew that the law was engraven upon the plates of brass.

		% ---
		\paragraph{1 Nephi 4:17} % *1NEPHI-4-17 
		% ---
			\label{1NEPHI-4-17}

			And again I knew that the Lord had delivered Laban into my hands for this cause
			that I might obtain the records according to his commandments.%
				\footnote{%
					\label{1NEPHI-4-17-NOTE-knew}%
					Here Nephi has a moment of clarity where everything in his mind clicks into place and everything makes sense. All the reasons and considerations lined up and it was clear to him what he needed to do, and why. That's great for him, but I think we also need to remember that it won't always be this way. We're not always going to be in a position where we can make the decisions \textit{after} the moment of clarity, or after having the reasons presented to us. Sometimes we just continue as in verse 6, without knowing exactly what we're doing or why. But again, like before, we should expect this, because we cannot hope to get new results by doing old things.
					}

		% ---
		\paragraph{1 Nephi 4:18} % *1NEPHI-4-18 
		% ---
			\label{1NEPHI-4-18}

			Therefore I did obey the voice of the Spirit
			and took Laban by the hair of the head,
			and I smote off his head with his own sword.

		% ---
		\paragraph{1 Nephi 4:19} % *1NEPHI-4-19 
		% ---
			\label{1NEPHI-4-19}

			And after that I had smote off his head with his own sword,
			I took the garments of Laban and put them
			upon mine own body, yea, even every whit.
			And I did gird on his armor about my loins.

		% ---
		\paragraph{1 Nephi 4:20} % *1NEPHI-4-20 
		% ---
			\label{1NEPHI-4-20}

			And after that I had done this, I went forth unto the treasury of Laban.
			And as I went forth towards the treasury of Laban,
			behold, I saw the servant of Laban which had the keys of the treasury,
			and I commanded him in the voice of Laban%
				\footnote{%
					\label{1NEPHI-4-20-NOTE-voice}%
					President Openshaw used to say that this was an example of how the miracle occurs after the test. After Nephi fulfilled the commandment, he was able to speak in another voice and deceive Zoram in order to be able to get the plates.
					}
			that he should go with me into the treasury.

		% ---
		\paragraph{1 Nephi 4:21} % *1NEPHI-4-21 
		% ---
			\label{1NEPHI-4-21}

			And he supposing me to be his master Laban
			---for he beheld the garments and also the sword girded about my loins---

		% ---
		\paragraph{1 Nephi 4:22} % *1NEPHI-4-22 
		% ---
			\label{1NEPHI-4-22}

			and he spake unto me concerning the elders of the Jews,
			he knowing that his master Laban had been out by night among them.

		% ---
		\paragraph{1 Nephi 4:23} % *1NEPHI-4-23 
		% ---
			\label{1NEPHI-4-23}

			And I spake unto him as if it had been Laban.

		% ---
		\paragraph{1 Nephi 4:24} % *1NEPHI-4-24 
		% ---
			\label{1NEPHI-4-24}

			And I also spake unto him that I should carry the engravings
			which were upon the plates of brass
			to my elder brethren, which were without the wall.

		% ---
		\paragraph{1 Nephi 4:25} % *1NEPHI-4-25 
		% ---
			\label{1NEPHI-4-25}

			And I also bade him that he should follow me.

		% ---
		\paragraph{1 Nephi 4:26} % *1NEPHI-4-26 
		% ---
			\label{1NEPHI-4-26}

			And he supposing that I spake of%
				\footnote{%
					% \label{}%
					How does Nephi know that this is what Zoram was thinking about what he said? Zoram must have told him later.
					}
			the brethren of the church%
				\footnote{%
					\label{1NEPHI-4-26-NOTE-church}%
					The word ``church'' occurs 259 times in the BOM, and this is the first occurrence. It's pretty unusual I think. I was just checking and \gr{ἐκκλησία} appears many times in the LXX, and refers to ``the Jewish congregation,'' according to the LSJ, and then refers to ``the Church'' in the NT. So maybe there is some precedent here, though note that ``church'' does not appear in the OT at all.
					}
			and that I was truly that Laban whom I had slew,
			wherefore he did follow me.

		% ---
		\paragraph{1 Nephi 4:27} % *1NEPHI-4-27 
		% ---
			\label{1NEPHI-4-27}

			And he spake unto me many times concerning the elders of the Jews
			as I went forth unto my brethren, which were without the wall.

		% ---
		\paragraph{1 Nephi 4:28} % *1NEPHI-4-28 
		% ---
			\label{1NEPHI-4-28}

				\footnote{%
					% \label{}%
					So the disguise is convincing enough to fool his own brothers, but it's also nighttime, so they may just not have been able to see things very well.
					}%
			And it came to pass that when Laman saw me,
			he was exceedingly frightened, and also Lemuel and Sam.
			And they fled from before my presence,
			for they supposed it was Laban
			and that he had slain me and had sought to take away their lives also.

		% ---
		\paragraph{1 Nephi 4:29} % *1NEPHI-4-29 
		% ---
			\label{1NEPHI-4-29}

			And it came to pass that I called after them and they did hear me;
			wherefore they did cease to flee from my presence.

		% ---
		\paragraph{1 Nephi 4:30} % *1NEPHI-4-30 
		% ---
			\label{1NEPHI-4-30}

			And it came to pass that when the servant of Laban beheld my brethren,
			he began to tremble and was about to flee from before me
			and return to the city of Jerusalem.

		% ---
		\paragraph{1 Nephi 4:31} % *1NEPHI-4-31 
		% ---
			\label{1NEPHI-4-31}

			And now I Nephi being a man large in stature,
			and also having received much strength of the Lord,
			therefore I did seize upon the servant of Laban
			and held him that he should not flee.

		% ---
		\paragraph{1 Nephi 4:32} % *1NEPHI-4-32 
		% ---
			\label{1NEPHI-4-32}

			And it came to pass that I spake with him
			that if he would hearken unto my words
			---as the Lord liveth and as I live---
			even so that if he would hearken unto our words, we would spare his life.

		% ---
		\paragraph{1 Nephi 4:33} % *1NEPHI-4-33 
		% ---
			\label{1NEPHI-4-33}

			And I spake unto him, even with an oath, that he need not fear,
			that he should be a free man%
				\footnote{%
					\label{1NEPHI-4-33-NOTE-free}%
					Zoram is described as a ``servant,'' but maybe he was actually a slave?
					}
			like unto us
			if he would go down into the wilderness with us.

		% ---
		\paragraph{1 Nephi 4:34} % *1NEPHI-4-34 
		% ---
			\label{1NEPHI-4-34}

			And I also spake unto him, saying:
			Surely the Lord hath commanded us to do this thing.
			And shall we not be diligent in keeping the commandment of the Lord?
			Therefore if thou wilt go down into the wilderness to my father,
			thou shalt have place with us.

		% ---
		\paragraph{1 Nephi 4:35} % *1NEPHI-4-35 
		% ---
			\label{1NEPHI-4-35}

			And it came to pass that Zoram did take courage at the words which I spake%
				\footnote{%
					% \label{}%
					Nephi does not tell us that the Spirit softened Zoram's heart, or anything like that; compare what he does for Ishamel in \hyperref[1NEPHI-7-5]{1 Nephi 7:5}. Here Nephi just says, he took courage at the words. So it's possible that the Spirit was moving him, but it's also possible that Zoram's situation working for Laban and others was just so bad or intolerable that he decided he would rather strike out into the wilderness than go back. Another thought here: we learn in \hyperref[1NEPHI-5-13]{1 Nephi 5:13} that the plates had many of Jeremiah's prophecies. From \hyperref[1NEPHI-7-14]{7:14} it seems that Jeremiah was a contemporary of Lehi's family, so how did his prophecies end up on the plates? Someone, either Laban or someone else, must have put them there. I wonder if it was Zoram---he had access to the records and so could have done it, and he also asks Nephi (as Laban) about the brothers of the church (4:26--27) a bunch of times, so Zoram was really interested in that and concerned about it, at least it seems. So maybe he also saw the trip with Nephi as a chance to pursue something important to him in a way that he wouldn't have been able to otherwise.
					}
			---now ` was the name of the servant---%
				\footnote{%
					% \label{}%
					This is the kind of remark where I wonder if it was in the original text, or if it's just something that JS inserted to clarify something else that was just said. Seems like a place where you can detect some insight into how the translation process occurred.
					}
			and he promised that he would go down into the wilderness unto our father;
			yea, and he also made an oath unto us
			that he would tarry with us from that time forth.

		% ---
		\paragraph{1 Nephi 4:36} % *1NEPHI-4-36 
		% ---
			\label{1NEPHI-4-36}

			Now we were desirous that he should tarry with us for this cause
			that the Jews might not know concerning our flight into the wilderness,
			lest they should pursue us and destroy us.

		% ---
		\paragraph{1 Nephi 4:37} % *1NEPHI-4-37 
		% ---
			\label{1NEPHI-4-37}

			And it came to pass that when Zoram had made an oath unto us,
			our fears did cease concerning him.

		% ---
		\paragraph{1 Nephi 4:38} % *1NEPHI-4-38 
		% ---
			\label{1NEPHI-4-38}

			And it came to pass that we took the plates of brass and the servant of Laban
			and departed into the wilderness and journeyed unto the tent of our father.
			
	% -----
	\subsubsection{1 Nephi 5} % *1NEPHI-5
	% -----
		\label{1NEPHI-5}

		% ---
		\paragraph{1 Nephi 5:1} % *1NEPHI-5-1 
		% ---
			\label{1NEPHI-5-1}

			And it came to pass that after we had came down
			into the wilderness unto our father,
			behold, he was filled with joy.
			And also my mother Sariah was exceeding glad,
			for she truly had mourned because of us,

		% ---
		\paragraph{1 Nephi 5:2} % *1NEPHI-5-2 
		% ---
			\label{1NEPHI-5-2}

			for she had supposed that we had perished in the wilderness.
			And she also had complained against my father,
			telling him that he was a visionary man, saying:
			Behold, thou hast led us forth from the land of our inheritance,
			and my sons are no more, and we perish in the wilderness.

		% ---
		\paragraph{1 Nephi 5:3} % *1NEPHI-5-3 
		% ---
			\label{1NEPHI-5-3}

			And after this manner of language had my mother complained against my father.

		% ---
		\paragraph{1 Nephi 5:4} % *1NEPHI-5-4 
		% ---
			\label{1NEPHI-5-4}

			And it had came to pass that my father spake unto her, saying:
			I know that I am a visionary man,
			for if I had not seen the things of God in a vision,
			I should not have known the goodness%
				\footnote{%
					% \label{}%
					Lehi mentions God's goodness in \hyperref[1NEPHI-1-14]{1 Nephi 1:14}, as part of his visionary experience.
					}
			of God
			but had tarried at Jerusalem and had perished with my brethren.

		% ---
		\paragraph{1 Nephi 5:5} % *1NEPHI-5-5 
		% ---
			\label{1NEPHI-5-5}

			But behold, I have obtained a land of promise,%
				\footnote{%
					\label{1NEPHI-5-5-NOTE-have}%
					Why the perfect tense here, ``have obtained''? The perfect tense indicates that an action has been completed, so Lehi's obtaining of the ``land of promise'' he's talking about here is already done. Since they're still very far from the actual land of promise and haven't even gotten to the beach yet, he must be referring to something else: the new life and freedom he has been given through his interactions with God. For him, giving up all those possessions and his property was a good trade to receive that personal land of promise, and the serenity and self-respect that come from living the way we know we ought to and really should are priceless. In \textit{The Fountainhead}, getting on toward the end, Keating says, ``Katie, why do they always teach us that it's easy and evil to do what we want and that we need discipline to restrain ourselves? It's the hardest thing in the world---to do what we want. And it takes the greatest kind of courage. I mean, what we really want. As I wanted to marry you. Not as I want to sleep with some woman or get drunk or get my name in the papers. Those things---they're not even desires---they're things people do to escape from desires---because it's such a big responsibility, really to want something'' (598--599). The self-respect that comes from righteous living is living the way you really want.
					
					See also \hyperref[DC38-18]{DC 38:18--20};
					}
			in the which things I do rejoice.
			Yea, and I know that the Lord will deliver my sons out of the hands of Laban
			and bring them down again unto us in the wilderness.

		% ---
		\paragraph{1 Nephi 5:6} % *1NEPHI-5-6 
		% ---
			\label{1NEPHI-5-6}

			And after this manner of language
			did my father Lehi comfort my mother Sariah concerning us
			while we journeyed in the wilderness up to the land of Jerusalem
			to obtain the record of the Jews.

		% ---
		\paragraph{1 Nephi 5:7} % *1NEPHI-5-7 
		% ---
			\label{1NEPHI-5-7}

			And when we had returned to the tent of my father,
			behold, their joy was full, and my mother was comforted.

		% ---
		\paragraph{1 Nephi 5:8} % *1NEPHI-5-8 
		% ---
			\label{1NEPHI-5-8}

			And she spake, saying:
			Now I know of a surety that the Lord hath commanded
			my husband to flee into the wilderness;
			yea, and I also know of a surety
			that the Lord hath protected my sons
			and delivered them out of the hands of Laban
			and gave them power whereby they could accomplish the thing
			which the Lord hath commanded them.%
				\footnote{%
					\label{1NEPHI-5-8-NOTE-sariah}%
					Each person in Lehi's family experienced their journey differently, and faced different tests. Sariah's test had to do with her children. She could not bear losing them and struggled when she thought they were gone, but when they returned, she came to a new conviction about what they were doing: ``Now I know.'' But sending his sons away was not hard for Lehi; his test came later, and had to do with food. Each person had a different test and the journey brought out challenges for whatever was deepest and most important or most crucial to them as individuals.
					}
			And after this manner of language did she speak.

		% ---
		\paragraph{1 Nephi 5:9} % *1NEPHI-5-9 
		% ---
			\label{1NEPHI-5-9}

			And it came to pass that they did rejoice exceedingly
			and did offer sacrifice and burnt offerings unto the Lord,
			and they gave thanks unto the God of Israel.

		% ---
		\paragraph{1 Nephi 5:10} % *1NEPHI-5-10 
		% ---
			\label{1NEPHI-5-10}

			And after that they had given thanks unto the God of Israel,
			my father Lehi took the records which were engraven upon the plates of brass
			and he did \hyperref[SEARCH-1828]{search} them from the beginning.

		% ---
		\paragraph{1 Nephi 5:11} % *1NEPHI-5-11 
		% ---
			\label{1NEPHI-5-11}

			And he beheld that they did contain the five books of Moses,%
				\footnote{%
					% \label{}%
					This description of the books, which are really the Torah, is very Christian and seems to be a clear example of JS coming through in the translation, or rather, JS's context.
					}
			which gave an account of the creation of the world
			and also of Adam and Eve, which was our first parents,

		% ---
		\paragraph{1 Nephi 5:12} % *1NEPHI-5-12 
		% ---
			\label{1NEPHI-5-12}

			and also a record of the Jews from the beginning,
			even down to the commencement of the reign of Zedekiah, king of Judah,

		% ---
		\paragraph{1 Nephi 5:13} % *1NEPHI-5-13 
		% ---
			\label{1NEPHI-5-13}

			and also the prophecies of the holy prophets from the beginning,
			even down to the commencement of the reign of Zedekiah,
			and also many prophecies which have been spoken by the mouth of Jeremiah.

		% ---
		\paragraph{1 Nephi 5:14} % *1NEPHI-5-14 
		% ---
			\label{1NEPHI-5-14}

			And it came to pass that my father Lehi also found
			upon the plates of brass a genealogy of his fathers;
			wherefore he knew that he was a descendant of Joseph,%
				\footnote{%
					% \label{}%
					Lehi names his two younger children ``Joseph'' and ``Jacob'' in honor of what he learns here (\hyperref[1NEPHI-18-7]{1 Nephi 18:7}).
					}
			yea, even that Joseph which was the son of Jacob,
			which was sold into Egypt and which was preserved by the hand of the Lord
			that he might preserve his father Jacob and all his household
			from perishing with famine.

		% ---
		\paragraph{1 Nephi 5:15} % *1NEPHI-5-15 
		% ---
			\label{1NEPHI-5-15}

			And they were also led out of captivity and out of the land of Egypt
			by that same God who had preserved them.

		% ---
		\paragraph{1 Nephi 5:16} % *1NEPHI-5-16 
		% ---
			\label{1NEPHI-5-16}

			And thus my father Lehi did discover the genealogy of his fathers.
			And Laban also was a descendant of Joseph;
			wherefore he and his fathers had kept the records.

		% ---
		\paragraph{1 Nephi 5:17} % *1NEPHI-5-17 
		% ---
			\label{1NEPHI-5-17}

			And now when my father saw all these things,
			he was filled with the Spirit and began to prophesy concerning his seed,

		% ---
		\paragraph{1 Nephi 5:18} % *1NEPHI-5-18 
		% ---
			\label{1NEPHI-5-18}

			that these plates of brass should go forth
			unto all nations, kindreds, tongues, and people which were of his seed.

		% ---
		\paragraph{1 Nephi 5:19} % *1NEPHI-5-19 
		% ---
			\label{1NEPHI-5-19}

			Wherefore he said that these plates of brass should never perish,
			neither should they be dimmed any more by time.%
				\footnote{%
					% \label{}%
					Hard to know the meaning of this prophecy---one thing it might mean is that the \textit{meaning} of the words on the brass plates would begin to be clear, and wouldn't be ``dimmed'' any longer, maybe because of the light that the Lord was going to give future Nephite peoples, and through them, the restoration? Not sure. Compare the discussion of ``brightness'' in \hyperref[ALMA-37]{Alma 37}.
					}
			And he prophesied many things concerning his seed.

		% ---
		\paragraph{1 Nephi 5:20} % *1NEPHI-5-20 
		% ---
			\label{1NEPHI-5-20}

			And it came to pass that thus far I and my father had kept the commandments
			wherewith the Lord had commanded us.

		% ---
		\paragraph{1 Nephi 5:21} % *1NEPHI-5-21 
		% ---
			\label{1NEPHI-5-21}

			And we had obtained the record which the Lord had commanded us
			and searched them and found that they were desirable,
			yea, even of great worth unto us,
			insomuch that we could preserve the commandments of the Lord unto our children.%
				\footnote{%
					\label{1NEPHI-5-21-NOTE-commandments}%
					This bit about the ``commandments'' is one of the reasons Nephi mentions when he discusses his thought process before killing Laban in \hyperref[1NEPHI-4-15]{1 Nephi 4:15}. Given that he's writing this many years after he experiences the events, I wonder how much that reason was actually in his head at the time he took Laban's life, and how much it occured to him later as he realized the value of having the scriptures around to teach his children.
					}

		% ---
		\paragraph{1 Nephi 5:22} % *1NEPHI-5-22 
		% ---
			\label{1NEPHI-5-22}

			Wherefore it was wisdom in the Lord that we should carry them with us
			as we journeyed in the wilderness toward the land of promise.
			
	% -----
	\subsubsection{1 Nephi 6} % *1NEPHI-6
	% -----
		\label{1NEPHI-6}
		
		% ---
		\paragraph{1 Nephi 6:1} % *1NEPHI-6-1 
		% ---
			\label{1NEPHI-6-1}

			And now I Nephi do not give the genealogy of my fathers in this part of my record,
			neither at any time shall I give it after upon these plates which I am writing,
			for it is given in the record which has been kept by my father;
			wherefore I do not write it in this work.

		% ---
		\paragraph{1 Nephi 6:2} % *1NEPHI-6-2 
		% ---
			\label{1NEPHI-6-2}

			For it sufficeth me to say that we are a descendant of Joseph.

		% ---
		\paragraph{1 Nephi 6:3} % *1NEPHI-6-3 
		% ---
			\label{1NEPHI-6-3}

			And it mattereth not to me that I am particular to give a full account
			of all the things of my father,
			for they cannot be written upon these plates,
			for I desire the room that I may write of the things of God.

		% ---
		\paragraph{1 Nephi 6:4} % *1NEPHI-6-4 
		% ---
			\label{1NEPHI-6-4}

			For the fullness of mine intent is that I may persuade men to come
			unto the God of Abraham and the God of Isaac and the God of Jacob and be saved.

		% ---
		\paragraph{1 Nephi 6:5} % *1NEPHI-6-5 
		% ---
			\label{1NEPHI-6-5}

			Wherefore the things which are pleasing unto the world I do not write,
			but the things which are pleasing unto God
			and unto them which are not of the world.

		% ---
		\paragraph{1 Nephi 6:6} % *1NEPHI-6-6 
		% ---
			\label{1NEPHI-6-6}

			Wherefore I shall give commandment unto my seed
			that they shall not occupy these plates with things
			which are not of worth unto the children of men.
			
	% -----
	\subsubsection{1 Nephi 7} % *1NEPHI-7 
	% -----
		\label{1NEPHI-7}
		
		% ---
		\paragraph{1 Nephi 7:1} % *1NEPHI-7-1 
		% ---
			\label{1NEPHI-7-1}

			And now I would that ye might know that
			after my father Lehi had made an end of prophesying concerning his seed,
			it came to pass that the Lord spake unto him again
			that it was not meet for him Lehi
			that he should take his family into the wilderness alone,
			but that his sons should take daughters to wife
			that they might raise up seed unto the Lord in the land of promise.

		% ---
		\paragraph{1 Nephi 7:2} % *1NEPHI-7-2 
		% ---
			\label{1NEPHI-7-2}

			And it came to pass that the Lord commanded him
			that I Nephi and my brethren should again return into the land of Jerusalem
			and bring down Ishmael and his family into the wilderness.

		% ---
		\paragraph{1 Nephi 7:3} % *1NEPHI-7-3 
		% ---
			\label{1NEPHI-7-3}

			And it came to pass that I Nephi did again with my brethren
			go forth into the wilderness to go up to Jerusalem.

		% ---
		\paragraph{1 Nephi 7:4} % *1NEPHI-7-4 
		% ---
			\label{1NEPHI-7-4}

			And it came to pass that we went up unto the house of Ishmael,
			and we did gain favor in the sight of Ishmael,
			insomuch that we did speak unto him the words of the Lord.

		% ---
		\paragraph{1 Nephi 7:5} % *1NEPHI-7-5 
		% ---
			\label{1NEPHI-7-5}

			And it came to pass that the Lord did soften the heart of Ishmael
			and also his whole household,
			insomuch that they took their journey with us
			down into the wilderness to the tent of our father.%
				\footnote{%
					\label{1-NEPHI-7-5-NOTE-ishmael}%
					I've sometimes wondered whether Nephi doesn't give many details about what happened when they visited Ishmael because Laman and Lemuel did most of the work to convince him or his family. It's just speculation, and Nephi says that the Lord ``did soften the heart of Ishmael,'' but it could also be that Laman and Lemuel were just better at that kind of task.
					}

		% ---
		\paragraph{1 Nephi 7:6} % *1NEPHI-7-6 
		% ---
			\label{1NEPHI-7-6}

			And it came to pass that as we journeyed in the wilderness,
			behold, Laman and Lemuel and two of the daughters of Ishmael
			and the two sons of Ishmael and their families did rebel against us,
			yea, against I Nephi and Sam
			and their father Ishmael and his wife and his three other daughters.

		% ---
		\paragraph{1 Nephi 7:7} % *1NEPHI-7-7 
		% ---
			\label{1NEPHI-7-7}

			And it came to pass that in the which rebellion
			they were desirous to return unto the land of Jerusalem.

		% ---
		\paragraph{1 Nephi 7:8} % *1NEPHI-7-8 
		% ---
			\label{1NEPHI-7-8}

			And now I Nephi being grieved for the hardness of their hearts,
			therefore I spake unto them, saying,
			yea, even unto Laman and unto Lemuel:
			Behold, thou art mine elder brethren,
			and how is it that ye are so hard in your hearts and so blind in your minds
			that ye have need that I, your younger brother, should speak unto you?
			Yea, and set an example for you?

		% ---
		\paragraph{1 Nephi 7:9} % *1NEPHI-7-9 
		% ---
			\label{1NEPHI-7-9}

			How is it that ye have not hearkened unto the word of the Lord?

		% ---
		\paragraph{1 Nephi 7:10} % *1NEPHI-7-10 
		% ---
			\label{1NEPHI-7-10}

			How is it that ye have forgotten that ye have seen an angel of the Lord?

		% ---
		\paragraph{1 Nephi 7:11} % *1NEPHI-7-11 
		% ---
			\label{1NEPHI-7-11}

			Yea, and how is it that ye have forgotten
			how great things the Lord hath done for us
			in delivering us out of the hands of Laban
			and also that we should obtain the record?

		% ---
		\paragraph{1 Nephi 7:12} % *1NEPHI-7-12 
		% ---
			\label{1NEPHI-7-12}

			Yea, and how is it that ye have forgotten
			that the Lord is able to do all things
			according to his will for the children of men?
			---if it so be that they exercise faith in him.
			Wherefore let us be faithful in him.

		% ---
		\paragraph{1 Nephi 7:13} % *1NEPHI-7-13 
		% ---
			\label{1NEPHI-7-13}

			And if it so be that we are faithful in him,
			we shall obtain the land of promise.
			And ye shall know at some future period
			that the word of the Lord shall be fulfilled
			concerning the destruction of Jerusalem,
			for all things which the Lord hath spoken
			concerning the destruction of Jerusalem must be fulfilled.

		% ---
		\paragraph{1 Nephi 7:14} % *1NEPHI-7-14 
		% ---
			\label{1NEPHI-7-14}

			For behold, the Spirit of the Lord ceaseth soon to strive with them.
			For behold, they have rejected the prophets,
			and Jeremiah have they cast into prison,
			and they have sought to take away the life of my father,
			insomuch that they have driven him out of the land.

		% ---
		\paragraph{1 Nephi 7:15} % *1NEPHI-7-15 
		% ---
			\label{1NEPHI-7-15}

			Now behold, I say unto you
			that if ye will return unto Jerusalem, ye shall also perish with them.
			And now if ye have choice, go up to the land,
			and remember the words which I speak unto you,
			that if ye go, ye will also perish.
			For thus the Spirit of the Lord \hyperref[CONSTRAIN]{constraineth} me that I should speak.

		% ---
		\paragraph{1 Nephi 7:16} % *1NEPHI-7-16 
		% ---
			\label{1NEPHI-7-16}

			And it came to pass that when I Nephi had spoken these words unto my brethren,
			they were angry with me.
			And it came to pass that they did lay their hands upon me,
			for behold, they were exceeding wroth.
			And they did bind me with cords,
			for they sought to take away my life,
			that they might leave me in the wilderness to be devoured by wild beasts.

		% ---
		\paragraph{1 Nephi 7:17} % *1NEPHI-7-17 
		% ---
			\label{1NEPHI-7-17}

			But it came to pass that I prayed unto the Lord, saying:
			O Lord, according to my faith which is in me,
			wilt thou deliver me from the hands of my brethren?
			Yea, even give me strength%
				\footnote{%
					% \label{}%
					I think he prays for strength because he knows that he's a big guy and getting more strength would be the natural way for him to break out of the ropes or whatever is holding him. In addition, he's previously said (\hyperref[1NEPHI-4-31]{4:31}) that he had received ``much strength'' from the Lord, so it would have made sense for him to get strength to overcome this problem. But as we see in the next verse, that is not what happens.
					}
			that I may burst these bands with which I am bound.

		% ---
		\paragraph{1 Nephi 7:18} % *1NEPHI-7-18 
		% ---
			\label{1NEPHI-7-18}

			And it came to pass that when I had said these words,
			behold, the bands were loosed%
				\footnote{%
					\label{1NEPHI-7-18-NOTE-loosed}%
					He prays for strength to burst the bands, but the answer comes in the form of the bands being loosened so he can get out of them.
					}
			from off my hands and feet,
			and I stood before my brethren, and I spake unto them again.

		% ---
		\paragraph{1 Nephi 7:19} % *1NEPHI-7-19 
		% ---
			\label{1NEPHI-7-19}

			And it came to pass that they were angry with me again
			and sought to lay hands upon me.
			But behold, one of the daughters of Ishmael,
			yea, and also her mother and one of the sons of Ishmael,%
				\footnote{%
					\label{1NEPHI-7-19-NOTE-sons}%
					It's always kind of stuck out to me that one of Ishmael's sons is now defending Nephi. Back in \hyperref[1NEPHI-7-6]{verse 6}, both the sons are joining Laman and Lemuel in giving Nephi a hard time, but Nephi's speech, and possibly also his escape from the restraints, make one of the sons rethink what he was doing. See \hyperref[2NEPHI-4-13]{2 Nephi 4:13} as well, where that son is apparently back with Laman and Lemuel.
					}
			did plead with my brethren,
			insomuch that they did soften their hearts
			and they did cease striving to take away my life.

		% ---
		\paragraph{1 Nephi 7:20} % *1NEPHI-7-20 
		% ---
			\label{1NEPHI-7-20}

			And it came to pass that they were sorrowful because of their wickedness,
			insomuch that they did bow down before me
			and did plead with me that I would forgive them of the thing
			that they had done against me.

		% ---
		\paragraph{1 Nephi 7:21} % *1NEPHI-7-21 
		% ---
			\label{1NEPHI-7-21}

			And it came to pass that I did frankly%
				\footnote{%
					% \label{}%
					Compare \hyperref[MATTHEW-18-35]{Matthew 18:35}: \gr{ἐὰν μὴ ἀφῆτε ἕκαστος τῷ ἀδελφῷ αὐτοῦ ἀπὸ τῶν καρδιῶν ὑμῶν}.
					}
			forgive them all that they had done.
			And I did exhort them that they would pray
			unto the Lord their God for forgiveness.
			And it came to pass that they did so.
			And after that they had done praying unto the Lord,
			we did again travel on our journey toward the tent of our father.

		% ---
		\paragraph{1 Nephi 7:22} % *1NEPHI-7-22 
		% ---
			\label{1NEPHI-7-22}

			And it came to pass that we did come down unto the tent of our father.
			And after that I and my brethren and all the house of Ishmael
			had come down unto the tent of my father,
			they did give thanks unto the Lord their God,
			and they did offer sacrifice and burnt offerings unto him.

	% -----
	\subsubsection{1 Nephi 8} % *1NEPHI-8 
	% -----
		\label{1NEPHI-8}
		
		% ---
		\paragraph{1 Nephi 8:1} % *1NEPHI-8-1 
		% ---
			\label{1NEPHI-8-1}

			\footnote{%
				% \label{}%
				I wonder if the gathering and sorting of all these seeds and fruit was on Lehi's mind and contributed to him having the dream. They also talk about ``seed'' (verse 3) as offspring a lot.
				}%
			And it came to pass that we had gathered together
			all manner of seeds of every kind,
			both of grain of every kind and also of the seeds of fruits of every kind.

		% ---
		\paragraph{1 Nephi 8:2} % *1NEPHI-8-2 
		% ---
			\label{1NEPHI-8-2}

			And it came to pass that while my father tarried in the wilderness,
			he spake unto us, saying:
			Behold, I have dreamed a dream,
			or in other words, I have seen a vision.

		% ---
		\paragraph{1 Nephi 8:3} % *1NEPHI-8-3 
		% ---
			\label{1NEPHI-8-3}

			And behold, because of the thing which I have seen,
			I have reason to rejoice in the Lord because of Nephi and also of Sam.
			For I have reason to suppose
			that they and also many of their seed will be saved.

		% ---
		\paragraph{1 Nephi 8:4} % *1NEPHI-8-4 
		% ---
			\label{1NEPHI-8-4}

			But behold, Laman and Lemuel, I fear exceedingly because of you.
			For behold, methought I saw a dark and dreary wilderness.

		% ---
		\paragraph{1 Nephi 8:5} % *1NEPHI-8-5 
		% ---
			\label{1NEPHI-8-5}

			And it came to pass that I saw a man and he was dressed in a white robe.
			And he came and stood before me.

		% ---
		\paragraph{1 Nephi 8:6} % *1NEPHI-8-6 
		% ---
			\label{1NEPHI-8-6}

			And it came to pass that he spake unto me and bade me follow him.

		% ---
		\paragraph{1 Nephi 8:7} % *1NEPHI-8-7 
		% ---
			\label{1NEPHI-8-7}

			And it came to pass that I followed him.
			And after I had followed him,
			I beheld myself that I was in a dark and dreary waste.

		% ---
		\paragraph{1 Nephi 8:8} % *1NEPHI-8-8 
		% ---
			\label{1NEPHI-8-8}

			And after that I had traveled for the space of many hours in darkness,
			I began to pray unto the Lord that he would have mercy%
				\footnote{%
					% \label{}%
					``Have mercy...according to the multitude of his tender mercies'' is language that parallels \hyperref[PSALMS-51-1]{Psalms 51:1} and is also similar to \hyperref[PSALMS-69-16]{Psalms 69:16}. See ``tender mercy'' \hyperref[MERCY-PHRASES]{here}.
					}
			on me
			according to the multitude of his tender mercies.

		% ---
		\paragraph{1 Nephi 8:9} % *1NEPHI-8-9 
		% ---
			\label{1NEPHI-8-9}

			And it came to pass that after I had prayed unto the Lord,
			I beheld a large and spacious field.

		% ---
		\paragraph{1 Nephi 8:10} % *1NEPHI-8-10 
		% ---
			\label{1NEPHI-8-10}

			And it came to pass that I beheld a tree
			whose fruit was desirable to make one happy.%
				\footnote{%
					\label{1NEPHI-8-10-NOTE-to}%
					This is kind of a weird construction after ``desirable'': ``desirable to make one happy.'' That's just a bit of a weird way of talking, at least now. I think it means, the fruit is desirable \textit{because} it will make you happy, but it's still strange.
					}

		% ---
		\paragraph{1 Nephi 8:11} % *1NEPHI-8-11 
		% ---
			\label{1NEPHI-8-11}

			And it came to pass that I did go forth and partook of the fruit thereof
			and beheld that it was most sweet, above all that I ever had before tasted.
			Yea, and I beheld that the fruit thereof was white
			to exceed all the whiteness that I had ever seen.

		% ---
		\paragraph{1 Nephi 8:12} % *1NEPHI-8-12 
		% ---
			\label{1NEPHI-8-12}

			And as I partook of the fruit thereof,
			it filled my soul with exceeding great joy.
			Wherefore I began to be desirous that my family should partake of it also,
			for I knew that it was desirous above all other fruit.

		% ---
		\paragraph{1 Nephi 8:13} % *1NEPHI-8-13 
		% ---
			\label{1NEPHI-8-13}

			And as I cast my eyes around about
			that perhaps I might discover my family also,
			and I beheld a river of water and it ran along,
			and it was near the tree of which I was partaking the fruit.

		% ---
		\paragraph{1 Nephi 8:14} % *1NEPHI-8-14 
		% ---
			\label{1NEPHI-8-14}

			And I looked to behold from whence it came,
			and I saw the head%
				\footnote{%
					\label{1NEPHI-8-14-NOTE-head}%
					In the 1828 dictionary, meaning \#18 for ``head'' is: ``The principal source of a stream; as the head of the Nile.'' That definition helps but I still have a hard time getting a picture in my head of what it's like---is there a large body of water and that's where the river starts from? It carries the water in another direction? I'm not sure.
					}
			thereof a little way off.
			And at the head thereof I beheld your mother Sariah and Sam and Nephi,
			and they stood as if they knew not whither they should go.

		% ---
		\paragraph{1 Nephi 8:15} % *1NEPHI-8-15 
		% ---
			\label{1NEPHI-8-15}

			And it came to pass that I beckoned unto them.
			And I also did say unto them with a loud voice
			that they should come unto me and partake of the fruit,
			which was desirable above all other fruit.

		% ---
		\paragraph{1 Nephi 8:16} % *1NEPHI-8-16 
		% ---
			\label{1NEPHI-8-16}

			And it came to pass that they did come unto me and partake of the fruit also.

		% ---
		\paragraph{1 Nephi 8:17} % *1NEPHI-8-17 
		% ---
			\label{1NEPHI-8-17}

			And it came to pass that I was desirous
			that Laman and Lemuel should come and partake of the fruit also.
			Wherefore I cast mine eyes toward the head of the river
			that perhaps I might see them.

		% ---
		\paragraph{1 Nephi 8:18} % *1NEPHI-8-18 
		% ---
			\label{1NEPHI-8-18}

			And it came to pass that I saw them,
			but they would not come unto me and partake of the fruit.

		% ---
		\paragraph{1 Nephi 8:19} % *1NEPHI-8-19 
		% ---
			\label{1NEPHI-8-19}

			And I beheld a rod of iron,
			and it extended along the bank of the river
			and led to the tree by which I stood.

		% ---
		\paragraph{1 Nephi 8:20} % *1NEPHI-8-20 
		% ---
			\label{1NEPHI-8-20}

			And I also beheld a straight and narrow path
			which came along by the rod of iron,
			even to the tree by which I stood.
			And it also led by the head of the fountain unto a large and spacious field,
			as if it had been a world.

		% ---
		\paragraph{1 Nephi 8:21} % *1NEPHI-8-21 
		% ---
			\label{1NEPHI-8-21}

			And I saw numberless concourses%
				\footnote{%
					% \label{}%
					The phrase ``numberless concourses'' also appears at \hyperref[1NEPHI-1-8]{1 Nephi 1:8}.
					}
			of people,
			many of whom were pressing forward
			that they might obtain the path
			which led unto the tree by which I stood.

		% ---
		\paragraph{1 Nephi 8:22} % *1NEPHI-8-22 
		% ---
			\label{1NEPHI-8-22}

			And it came to pass that they did come forth
			and commenced in the path which led to the tree.

		% ---
		\paragraph{1 Nephi 8:23} % *1NEPHI-8-23 
		% ---
			\label{1NEPHI-8-23}

			And it came to pass that there arose a mist of darkness,
			yea, even an exceeding great mist of darkness,
			insomuch that they which had commenced in the path did lose their way,
			that they wandered off and were lost.

		% ---
		\paragraph{1 Nephi 8:24} % *1NEPHI-8-24 
		% ---
			\label{1NEPHI-8-24}

			And it came to pass that I beheld others pressing forward.
			And they came forth and caught hold of the end of the rod of iron.
			And they did press forward through the mists of darkness,
			clinging to the rod of iron,
			even until they did come forth and partook of the fruit of the tree.

		% ---
		\paragraph{1 Nephi 8:25} % *1NEPHI-8-25 
		% ---
			\label{1NEPHI-8-25}

			And after that they had partook of the fruit of the tree,
			they did cast their eyes about as if they were ashamed.

		% ---
		\paragraph{1 Nephi 8:26} % *1NEPHI-8-26 
		% ---
			\label{1NEPHI-8-26}

			And I also cast my eyes around about
			and beheld on the other side of the river of water a great and spacious building.
			And it stood as it were in the air, high above the earth.

		% ---
		\paragraph{1 Nephi 8:27} % *1NEPHI-8-27 
		% ---
			\label{1NEPHI-8-27}

			And it was filled with people, both old and young, both male and female,
			and their manner of dress was exceeding fine.
			And they were in the attitude of mocking and pointing their fingers
			towards those which had came up and were partaking of the fruit.

		% ---
		\paragraph{1 Nephi 8:28} % *1NEPHI-8-28 
		% ---
			\label{1NEPHI-8-28}

			And after that they had tasted of the fruit, they were ashamed
			because of those that were a scoffing at them,
			and they fell away into forbidden paths and were lost.

		% ---
		\paragraph{1 Nephi 8:29} % *1NEPHI-8-29 
		% ---
			\label{1NEPHI-8-29}

			And now I Nephi do not speak all the words of my father.

		% ---
		\paragraph{1 Nephi 8:30} % *1NEPHI-8-30 
		% ---
			\label{1NEPHI-8-30}

			But to be short in writing, behold, he saw other multitudes pressing forwards.
			And they came and caught hold of the end of the rod of iron.
			And they did press their way forward,
			continually holding fast to the rod of iron,%
				\footnote{%
					\label{1NEPHI-8-30-NOTE-hold-fast}%
					In the October 2011 general conference, David Bednar gave a talk in which he focused on the difference between ``clinging'' to the rod (\hyperref[1NEPHI-8-24]{verse 24}) and ``continually holding fast'' in this verse. One 1828 definition of ``cling'' is ``To adhere closely; to stick to; to hold fast upon,'' and the word ``cling'' is used to define two different senses of ``hold.'' So the meaningful difference isn't between ``cling'' and ``hold fast'' but in the fact that one includes the word ``continually.'' Compare ``lay hold'' at \hyperref[1TIMOTHY-6-12]{1 Timothy 6:12} and \hyperref[1TIMOTHY-6-19]{19}.
					}
			until they came forth and fell down and partook of the fruit of the tree.

		% ---
		\paragraph{1 Nephi 8:31} % *1NEPHI-8-31 
		% ---
			\label{1NEPHI-8-31}

			And he also saw other multitudes pressing their way
			towards that great and spacious building.

		% ---
		\paragraph{1 Nephi 8:32} % *1NEPHI-8-32 
		% ---
			\label{1NEPHI-8-32}

			And it came to pass that many were drowned in the depths of the fountain,
			and many were lost from his view, wandering in strange roads.

		% ---
		\paragraph{1 Nephi 8:33} % *1NEPHI-8-33 
		% ---
			\label{1NEPHI-8-33}

			And great was the multitude that did enter into that strange building.
			And after%
				\footnote{%
					\label{1NEPHI-8-33-NOTE-after}%
					According to Lehi's perception of what happened, the people in the building only mocked those eating the fruit \textit{after} they had entered into the building. The building isn't a place or even a determinate group of people; it's a mindset that you adopt toward those who are trying to live righteous lives. People from all backgrounds and beliefs are able to get to the point where they have that mindset.
					}
			that they did enter into that building,
			they did point the finger of scorn at me
			and those that were partaking of the fruit also,
			but we heeded them not.%
				\footnote{%
					\label{1NEPHI-8-33-NOTE-heeded-not}%
					This is the key, I think. We must push our independence of mind and life to the point where we no longer ``heed'' the world. ``\hyperref[HEED-1828]{Heed}'' means ``To mind; to regard with care; to take notice of; to attend to; to observe.'' I think it's more than just taking notice of or attending to, though; ``to regard with care'' is good, because it makes it clear that what you are hearing \textit{is playing a role in your choices, plans, and behaviors}. That is what we must reject.
					}

		% ---
		\paragraph{1 Nephi 8:34} % *1NEPHI-8-34 
		% ---
			\label{1NEPHI-8-34}

			Thus is the words of my father,
			for as many as heeded them had fallen away.

		% ---
		\paragraph{1 Nephi 8:35} % *1NEPHI-8-35 
		% ---
			\label{1NEPHI-8-35}

			And Laman and Lemuel partook not of the fruit, saith my father.

		% ---
		\paragraph{1 Nephi 8:36} % *1NEPHI-8-36 
		% ---
			\label{1NEPHI-8-36}

			And it came to pass that after my father had spoken
			all the words of his dream or vision, which were many,
			he said unto us,
			because of these things which he saw in a vision,
			he exceedingly feared for Laman and Lemuel.
			Yea, he feared lest they should be cast off from the presence of the Lord.

		% ---
		\paragraph{1 Nephi 8:37} % *1NEPHI-8-37 
		% ---
			\label{1NEPHI-8-37}

			And he did exhort them then with all the feeling of a tender parent
			that they would hearken to his words,
			in that perhaps the Lord would be merciful to them and not cast them off.
			Yea, my father did preach unto them.

		% ---
		\paragraph{1 Nephi 8:38} % *1NEPHI-8-38 
		% ---
			\label{1NEPHI-8-38}

			And after that he had preached unto them
			and also prophesied unto them of many things,
			he bade them to keep the commandments of the Lord.
			And he did cease speaking unto them.

	% -----
	\subsubsection{1 Nephi 9} % *1NEPHI-9 
	% -----
		\label{1NEPHI-9}
		
		% ---
		\paragraph{1 Nephi 9:1} % *1NEPHI-9-1 
		% ---
			\label{1NEPHI-9-1}

			And all these things did my father see and hear and speak
			as he dwelt in a tent in the valley of Lemuel,
			and also a great many more things which cannot be written upon these plates.

		% ---
		\paragraph{1 Nephi 9:2} % *1NEPHI-9-2 
		% ---
			\label{1NEPHI-9-2}

			And now as I have spoken concerning these plates,
			behold, they are not the plates
			upon which I make a full account of the history of my people,
			for the plates upon which I make a full account of my people
			I have given the name of Nephi;
			wherefore they are called the plates of Nephi after mine own name.
			And these plates also%
				\footnote{%
					% \label{1NEPHI-9-2-NOTE-plates}%
					Needlessly confusing. Why give them both the same name? It doesn't make a lot of sense.
					}
			are called the plates of Nephi.

		% ---
		\paragraph{1 Nephi 9:3} % *1NEPHI-9-3 
		% ---
			\label{1NEPHI-9-3}

			Nevertheless I have received a commandment of the Lord
			that I should make these plates for the special purpose
			that there should be an account engraven of the ministry of my people.

		% ---
		\paragraph{1 Nephi 9:4} % *1NEPHI-9-4 
		% ---
			\label{1NEPHI-9-4}

			And upon the other plates should be engraven
			an account of the reigns of the kings and the wars and contentions of my people.
			Wherefore these plates are for the more part of the ministry,
			and the other plates are for the more part of the reigns of the kings
			and the wars and contentions of my people.

		% ---
		\paragraph{1 Nephi 9:5} % *1NEPHI-9-5 
		% ---
			\label{1NEPHI-9-5}

			Wherefore the Lord hath commanded me to make these plates
			for a wise purpose in him, which purpose I know not.

		% ---
		\paragraph{1 Nephi 9:6} % *1NEPHI-9-6 
		% ---
			\label{1NEPHI-9-6}

			But the Lord knoweth all things from the beginning.
			Wherefore he prepareth a way
			to accomplish all his works among the children of men.
			For behold, he hath all power unto the fulfilling of all his words.
			And thus it is.
			Amen.

	% -----
	\subsubsection{1 Nephi 10} % *1NEPHI-10 
	% -----
		\label{1NEPHI-10}
		
		% ---
		\paragraph{1 Nephi 10:1} % *1NEPHI-10-1 
		% ---
			\label{1NEPHI-10-1}

			And now I Nephi proceed to give an account upon these plates
			of my proceedings and my reign and ministry.
			Wherefore to proceed with mine account,
			I must speak somewhat of the things of my father and also of my brethren.

		% ---
		\paragraph{1 Nephi 10:2} % *1NEPHI-10-2 
		% ---
			\label{1NEPHI-10-2}

			For behold, it came to pass that after my father had made an end
			of speaking the words of his dream and also of exhorting them to all diligence,
			he spake unto them concerning the Jews,

		% ---
		\paragraph{1 Nephi 10:3} % *1NEPHI-10-3 
		% ---
			\label{1NEPHI-10-3}

			how that after they were destroyed
			---yea, even that great city Jerusalem---
			and that many were carried away captive into Babylon,
			that according to the own due time of the Lord they should return again
			---yea, even be brought back out of captivity---
			and after that they are brought back out of captivity,
			to possess again their land of inheritance

		% ---
		\paragraph{1 Nephi 10:4} % *1NEPHI-10-4 
		% ---
			\label{1NEPHI-10-4}

			---yea, even six hundred years from the time that my father left Jerusalem---
			a prophet would the Lord God raise up among the Jews,
			yea, even a Messiah, or in other words, a Savior of the world.

		% ---
		\paragraph{1 Nephi 10:5} % *1NEPHI-10-5 
		% ---
			\label{1NEPHI-10-5}

			And he also spake concerning the prophets,
			how great a number had testified of these things
			concerning this Messiah of which he had spoken,
			or this Redeemer of the world.%
				\footnote{%
					\label{1NEPHI-10-5-NOTE-or}%
					In the previous verse Nephi mentions that his father is teaching about a ``Messiah,'' and then he glosses that term as ``Savior of the world.'' Then in this verse he does the same thing---mentions the title ``Messiah'' and then glosses it immediately as ``Redeemer of the world.'' Why does he do that? Nephi would have been familiar with the term \he{mA+siy/a.h}, and perhaps based on \hyperref[MESSIAH-HEBREW]{these meanings}, would have thought of that title as describing a king, priest, or prophet. But although Christ is all three of those things, he also plays a more fundamental role in the plan of salvation than any other king, priest, or prophet, and so it seems that Lehi wanted to emphasize that the ``Messiah'' or \he{mA+siy/a.h} he was referring to was not just some person or even some great person, but the figure who would actually bring salvation to people through his actions. This also makes sense of the next verse, verse six, where he doesn't just use a title but gives a real description of why the figure would be important. Lehi was imparting new information here, and giving everyone in his family a new way of seeing their relationship to God and the plan of salvation. He was using a familiar term in a new way and needed to explain why. (Note that something like this occurs again in \hyperref[1NEPHI-10-14]{verse 14} and in \hyperref[1NEPHI-10-17]{verse 17}. He's sharing his new understanding of the Messiah.)
					}

		% ---
		\paragraph{1 Nephi 10:6} % *1NEPHI-10-6 
		% ---
			\label{1NEPHI-10-6}

			Wherefore all mankind was in a lost and in a fallen state and ever would be
			save they should rely on this Redeemer.

		% ---
		\paragraph{1 Nephi 10:7} % *1NEPHI-10-7 
		% ---
			\label{1NEPHI-10-7}

			And he spake also concerning a prophet
			which should come before the Messiah to prepare the way of the Lord,

		% ---
		\paragraph{1 Nephi 10:8} % *1NEPHI-10-8 
		% ---
			\label{1NEPHI-10-8}

			yea, even he should go forth and cry in the wilderness:
			Prepare ye the way of the Lord and make his paths straight,
			for there standeth one among you whom ye know not
			---and he is mightier than I---
			whose shoe’s latchet I am not worthy to unloose.
			And much spake my father concerning this thing.

		% ---
		\paragraph{1 Nephi 10:9} % *1NEPHI-10-9 
		% ---
			\label{1NEPHI-10-9}

			And my father saith that he should baptize in Bethabara beyond Jordan.
			And he also spake that he should baptize with water,
			yea, even that he should baptize the Messiah with water.

		% ---
		\paragraph{1 Nephi 10:10} % *1NEPHI-10-10 
		% ---
			\label{1NEPHI-10-10}

			And after that he had baptized the Messiah with water
			he should behold and bear record that he had baptized the Lamb of God,
			which should take away the sin of the world.

		% ---
		\paragraph{1 Nephi 10:11} % *1NEPHI-10-11 
		% ---
			\label{1NEPHI-10-11}

			And it came to pass that after my father had spoken these words,
			he spake unto my brethren concerning the gospel
			which should be preached among the Jews,
			and also concerning the dwindling of the Jews in unbelief.
			And after that they had slain the Messiah which should come---
			and after that he had been slain,
			he should rise from the dead
			and should make himself manifest by the Holy Ghost unto the Gentiles.%
				\footnote{%
					% \label{1NEPHI-10-11-NOTE-manifest}%
					See \hyperref[2NEPHI-26-13]{2 Nephi 26:13}.
					}

		% ---
		\paragraph{1 Nephi 10:12} % *1NEPHI-10-12 
		% ---
			\label{1NEPHI-10-12}

			Yea, even my father spake much concerning the Gentiles
			and also concerning the house of Israel,
			that they should be compared like unto an olive tree
			whose branches should be broken off
			and should be scattered upon all the face of the earth.

		% ---
		\paragraph{1 Nephi 10:13} % *1NEPHI-10-13 
		% ---
			\label{1NEPHI-10-13}

			Wherefore he said it must needs be
			that we should be led with one accord into the land of promise,
			unto the fulfilling of the word of the Lord
			that we should be scattered upon all the face of the earth.

		% ---
		\paragraph{1 Nephi 10:14} % *1NEPHI-10-14 
		% ---
			\label{1NEPHI-10-14}

			And after that the house of Israel should be scattered,
			they should be gathered together again,
			or in fine, that after the Gentiles had received the fullness of the gospel,
			the natural branches of the olive tree or the remnants of the house of Israel
			should be grafted in or come to the knowledge of the true Messiah,
			their Lord and their Redeemer.

		% ---
		\paragraph{1 Nephi 10:15} % *1NEPHI-10-15 
		% ---
			\label{1NEPHI-10-15}

			And after this manner of language did my father prophesy
			and speak unto my brethren,
			and also many more things which I do not write in this book;
			for I have written as many of them as were expedient for me in mine other book.

		% ---
		\paragraph{1 Nephi 10:16} % *1NEPHI-10-16 
		% ---
			\label{1NEPHI-10-16}

			And all these things of which I have spoken was done
			as my father dwelt in a tent in the valley of Lemuel.

		% ---
		\paragraph{1 Nephi 10:17} % *1NEPHI-10-17 
		% ---
			\label{1NEPHI-10-17}

			And it came to pass that after I Nephi having heard all the words of my father
			concerning the things which he saw in a vision
			and also the things which he spake by the power of the Holy Ghost,
			which power he received by faith on the Son of God
			---and the Son of God was the Messiah which should come---
			and it came to pass that I Nephi was desirous also
			that I might see and hear and know of these things by the power of the Holy Ghost,
			which is the gift of God unto all those who diligently seek him%
				\footnote{%
					\label{1NEPHI-10-17-NOTE-gift}%
					Here we read that the Holy Ghost (or the power of the Holy Ghost) is the ``gift of God unto all those who diligently seek him.'' This statement is made without qualification and so therefore applies to everyone, everywhere, in a similar way to the promise made in \hyperref[MORMON-9-21]{Mormon 9:21}. \textit{Anyone} can get the Holy Ghost when they seek God; and if we broaden the possibility of what to seek by including just truth in general, then the Holy Ghost is going to be poured out upon many different people in many different circumstances.
					}
			as well in times of old
			as in the time that he should manifest himself unto the children of men,

		% ---
		\paragraph{1 Nephi 10:18} % *1NEPHI-10-18 
		% ---
			\label{1NEPHI-10-18}

			for he is the same yesterday and today and forever.
			And the way is prepared for all men from the foundation of the world
			if it so be that they repent and come unto him.

		% ---
		\paragraph{1 Nephi 10:19} % *1NEPHI-10-19 
		% ---
			\label{1NEPHI-10-19}

			For he that diligently seeketh shall find,
			and the mysteries of God shall be unfolded%
				\footnote{%
					% \label{1NEPHI-10-19-NOTE-unfolded}%
					This is the first use of ``unfold'' in the BOM. Unfortunately this key word does not appear at all in the NT. Ten times in the BOM, seven in the DC. As I wrote in a handout to my JS class in fall 2020: ``When you unfold a piece of cloth, parts that were once visible pass out of sight, and other parts become visible that you couldn't see before; in addition, when an entire cloth is unfolded, it is not possible to view both its sides at the same time (without a camera, or something like that).'' Scriptures that elaborate on this point include \hyperref[JACOB-4-8]{Jacob 4:8}, 
					}
			to them by the power of the Holy Ghost
			as well in this time as in times of old
			and as well in times of old as in times to come;
			\hyperref[WHEREFORE]{wherefore} the course of the Lord is one eternal round.%
			\footnote{%
				% \label{1NEPHI-10-19-NOTE-round}%
				Obligatory mention of how ``round'' means ``the step of a ladder''; see \hyperref[ROUND-1828]{here}, and also the phrase ``\hyperref[ROUND-PHRASES]{one eternal round}.''
				} 

		% ---
		\paragraph{1 Nephi 10:20} % *1NEPHI-10-20 
		% ---
			\label{1NEPHI-10-20}

			Therefore remember, O man:
			for all thy doings thou shalt be brought into judgment.

		% ---
		\paragraph{1 Nephi 10:21} % *1NEPHI-10-21 
		% ---
			\label{1NEPHI-10-21}

			Wherefore if ye have sought to do wickedly in the days of your probation,
			then ye are found unclean before the judgment seat of God.
			And no unclean thing can dwell with God;
			wherefore ye must be cast off forever.

		% ---
		\paragraph{1 Nephi 10:22} % *1NEPHI-10-22 
		% ---
			\label{1NEPHI-10-22}

			And the Holy Ghost giveth authority
			that I should speak these things and deny them not.

	% -----
	\subsubsection{1 Nephi 11} % *1NEPHI-11 
	% -----
		\label{1NEPHI-11}
		
		% ---
		\paragraph{1 Nephi 11:1} % *1NEPHI-11-1 
		% ---
			\label{1NEPHI-11-1}

			For it came to pass that after I had desired
			to know the things that my father had seen,
			and believing that the Lord was able to make them known unto me,
			wherefore as I sat pondering in mine heart,
			I was caught away in the Spirit of the Lord,
			yea, into an exceeding high mountain,
			a mountain which I never had before seen
			and upon which I never had before sat my foot.

		% ---
		\paragraph{1 Nephi 11:2} % *1NEPHI-11-2 
		% ---
			\label{1NEPHI-11-2}

			And the Spirit saith unto me:
			Behold, what desirest thou?

		% ---
		\paragraph{1 Nephi 11:3} % *1NEPHI-11-3 
		% ---
			\label{1NEPHI-11-3}

			And I saith:
			I desire to behold the things which my father saw.

		% ---
		\paragraph{1 Nephi 11:4} % *1NEPHI-11-4 
		% ---
			\label{1NEPHI-11-4}

			And the Spirit saith unto me:
			Believest thou that thy father saw the tree of which he hath spoken?

		% ---
		\paragraph{1 Nephi 11:5} % *1NEPHI-11-5 
		% ---
			\label{1NEPHI-11-5}

			And I said:
			Yea, thou knowest that I believe all the words of my father.

		% ---
		\paragraph{1 Nephi 11:6} % *1NEPHI-11-6 
		% ---
			\label{1NEPHI-11-6}

			And when I had spake these words,
			the Spirit cried with a loud voice, saying:
			Hosanna to the Lord, the Most High God,
			for he is God over all the earth, yea, even above all.
			And blessed art thou Nephi because thou believest in the Son of the Most High;
			wherefore thou shalt behold the things which thou hast desired.

		% ---
		\paragraph{1 Nephi 11:7} % *1NEPHI-11-7 
		% ---
			\label{1NEPHI-11-7}

			And behold, this thing shall be given unto thee for a sign,
			that after thou hast beheld the tree
			which bare the fruit of which thy father tasted,
			thou shalt also behold a man descending out of heaven%
				\footnote{%
					\label{1NEPHI-11-7-NOTE-man}%
					At least somewhat similar to what Lehi sees in \hyperref[1NEPHI-1-9]{1 Nephi 1:9}. I ended up concluding that the ``One'' spoken of in that verse is \hyperref[1NEPHI-1-9-NOTE-one]{probably Christ}, but here it's going to be clearer that it actually is Christ.
					}
			and him shall ye witness.
			And after that ye shall have witnessed him,
			ye shall bear record that it is the Son of God.

		% ---
		\paragraph{1 Nephi 11:8} % *1NEPHI-11-8 
		% ---
			\label{1NEPHI-11-8}

			And it came to pass that the Spirit saith unto me: Look!
			And I looked and beheld a tree,
			and it was like unto the tree which my father had seen.
			And the beauty thereof was far beyond, yea, exceeding of all beauty,
			and the whiteness thereof did exceed the whiteness of the driven snow.%
				\footnote{%
					\label{1NEPHI-11-8-NOTE-snow}%
					When I read this I wondered, how does Nephi know about snow? Obviously someone could have just told him about it, but it could be also that this is a way of describing it that he picked up after they arrived in the promised land. It does snow in Jerusalem too, though. It's rare but it happens. I read this verse on 21 February 2021, and two days before that it snowed a fair bit, more than enough for people to play in and make snowmen. So Nephi could have experienced snow for himself by this time. What's interesting is that this is the \textit{only} occurrence of the word ``snow'' in the entire BOM. There are three occurrences in the NT, all describing the transfigured or resurrected Christ, and then twenty-two uses in the OT (four in the DC). So if the NT uses are any guide, as they sometimes are for the BOM, then by ``snow'' here we're supposed to understand something celestial or celestially pure, maybe. Yet another interesting thing about this is why Nephi describes it as whiter than the \textit{driven} snow. Does snow become whiter when it's driven, when it's flurrying around you? I'm not sure, but it does change your perception of it---there's white flashing everywhere all around ou. Maybe that is what Nephi intends here? The word ``driven'' appears many times in the scriptures but in my brief and not exhaustive looking I did not find another use of it as an attributive adjective like this; the other uses I saw are either predicate adjectives or the participles in perfect-tense verbs.
					}

		% ---
		\paragraph{1 Nephi 11:9} % *1NEPHI-11-9 
		% ---
			\label{1NEPHI-11-9}

			And it came to pass that after that I had seen the tree,
			I said unto the Spirit:
			I behold thou hast shewn unto me the tree
			which is most precious above all.%
				\footnote{%
					\label{1NEPHI-11-9-NOTE-precious}%
					See \hyperref[1NEPHI-15-36]{1 Nephi 15:36} and \hyperref[HELAMAN-5-8]{Helaman 5:8}. Definitions \#2 and \#3 of \hyperref[PRECIOUS-1828]{precious} in the 1828 dictionary refer to a thing's objective and subjective value, respectively. Is it right to say that the tree has both kinds of value---it is objectively valuable, but also valuable to any given individual who experiences it? I'm not sure. 
					}

		% ---
		\paragraph{1 Nephi 11:10} % *1NEPHI-11-10 
		% ---
			\label{1NEPHI-11-10}

			And he saith unto me:
			What desirest thou?

		% ---
		\paragraph{1 Nephi 11:11} % *1NEPHI-11-11 
		% ---
			\label{1NEPHI-11-11}

			And I said unto him:
			To know the interpretation thereof.
			For I spake unto him as a man speaketh,
			for I beheld that he was in the form of a man,%
				\footnote{%
					% \label{}%
					See \hyperref[DC77-2]{DC 77:2}. That verse and this one here are some of the strongest scriptural evidences that spirits are already ``humanized,'' or take a human form, prior to actually inhabiting a human body.
					}
			yet nevertheless I knew that it was the Spirit of the Lord,
			and he spake unto me as a man speaketh with another.

		% ---
		\paragraph{1 Nephi 11:12} % *1NEPHI-11-12 
		% ---
			\label{1NEPHI-11-12}

			And it came to pass that he said unto me: Look!
			And I looked as if to look upon him and I saw him not,
			for he had gone from before my presence.

		% ---
		\paragraph{1 Nephi 11:13} % *1NEPHI-11-13 
		% ---
			\label{1NEPHI-11-13}

			And it came to pass that I looked
			and beheld the great city Jerusalem and also other cities.
			And I beheld the city of Nazareth.%
				\footnote{%
					\label{1NEPHI-11-13-NOTE-nazareth}%
					Nazareth is about 90 miles by car from Jerusalem, and about 63 miles directly distant. That would probably be a journey of three days for Nephi if you were making good time and had good conditions.
					}
			And in the city of Nazareth I beheld a virgin,%
				\footnote{%
					\label{1NEPHI-11-13-NOTE-virgin}%
					This is the first occurrence out of seven in the BOM. Five are here in 1 Nephi 11, then one in the Isaiah chapters of 2 Nephi and one in Alma 7 (see \hyperref[VIRGIN-BOM]{here}). The word occurs in the NT fourteen times, and you can see those occurrences \hyperref[VIRGIN-NT-KJV]{here}. Every one of those is a translation of \gr{παρθένος}, which is a word that \hyperref[PARTHENOS]{focuses} on the fact that a younger woman has not had sex. I think it's a focus that is not very salutary to many nowadays, so we would have to understand more about earlier cultures to make sense of it. The word does also just mean ``young woman,'' but according to BDAG, the focus in NT-related literature is on the fact that she hasn't had sex. The primary definition in the \hyperref[VIRGIN-1828]{1828 dictionary} for the word also focuses on the sex aspect. What's curious to me here in 1 Nephi 11:13 is how Nephi knew the girl, who is Mary, hadn't had sex. It's not like you can just tell that from someone's behavior on the outside. It must have been something the Spirit gave him to understand, which itself seems a little odd, because isn't it a little private, whether she's had sex or not? In any case I think that Nephi learns this about her so that he can learn a further lesson about the ``condescension of God,'' or the miraculous manner in which Christ came to be on the Earth.
					
					The other option is that Nephi himself is using the term to refer to her just as a younger woman, which would be a legitimate use of it. In that case he wouldn't have any special access to her sexual history, and I think that's plausible too.
					}
			and she was exceeding fair and white.%
				\footnote{%
					\label{1NEPHI-11-13-NOTE-white}%
					I'm assuming he means ``white'' as a way of describing her purity or holiness, in the way he described the tree above in \hyperref[1NEPHI-11-8]{verse 8}. I don't see how he could mean that she had light-colored skin.
					}

		% ---
		\paragraph{1 Nephi 11:14} % *1NEPHI-11-14 
		% ---
			\label{1NEPHI-11-14}

			And it came to pass that I saw the heavens open,
			and an angel came down and stood before me.
			And he saith unto me:
			Nephi, what beholdest thou?

		% ---
		\paragraph{1 Nephi 11:15} % *1NEPHI-11-15 
		% ---
			\label{1NEPHI-11-15}

			And I saith unto him:
			A virgin most beautiful and fair above all other virgins.

		% ---
		\paragraph{1 Nephi 11:16} % *1NEPHI-11-16 
		% ---
			\label{1NEPHI-11-16}

			And he saith unto me:
			Knowest thou the condescension of God?%
				\footnote{%
					\label{1NEPHI-11-16-NOTE-condescension}%
					See \hyperref[1NEPHI-11-26]{verse 26} and \hyperref[CONDESCENSION-1nephi]{this study} for more on ``condescension'' here.
					}

		% ---
		\paragraph{1 Nephi 11:17} % *1NEPHI-11-17 
		% ---
			\label{1NEPHI-11-17}

			And I said unto him:
			I know that he loveth his children;
			nevertheless I do not know the meaning of all things.

		% ---
		\paragraph{1 Nephi 11:18} % *1NEPHI-11-18 
		% ---
			\label{1NEPHI-11-18}

			And he said unto me:
			Behold, the virgin which thou seest is the mother of God
			after the manner of the flesh.

		% ---
		\paragraph{1 Nephi 11:19} % *1NEPHI-11-19 
		% ---
			\label{1NEPHI-11-19}

			And it came to pass that I beheld that she was carried away in the spirit.
			And after that she had been carried away in the spirit for the space of a time,
			the angel spake unto me, saying: Look!

		% ---
		\paragraph{1 Nephi 11:20} % *1NEPHI-11-20 
		% ---
			\label{1NEPHI-11-20}

			And I looked and beheld the virgin%
				\footnote{%
					% \label{}%
					Even after he sees the baby, Nephi still refers to Mary as a ``virgin.'' Two thoughts: first, it's possible that he is not using the term ``virgin'' to refer to her as someone who hasn't had intercourse (see the note to verse 13), and he's just using it to mean ``girl'' or ``young lady.'' Second, if the term does have a sexual connotation here, then it would speak against BY's view that God had intercourse with Mary and impregnated her that way.
					}
			again,
			bearing a child in her arms.

		% ---
		\paragraph{1 Nephi 11:21} % *1NEPHI-11-21 
		% ---
			\label{1NEPHI-11-21}

			And the angel said unto me:
			Behold the Lamb of God, yea, even the Eternal Father.
			Knowest thou the meaning of the tree which thy father saw?

		% ---
		\paragraph{1 Nephi 11:22} % *1NEPHI-11-22 
		% ---
			\label{1NEPHI-11-22}

			And I answered him, saying:
			Yea, it is the love of God,
			which sheddeth itself \hyperref[ABROAD]{abroad}%
				\footnote{%
					% \label{1NEPHI-11-22-NOTE-shed}%
					This part of the verse is a reference to \hyperref[ROMANS-5-5]{Romans 5:5}. See \hyperref[SHED]{Shed} and especially \hyperref[SHED-abroad]{this study}.
					}
			in the hearts of the children of men;
			wherefore it is the most desirable above all things.

		% ---
		\paragraph{1 Nephi 11:23} % *1NEPHI-11-23 
		% ---
			\label{1NEPHI-11-23}

			And he spake unto me, saying:
			Yea, and the most joyous to the soul.

		% ---
		\paragraph{1 Nephi 11:24} % *1NEPHI-11-24 
		% ---
			\label{1NEPHI-11-24}

			And after that he had said these words,
			he said unto me: Look!
			And I looked and I beheld the Son of God
			a going forth among the children of men.
			And I saw many fall down at his feet and worship him.

		% ---
		\paragraph{1 Nephi 11:25} % *1NEPHI-11-25 
		% ---
			\label{1NEPHI-11-25}

			And it came to pass that I beheld
			that the rod of iron which my father had seen was the word of God,
			which led to the fountain of living waters%
				\footnote{%
					\label{1NEPHI-11-25-NOTE-waters}%
					Both the OT and the NT use the phrase ``living waters,'' and the NT also includes the similar phrase ``waters of life.'' See \hyperref[WATER-PHRASES]{this list}.
					}
			or to the tree of life,
			which waters are a representation of the love of God.
			And I also beheld that the tree of life was a representation of the love of God.

		% ---
		\paragraph{1 Nephi 11:26} % *1NEPHI-11-26 
		% ---
			\label{1NEPHI-11-26}

			And the angel said unto me again:
			Look and behold the condescension of God!

		% ---
		\paragraph{1 Nephi 11:27} % *1NEPHI-11-27 
		% ---
			\label{1NEPHI-11-27}

			And I looked and beheld the Redeemer of the world,
			of which my father had spoken.
			And I also beheld the prophet which should prepare the way before him.
			And the Lamb of God went forth and was baptized of him.
			And after that he was baptized,
			I beheld the heavens open
			and the Holy Ghost came down out of heaven
			and abode upon him in the form of a dove.

		% ---
		\paragraph{1 Nephi 11:28} % *1NEPHI-11-28 
		% ---
			\label{1NEPHI-11-28}

			And I beheld that he went forth
			ministering unto the people in power and great glory,
			and the multitudes were gathered together to hear him.
			And I beheld that they cast him out from among them.

		% ---
		\paragraph{1 Nephi 11:29} % *1NEPHI-11-29 
		% ---
			\label{1NEPHI-11-29}

			And I also beheld twelve others following him.
			And it came to pass that they were carried away in the spirit from before my face,
			that I saw them not.

		% ---
		\paragraph{1 Nephi 11:30} % *1NEPHI-11-30 
		% ---
			\label{1NEPHI-11-30}

			And it came to pass that the angel spake unto me, saying: Look!
			And I looked and I beheld the heavens open again.
			And I saw angels descending upon the children of men,
			and they did minister unto them.

		% ---
		\paragraph{1 Nephi 11:31} % *1NEPHI-11-31 
		% ---
			\label{1NEPHI-11-31}

			And he spake unto me again, saying: Look!
			And I looked and I beheld the Lamb of God
			going forth among the children of men.
			And I beheld multitudes of people which were sick
			and which were afflicted of all manner of diseases
			and with devils and unclean spirits
			---and the angel spake and shewed all these things unto me---
			and they were healed by the power of the Lamb of God,
			and the devils and the unclean spirits were cast out.

		% ---
		\paragraph{1 Nephi 11:32} % *1NEPHI-11-32 
		% ---
			\label{1NEPHI-11-32}

			And it came to pass that the angel spake unto me again, saying: Look!
			And I looked and beheld the Lamb of God,
			that he was taken by the people,
			yea, the everlasting God was judged of the world.
			And I saw and bare record.

		% ---
		\paragraph{1 Nephi 11:33} % *1NEPHI-11-33 
		% ---
			\label{1NEPHI-11-33}

			And I Nephi saw that he was lifted up upon the cross
			and slain for the sins of the world.

		% ---
		\paragraph{1 Nephi 11:34} % *1NEPHI-11-34 
		% ---
			\label{1NEPHI-11-34}

			And after that he was slain, I saw the multitudes of the earth,
			that they were gathered together to fight against the apostles of the Lamb,
			for thus were the twelve called by the angel of the Lord.

		% ---
		\paragraph{1 Nephi 11:35} % *1NEPHI-11-35 
		% ---
			\label{1NEPHI-11-35}

			And the multitude of the earth was gathered together,
			and I beheld that they were in a large and spacious building,
			like unto the building which my father saw.
			And the angel of the Lord spake unto me, saying:
			Behold the world and the wisdom thereof;
			yea, behold, the house of Israel hath gathered together
			to fight against the twelve apostles of the Lamb.

		% ---
		\paragraph{1 Nephi 11:36} % *1NEPHI-11-36 
		% ---
			\label{1NEPHI-11-36}

			And it came to pass that I saw and bare record
			that the great and spacious building was the pride of the world;
			and the fall thereof was exceeding great.
			And the angel of the Lord spake unto me, saying:
			Thus shall be the destruction of all nations, kindreds, tongues, and people
			that shall fight against the twelve apostles of the Lamb.

	% -----
	\subsubsection{1 Nephi 12} % *1NEPHI-12 
	% -----
		\label{1NEPHI-12}
		
		% ---
		\paragraph{1 Nephi 12:1} % *1NEPHI-12-1 
		% ---
			\label{1NEPHI-12-1}

			And it came to pass that the angel said unto me:
			Look and behold thy seed and also the seed of thy brethren.
			And I looked and beheld the land of promise.
			And I beheld multitudes of people,
			yea, even as it were in number as many as the sand of the sea.

		% ---
		\paragraph{1 Nephi 12:2} % *1NEPHI-12-2 
		% ---
			\label{1NEPHI-12-2}

			And it came to pass that I beheld multitudes gathered together
			to battle one against the other.
			And I beheld wars and rumors of wars
			and great slaughters with the sword among my people.

		% ---
		\paragraph{1 Nephi 12:3} % *1NEPHI-12-3 
		% ---
			\label{1NEPHI-12-3}

			And it came to pass that I beheld many generations pass away
			after the manner of wars and contentions in the land.
			And I beheld many cities,
			yea, even that I did not number them.

		% ---
		\paragraph{1 Nephi 12:4} % *1NEPHI-12-4 
		% ---
			\label{1NEPHI-12-4}

			And it came to pass that I saw a mist of darkness
			on the face of the land of promise.
			And I saw lightnings and I heard thunderings and earthquakes
			and all manner of tumultuous noises.
			And I saw the earth that it rent the rocks,
			and I saw mountains tumbling into pieces,
			and I saw the plains of the earth that they were broken up.
			And I saw many cities that they were sunk,
			and I saw many that they were burnt with fire,
			and I saw many that they did tumble to the earth because of the quaking thereof.

		% ---
		\paragraph{1 Nephi 12:5} % *1NEPHI-12-5 
		% ---
			\label{1NEPHI-12-5}

			And it came to pass that after I saw these things,
			I saw the vapor of darkness,
			that it passed from off the face of the earth.
			And behold, I saw the multitudes which had not fallen
			because of the great and terrible judgments of the Lord.

		% ---
		\paragraph{1 Nephi 12:6} % *1NEPHI-12-6 
		% ---
			\label{1NEPHI-12-6}

			And I saw the heavens open
			and the Lamb of God descending out of heaven.
			And he came down and he shewed himself unto them.

		% ---
		\paragraph{1 Nephi 12:7} % *1NEPHI-12-7 
		% ---
			\label{1NEPHI-12-7}

			And I also saw and bare record
			that the Holy Ghost fell upon twelve others,
			and they were ordained of God and chosen.

		% ---
		\paragraph{1 Nephi 12:8} % *1NEPHI-12-8 
		% ---
			\label{1NEPHI-12-8}

			And the angel spake unto me, saying:
			Behold the twelve disciples of the Lamb
			which are chosen to minister unto thy seed.

		% ---
		\paragraph{1 Nephi 12:9} % *1NEPHI-12-9 
		% ---
			\label{1NEPHI-12-9}

			And he saith unto me:
			Thou rememberest the twelve apostles of the Lamb.
			Behold, they are they which shall judge the twelve tribes of Israel;
			wherefore the twelve ministers of thy seed shall be judged of them,
			for ye are of the house of Israel.

		% ---
		\paragraph{1 Nephi 12:10} % *1NEPHI-12-10 
		% ---
			\label{1NEPHI-12-10}

			And these twelve ministers which thou beholdest shall judge thy seed.
			And behold, they are righteous forever,
			for because of their faith in the Lamb of God,
			their garments are made white in his blood.%
				\footnote{%
					\label{1NEPHI-12-10-NOTE-righteous}%
					This is an interesting explanation of the way in which the twelve disciples are ``righteous'' forever---they are called so here because their garments are cleansed in the blood of Christ. What's interesting is that you would expect them to be ``clean'' or ``pure'' or something like that as a result of that washing, because you normally think of ``righteous'' as a quality of \textit{action}, or a quality that you show in your actions. You don't normally think of being ``righteous'' as the result of the cleansing process but that is how it's talked about here. See the next verse for more.
					}

		% ---
		\paragraph{1 Nephi 12:11} % *1NEPHI-12-11 
		% ---
			\label{1NEPHI-12-11}

			And the angel saith unto me: Look!
			And I looked and beheld three generations did pass away in righteousness,
			and their garments were white, even like unto the Lamb of God.
			And the angel said unto me:
			These are made white in the blood of the Lamb
			because of their faith in him.

		% ---
		\paragraph{1 Nephi 12:12} % *1NEPHI-12-12 
		% ---
			\label{1NEPHI-12-12}

			And I Nephi also saw many of the fourth generation
			which did pass away in righteousness.

		% ---
		\paragraph{1 Nephi 12:13} % *1NEPHI-12-13 
		% ---
			\label{1NEPHI-12-13}

			And it came to pass that I saw the multitudes of the earth gathered together.

		% ---
		\paragraph{1 Nephi 12:14} % *1NEPHI-12-14 
		% ---
			\label{1NEPHI-12-14}

			And the angel said unto me:
			Behold thy seed and also the seed of thy brethren.

		% ---
		\paragraph{1 Nephi 12:15} % *1NEPHI-12-15 
		% ---
			\label{1NEPHI-12-15}

			And it came to pass that I looked and beheld the people of my seed
			gathered together in multitudes against the seed of my brethren;
			and they were gathered together to battle.

		% ---
		\paragraph{1 Nephi 12:16} % *1NEPHI-12-16 
		% ---
			\label{1NEPHI-12-16}

			And the angel spake unto me, saying:
			Behold the fountain of filthy water which thy father saw,
			yea, even the river of which he spake;
			and the depths thereof are the depths of hell.

		% ---
		\paragraph{1 Nephi 12:17} % *1NEPHI-12-17 
		% ---
			\label{1NEPHI-12-17}

			And the mists of darkness are the temptations of the devil,
			which blindeth the eyes and hardeneth the hearts of the children of men
			and leadeth them away into broad roads
			that they perish and are lost.

		% ---
		\paragraph{1 Nephi 12:18} % *1NEPHI-12-18 
		% ---
			\label{1NEPHI-12-18}

			And the large and spacious building%
			\footnote{%
				% \label{}%
				Notice that the building is \textit{not} the ``world,'' or some other way of referring to a group of unrighteous people as a whole. It's an attitude that characterizes you and your actions, not a group.
				}
			which thy father saw
			is
				\footnote{%
					% \label{}%
					The three biblical precedents for ``vain imaginations'' are \hyperref[PSALMS-2-1]{Psalms 2:1}, \hyperref[ACTS-4-25]{Acts 4:25}, and \hyperref[ROMANS-1-21]{Romans 1:21}.
					}%
			\hyperref[VANITY-1828]{vain}%
				\footnote{%
					% \label{}%
					See the 1828 definitions for \hyperref[VANITY-1828]{vain}; the word doesn't mean ``prideful'' or ``conceited'' here, both because those aren't the primary meanings of the word in the 1828, and also because the text immediately adds ``pride.'' So ``vain'' here means ``empty, worthless, having no value, fruitless, ineffectual, useless.''
					}
			imaginations%
				\footnote{%
					\label{1NEPHI-12-18-NOTE-imagination}%
					See the 1828 definitions of ``imagination'' \hyperref[IMAGINATION-1828]{here}; the two relevant ones I think are \#2 and \#3: ``Contrivance; scheme formed in the mind; device'' and ``Conceit; an unsolid or fanciful opinion.''
					} 
			and the pride of the children of men.
			And a great and a terrible gulf divideth them,
			yea, even the sword of the justice of the Eternal God and Jesus Christ,
			which is the Lamb of God,
			of whom the Holy Ghost beareth record
			from the beginning of the world until this time
			and from this time henceforth and forever.

		% ---
		\paragraph{1 Nephi 12:19} % *1NEPHI-12-19 
		% ---
			\label{1NEPHI-12-19}

			And while the angel spake these words,
			I beheld and saw that the seed of my brethren did contend against my seed,
			according to the word of the angel.
			And because of the pride of my seed and the temptations of the devil,
			I beheld that the seed of my brethren did overpower the people of my seed.

		% ---
		\paragraph{1 Nephi 12:20} % *1NEPHI-12-20 
		% ---
			\label{1NEPHI-12-20}

			And it came to pass that I beheld and saw the people of the seed of my brethren,
			that they had overcome my seed.
			And they went forth in multitudes upon the face of the land;

		% ---
		\paragraph{1 Nephi 12:21} % *1NEPHI-12-21 
		% ---
			\label{1NEPHI-12-21}

			and I saw them gathered together in multitudes.
			And I saw wars and rumors of wars among them,
			and in wars and rumors of wars I saw many generations pass away.

		% ---
		\paragraph{1 Nephi 12:22} % *1NEPHI-12-22 
		% ---
			\label{1NEPHI-12-22}

			And the angel said unto me:
			Behold, these shall dwindle in unbelief.

		% ---
		\paragraph{1 Nephi 12:23} % *1NEPHI-12-23 
		% ---
			\label{1NEPHI-12-23}

			And it came to pass that I beheld that after they had dwindled in unbelief,
			they became a dark and loathsome and a filthy people,
			full of idleness and all manner of abominations.

	% -----
	\subsubsection{1 Nephi 13} % *1NEPHI-13 
	% -----
		\label{1NEPHI-13}
		
		% ---
		\paragraph{1 Nephi 13:1} % *1NEPHI-13-1 
		% ---
			\label{1NEPHI-13-1}

			And it came to pass that the angel spake unto me, saying: Look!
			And I looked and beheld many nations and kingdoms.

		% ---
		\paragraph{1 Nephi 13:2} % *1NEPHI-13-2 
		% ---
			\label{1NEPHI-13-2}

			And the angel saith unto me:
			What beholdest thou?
			And I said:
			I behold many nations and kingdoms.

		% ---
		\paragraph{1 Nephi 13:3} % *1NEPHI-13-3 
		% ---
			\label{1NEPHI-13-3}

			And he saith unto me:
			These are the nations and kingdoms of the Gentiles.

		% ---
		\paragraph{1 Nephi 13:4} % *1NEPHI-13-4 
		% ---
			\label{1NEPHI-13-4}

			And it came to pass that I saw among the nations of the Gentiles
			the formation of a great church.

		% ---
		\paragraph{1 Nephi 13:5} % *1NEPHI-13-5 
		% ---
			\label{1NEPHI-13-5}

			And the angel said unto me:
			Behold the formation of a church
			which is most abominable above all other churches,
			which slayeth the saints of God,
			yea, and tortureth them and bindeth them down
			and yoketh them with a yoke of iron
			and bringeth them down into captivity.

		% ---
		\paragraph{1 Nephi 13:6} % *1NEPHI-13-6 
		% ---
			\label{1NEPHI-13-6}

			And it came to pass that I beheld this great and abominable church,
			and I saw the devil that he was the founder of it.

		% ---
		\paragraph{1 Nephi 13:7} % *1NEPHI-13-7 
		% ---
			\label{1NEPHI-13-7}

			And I also saw gold and silver and silks and scarlets
			and fine-twined linen and all manner of precious clothing,
			and I saw many harlots.

		% ---
		\paragraph{1 Nephi 13:8} % *1NEPHI-13-8 
		% ---
			\label{1NEPHI-13-8}

			And the angel spake unto me, saying:
			Behold, the gold and the silver and the silks and the scarlets
			and the fine-twined linen and the precious clothing and the harlots
			are the desires of this great and abominable church.

		% ---
		\paragraph{1 Nephi 13:9} % *1NEPHI-13-9 
		% ---
			\label{1NEPHI-13-9}

			And also for the praise of the world do they destroy the saints of God
			and bring them down into captivity.

		% ---
		\paragraph{1 Nephi 13:10} % *1NEPHI-13-10 
		% ---
			\label{1NEPHI-13-10}

			And it came to pass that I looked and beheld many waters,
			and they divided the Gentiles from the seed of my brethren.

		% ---
		\paragraph{1 Nephi 13:11} % *1NEPHI-13-11 
		% ---
			\label{1NEPHI-13-11}

			And it came to pass that the angel saith unto me:
			Behold, the wrath of God is upon the seed of thy brethren.

		% ---
		\paragraph{1 Nephi 13:12} % *1NEPHI-13-12 
		% ---
			\label{1NEPHI-13-12}

			And I looked and beheld a man among the Gentiles,
			which were separated from the seed of my brethren by the many waters.
			And I beheld the Spirit of God,
			that it came down and wrought upon the man,
			and he went forth upon the many waters,
			even unto the seed of my brethren, which were in the promised land.

		% ---
		\paragraph{1 Nephi 13:13} % *1NEPHI-13-13 
		% ---
			\label{1NEPHI-13-13}

			And it came to pass that I beheld the Spirit of God,
			that it wrought upon other Gentiles,
			and they went forth out of captivity upon the many waters.

		% ---
		\paragraph{1 Nephi 13:14} % *1NEPHI-13-14 
		% ---
			\label{1NEPHI-13-14}

			And it came to pass that I beheld
			many multitudes of the Gentiles upon the land of promise.
			And I beheld the wrath of God,
			that it was upon the seed of my brethren.
			And they were scattered before the Gentiles and they were smitten.

		% ---
		\paragraph{1 Nephi 13:15} % *1NEPHI-13-15 
		% ---
			\label{1NEPHI-13-15}

			And I beheld the Spirit of the Lord, that it was upon the Gentiles,
			that they did prosper and obtain the land for their inheritance.
			And I beheld that they were white and exceeding fair and beautiful,
			like unto my people before that they were slain.

		% ---
		\paragraph{1 Nephi 13:16} % *1NEPHI-13-16 
		% ---
			\label{1NEPHI-13-16}

			And it came to pass that I Nephi beheld that the Gentiles
			which had gone forth out of captivity
			did humble themselves before the Lord,
			and the power of the Lord was with them.

		% ---
		\paragraph{1 Nephi 13:17} % *1NEPHI-13-17 
		% ---
			\label{1NEPHI-13-17}

			And I beheld that their mother Gentiles was gathered together
			upon the waters and upon the land also,
			to battle against them.

		% ---
		\paragraph{1 Nephi 13:18} % *1NEPHI-13-18 
		% ---
			\label{1NEPHI-13-18}

			And I beheld that the power of God was with them,
			and also that the wrath of God was upon all those
			that were gathered together against them to battle.

		% ---
		\paragraph{1 Nephi 13:19} % *1NEPHI-13-19 
		% ---
			\label{1NEPHI-13-19}

			And I Nephi beheld that the Gentiles which had gone out of captivity
			were delivered by the power of God out of the hands of all other nations.

		% ---
		\paragraph{1 Nephi 13:20} % *1NEPHI-13-20 
		% ---
			\label{1NEPHI-13-20}

			And it came to pass that I Nephi beheld that they did prosper in the land.
			And I beheld a book and it was carried forth among them.

		% ---
		\paragraph{1 Nephi 13:21} % *1NEPHI-13-21 
		% ---
			\label{1NEPHI-13-21}

			And the angel saith unto me:
			Knowest thou the meaning of the book?

		% ---
		\paragraph{1 Nephi 13:22} % *1NEPHI-13-22 
		% ---
			\label{1NEPHI-13-22}

			And I saith:
			I know not.

		% ---
		\paragraph{1 Nephi 13:23} % *1NEPHI-13-23 
		% ---
			\label{1NEPHI-13-23}

			And he saith:
			Behold, it proceedeth out of the mouth of a Jew.
			And I Nephi beheld it.
			And he saith unto me:
			The book which thou beholdest is a record of the Jews,
			which contain the covenants of the Lord
			which he hath made unto the house of Israel.
			And it also containeth many of the prophecies of the holy prophets.
			And it is a record like unto the engravings
			which are upon the plates of brass,
			save there are not so many.
			Nevertheless they contain the covenants of the Lord
			which he hath made unto the house of Israel;
			wherefore they are of great worth unto the Gentiles.

		% ---
		\paragraph{1 Nephi 13:24} % *1NEPHI-13-24 
		% ---
			\label{1NEPHI-13-24}

			And the angel of the Lord said unto me:
			Thou hast beheld that the book proceeded forth from the mouth of a Jew.
			And when it proceeded forth from the mouth of a Jew,
			it contained the fullness of the gospel of the Lamb,
			of whom the twelve apostles bare record.
			And they bare record according to the truth which is in the Lamb of God.

		% ---
		\paragraph{1 Nephi 13:25} % *1NEPHI-13-25 
		% ---
			\label{1NEPHI-13-25}

			Wherefore these things go forth from the Jews in purity unto the Gentiles,
			according to the truth which is in God.

		% ---
		\paragraph{1 Nephi 13:26} % *1NEPHI-13-26 
		% ---
			\label{1NEPHI-13-26}

			And after that they go forth by the hand of the twelve apostles of the Lamb
			from the Jews unto the Gentiles,
			behold, after this thou seest the formation of that great and abominable church,
			which is the most abominable of all other churches.
			For behold, they have taken away from the gospel of the Lamb
			many parts which are plain and most precious;
			and also many covenants of the Lord have they taken away.

		% ---
		\paragraph{1 Nephi 13:27} % *1NEPHI-13-27 
		% ---
			\label{1NEPHI-13-27}

			And all this have they done that they might pervert the right ways of the Lord,
			that they might blind the eyes and harden the hearts of the children of men.

		% ---
		\paragraph{1 Nephi 13:28} % *1NEPHI-13-28 
		% ---
			\label{1NEPHI-13-28}

			Wherefore thou seest that after the book hath gone forth
			through the hands of the great and abominable church
			that there are many plain and most precious things taken away from the book,
			which is the book of the Lamb of God.

		% ---
		\paragraph{1 Nephi 13:29} % *1NEPHI-13-29 
		% ---
			\label{1NEPHI-13-29}

			And after that these plain and precious things were taken away,
			it goeth forth unto all the nations of the Gentiles.
			And after it goeth forth unto all the nations of the Gentiles,
			yea, even across the many waters---which thou hast seen---
			with the Gentiles which have gone forth out of captivity,
			and thou seest because of the many plain and precious things
			which have been taken out of the book,
			which were plain unto the understanding of the children of men
			according to the plainness which is in the Lamb of God---
			and because of these things which are taken away out of the gospel of the Lamb,
			an exceeding great many do stumble,
			yea, insomuch that Satan hath great power over them.

		% ---
		\paragraph{1 Nephi 13:30} % *1NEPHI-13-30 
		% ---
			\label{1NEPHI-13-30}

			Nevertheless thou beholdest that the Gentiles
			which have gone forth out of captivity
			and have been lifted up by the power of God
			above all other nations upon the face of the land
			which is choice above all other lands,
			which is the land which the Lord God hath covenanted with thy father
			that his seed should have for the land of their inheritance,
			wherefore thou seest that the Lord God will not suffer
			that the Gentiles will utterly destroy the mixture of thy seed
			which is among thy brethren.

		% ---
		\paragraph{1 Nephi 13:31} % *1NEPHI-13-31 
		% ---
			\label{1NEPHI-13-31}

			Neither will he suffer that the Gentiles shall destroy the seed of thy brethren.

		% ---
		\paragraph{1 Nephi 13:32} % *1NEPHI-13-32 
		% ---
			\label{1NEPHI-13-32}

			Neither will the Lord God suffer that the Gentiles shall forever remain
			in that state of awful wickedness which thou beholdest that they are in
			because of the plain and most precious parts of the gospel of the Lamb
			which hath been kept back by that abominable church,
			whose formation thou hast seen.

		% ---
		\paragraph{1 Nephi 13:33} % *1NEPHI-13-33 
		% ---
			\label{1NEPHI-13-33}

			Wherefore, saith the Lamb of God,
			I will be merciful unto the Gentiles,
			unto the visiting of the remnant of the house of Israel in great judgment.

		% ---
		\paragraph{1 Nephi 13:34} % *1NEPHI-13-34 
		% ---
			\label{1NEPHI-13-34}

			And it came to pass that the angel of the Lord spake unto me, saying:
			Behold, saith the Lamb of God,
			after that I have visited the remnant of the house of Israel
			---and this remnant of which I speak is the seed of thy father---
			wherefore after that I have visited them in judgment
			and smitten them by the hand of the Gentiles,
			and after that the Gentiles do stumble exceedingly
			because of the most plain and precious parts of the gospel of the Lamb
			which hath been kept back by that abominable church,
			which is the mother of harlots, saith the Lamb,
			wherefore I will be merciful unto the Gentiles in that day, saith the Lamb,
			insomuch that I will bring forth unto them in mine own power
			much of my gospel,%
				\footnote{%
					\label{1NEPHI-13-34-NOTE-much}%
					Interesting comment, to explicitly say that only ``much'' of the gospel will be brought forth. ``Much'' $\neq$ ``all,'' but this bit here isn't necessarily talking about the entire restoration.
					}
			which shall be plain and precious, saith the Lamb.

		% ---
		\paragraph{1 Nephi 13:35} % *1NEPHI-13-35 
		% ---
			\label{1NEPHI-13-35}

			For behold, saith the Lamb, I will manifest myself unto thy seed
			that they shall write many things which I shall minister unto them,
			which shall be plain and precious.
			And after that thy seed shall be destroyed and dwindle in unbelief,
			and also the seed of thy brethren,
			behold, these things shall be hid up to come forth unto the Gentiles
			by the gift and power of the Lamb.

		% ---
		\paragraph{1 Nephi 13:36} % *1NEPHI-13-36 
		% ---
			\label{1NEPHI-13-36}

			And in them shall be written my gospel, saith the Lamb,
			and my rock and my salvation.

		% ---
		\paragraph{1 Nephi 13:37} % *1NEPHI-13-37 
		% ---
			\label{1NEPHI-13-37}

			And blessed are they which shall seek to bring forth my Zion at that day,
			for they shall have the gift and the power of the Holy Ghost.
			And if they endure unto the end,
			they shall be lifted up at the last day
			and shall be saved in the everlasting kingdom of the Lamb.
			Yea, whoso shall publish peace
			---that shall publish tidings of great joy ---
			how beautiful upon the mountains shall they be!

		% ---
		\paragraph{1 Nephi 13:38} % *1NEPHI-13-38 
		% ---
			\label{1NEPHI-13-38}

			And it came to pass that I beheld the remnant of the seed of my brethren
			and also the book of the Lamb of God
			which had proceeded forth from the mouth of the Jew.
			And I beheld that it came forth from the Gentiles
			unto the remnant of the seed of my brethren.

		% ---
		\paragraph{1 Nephi 13:39} % *1NEPHI-13-39 
		% ---
			\label{1NEPHI-13-39}

			And after it had come forth unto them, I beheld other books
			which came forth by the power of the Lamb from the Gentiles unto them,
			unto the convincing of the Gentiles and the remnant of the seed of my brethren
			---and also to the Jews, which were scattered upon all the face of the earth---
			that the records of the prophets and of the twelve apostles of the Lamb are true.

		% ---
		\paragraph{1 Nephi 13:40} % *1NEPHI-13-40 
		% ---
			\label{1NEPHI-13-40}

			And the angel spake unto me, saying:
			These last records which thou hast seen among the Gentiles
			shall establish the truth of the first,
			which is of the twelve apostles of the Lamb,
			and shall make known the plain and precious things
			which have been taken away from them
			and shall make known to all kindreds, tongues, and people
			that the Lamb of God is the Eternal Father and the Savior of the world
			and that all men must come unto him or they cannot be saved.

		% ---
		\paragraph{1 Nephi 13:41} % *1NEPHI-13-41 
		% ---
			\label{1NEPHI-13-41}

			And they must come according to the words
			which shall be established by the mouth of the Lamb.
			And the words of the Lamb shall be made known in the records of thy seed
			as well as in the records of the twelve apostles of the Lamb.
			Wherefore they both shall be established in one,
			for there is one God and one Shepherd over all the earth.

		% ---
		\paragraph{1 Nephi 13:42} % *1NEPHI-13-42 
		% ---
			\label{1NEPHI-13-42}

			And the time cometh that he shall manifest himself unto all nations,
			both unto the Jews and also unto the Gentiles.
			And after that he hath manifested himself unto the Jews and also unto the Gentiles,
			then he shall manifest himself unto the Gentiles and also unto the Jews.
			And the last shall be first and the first shall be last.

	% -----
	\subsubsection{1 Nephi 14} % *1NEPHI-14 
	% -----
		\label{1NEPHI-14}
		
		% ---
		\paragraph{1 Nephi 14:1} % *1NEPHI-14-1 
		% ---
			\label{1NEPHI-14-1}

			And it shall come to pass that
			if the Gentiles shall hearken unto the Lamb of God in that day,
			that he shall manifest himself unto them in word and also in power, in very deed,
			unto the taking away of their stumbling blocks,
			if it so be that they harden not their hearts against the Lamb.

		% ---
		\paragraph{1 Nephi 14:2} % *1NEPHI-14-2 
		% ---
			\label{1NEPHI-14-2}

			And if it so be that they harden not their hearts against the Lamb of God,
			they shall be numbered among the seed of thy father;
			yea, they shall be numbered among the house of Israel.
			And they shall be a blessed people upon the promised land forever.
			They shall be no more brought down into captivity,
			and the house of Israel shall no more be confounded.

		% ---
		\paragraph{1 Nephi 14:3} % *1NEPHI-14-3 
		% ---
			\label{1NEPHI-14-3}

			And that great pit which hath been digged for them
			by that great and abominable church,
			which was founded by the devil and his children
			that he might lead away the souls of men down to hell---
			yea, that great pit which hath been digged for the destruction of men
			shall be filled by those who digged it,
			unto their utter destruction, saith the Lamb of God---
			not the destruction of the soul,
			save it be the casting of it into that hell%
				\footnote{%
					% \label{}%
					See \hyperref[1NEPHI-15-35]{1 Nephi 15:35}, \hyperref[DC63-4]{DC 63:4}, \hyperref[DC64-17]{DC 64:17}, 
					}
			which hath no end.

		% ---
		\paragraph{1 Nephi 14:4} % *1NEPHI-14-4 
		% ---
			\label{1NEPHI-14-4}

			For behold, this is according to the captivity of the devil,
			and also according to the justice of God,
			upon all those who will work wickedness and abomination before him.

		% ---
		\paragraph{1 Nephi 14:5} % *1NEPHI-14-5 
		% ---
			\label{1NEPHI-14-5}

			And it came to pass that the angel spake unto me Nephi, saying:
			Thou hast beheld that if the Gentiles repent, it shall be well with them.
			And thou also knowest concerning the covenants of the Lord
			unto the house of Israel.
			And thou also hast heard that whoso repenteth not must perish.

		% ---
		\paragraph{1 Nephi 14:6} % *1NEPHI-14-6 
		% ---
			\label{1NEPHI-14-6}

			Therefore woe be unto the Gentiles
			if it so be that they harden their hearts against the Lamb of God.

		% ---
		\paragraph{1 Nephi 14:7} % *1NEPHI-14-7 
		% ---
			\label{1NEPHI-14-7}

			For the time cometh, saith the Lamb of God,
			that I will work a great and a marvelous work among the children of men,
			a work which shall be everlasting,
			either on the one hand or on the other,
			either to the convincing of them unto peace and life eternal
			or unto the deliverance of them to the hardness of their hearts
			and the blindness of their minds,
			unto their being brought down into captivity,
			and also unto destruction both temporally and spiritually,
			according to the captivity of the devil of which I have spoken.

		% ---
		\paragraph{1 Nephi 14:8} % *1NEPHI-14-8 
		% ---
			\label{1NEPHI-14-8}

			And it came to pass that
			when the angel had spoken these words,
			he saith unto me:
			Remember thou the covenants of the Father unto the house of Israel?
			I saith unto him: Yea.

		% ---
		\paragraph{1 Nephi 14:9} % *1NEPHI-14-9 
		% ---
			\label{1NEPHI-14-9}

			And it came to pass that he saith unto me:
			Look and behold that great and abominable church,
			which is the mother of abominations, whose founder is the devil.

		% ---
		\paragraph{1 Nephi 14:10} % *1NEPHI-14-10 
		% ---
			\label{1NEPHI-14-10}

			And he saith unto me:
			Behold, there is save it be two churches;
			the one is the church of the Lamb of God
			and the other is the church of the devil.
			Wherefore whoso belongeth not to the church of the Lamb of God
			belongeth to that great church which is the mother of abominations,
			and she is the whore of all the earth.

		% ---
		\paragraph{1 Nephi 14:11} % *1NEPHI-14-11 
		% ---
			\label{1NEPHI-14-11}

			And it came to pass that I looked and beheld the whore of all the earth.
			And she sat upon many waters,
			and she had dominion over all the earth
			among all nations, kindreds, tongues, and people.

		% ---
		\paragraph{1 Nephi 14:12} % *1NEPHI-14-12 
		% ---
			\label{1NEPHI-14-12}

			And it came to pass that I beheld the church of the Lamb of God;
			and its numbers were few
			because of the wickedness and abominations of the whore
			which sat upon many waters.
			Nevertheless I beheld that the church of the Lamb, which were the saints of God,
			were also upon all the face of the earth;
			and their dominions upon the face of the earth were small
			because of the wickedness of the great whore which I saw.

		% ---
		\paragraph{1 Nephi 14:13} % *1NEPHI-14-13 
		% ---
			\label{1NEPHI-14-13}

			And it came to pass that I beheld that the great mother of abominations
			did gather together in multitudes upon the face of all the earth
			among all the nations of the Gentiles
			to fight against the Lamb of God.

		% ---
		\paragraph{1 Nephi 14:14} % *1NEPHI-14-14 
		% ---
			\label{1NEPHI-14-14}

			And it came to pass that I Nephi beheld the power of the Lamb of God,
			that it descended upon the saints of the church of the Lamb
			and upon the covenant people of the Lord,
			which were scattered upon all the face of the earth.
			And they were armed with righteousness
			and with the power of God in great glory.

		% ---
		\paragraph{1 Nephi 14:15} % *1NEPHI-14-15 
		% ---
			\label{1NEPHI-14-15}

			And it came to pass that I beheld that
			the wrath of God was poured out upon that great and abominable church,
			insomuch that there were wars and rumors of wars
			among all the nations and kindreds of the earth.

		% ---
		\paragraph{1 Nephi 14:16} % *1NEPHI-14-16 
		% ---
			\label{1NEPHI-14-16}

			And as there began to be wars and rumors of wars
			among all the nations which belonged to the mother of abominations,
			the angel spake unto me, saying:
			Behold, the wrath of God is upon the mother of harlots.
			And behold, thou seest all these things.

		% ---
		\paragraph{1 Nephi 14:17} % *1NEPHI-14-17 
		% ---
			\label{1NEPHI-14-17}

			And when the day cometh
			that the wrath of God is poured out upon the mother of harlots
			---which is the great and abominable church of all the earth,
			whose founder is the devil---
			then at that day the work of the Father shall commence
			in preparing the way for the fulfilling of his covenants
			which he hath made to his people which are of the house of Israel.

		% ---
		\paragraph{1 Nephi 14:18} % *1NEPHI-14-18 
		% ---
			\label{1NEPHI-14-18}

			And it came to pass that the angel spake unto me, saying: Look!

		% ---
		\paragraph{1 Nephi 14:19} % *1NEPHI-14-19 
		% ---
			\label{1NEPHI-14-19}

			And I looked and beheld a man, and he was dressed in a white robe.

		% ---
		\paragraph{1 Nephi 14:20} % *1NEPHI-14-20 
		% ---
			\label{1NEPHI-14-20}

			And the angel said unto me:
			Behold one of the twelve apostles of the Lamb.

		% ---
		\paragraph{1 Nephi 14:21} % *1NEPHI-14-21 
		% ---
			\label{1NEPHI-14-21}

			Behold, he shall see and write the remainder of these things,
			yea, and also many things which have been.

		% ---
		\paragraph{1 Nephi 14:22} % *1NEPHI-14-22 
		% ---
			\label{1NEPHI-14-22}

			And he shall also write concerning the end of the world.

		% ---
		\paragraph{1 Nephi 14:23} % *1NEPHI-14-23 
		% ---
			\label{1NEPHI-14-23}

			Wherefore the things which he shall write are \hyperref[TRUTH-PHRASES]{just and true}.
			And behold, they are written in the book
			which thou beheld proceeding out of the mouth of the Jew.
			And at the time they proceeded out of the mouth of the Jew,
			or at the time the book proceeded out of the mouth of the Jew,
			the things which were written were plain and pure and most precious
			and easy to the understanding of all men.

		% ---
		\paragraph{1 Nephi 14:24} % *1NEPHI-14-24 
		% ---
			\label{1NEPHI-14-24}

			And behold, the things which this apostle of the Lamb shall write
			are many things which thou hast seen.
			And behold, the remainder shalt thou see.

		% ---
		\paragraph{1 Nephi 14:25} % *1NEPHI-14-25 
		% ---
			\label{1NEPHI-14-25}

			But the things which thou shalt see hereafter
			thou shalt not write,
			for the Lord God hath ordained the apostle of the Lamb of God
			that he should write them.

		% ---
		\paragraph{1 Nephi 14:26} % *1NEPHI-14-26 
		% ---
			\label{1NEPHI-14-26}

			And also others which have been,
			to them hath he shown all things,
			and they have written them.
			And they are sealed up to come forth in their purity,
			according to the truth which is in the Lamb,
			in the own due time of the Lord,
			unto the house of Israel.

		% ---
		\paragraph{1 Nephi 14:27} % *1NEPHI-14-27 
		% ---
			\label{1NEPHI-14-27}

			And I Nephi heard and bare record
			that the name of the apostle of the Lamb was John,
			according to the word of the angel.

		% ---
		\paragraph{1 Nephi 14:28} % *1NEPHI-14-28 
		% ---
			\label{1NEPHI-14-28}

			And behold, I Nephi am forbidden
			that I should write the remainder of the things which I saw.
			Wherefore the things which I have written sufficeth me,
			and I have not written but a small part of the things which I saw.

		% ---
		\paragraph{1 Nephi 14:29} % *1NEPHI-14-29 
		% ---
			\label{1NEPHI-14-29}

			And I bear record that I saw the things which my father saw,
			and the angel of the Lord did make them known unto me.

		% ---
		\paragraph{1 Nephi 14:30} % *1NEPHI-14-30 
		% ---
			\label{1NEPHI-14-30}

			And now I make an end of speaking concerning the things
			which I saw while I was carried away in the spirit.
			And if all the things which I saw are not written,
			the things which I have written are true.
			And thus it is.
			Amen.

	% -----
	\subsubsection{1 Nephi 15} % *1NEPHI-15 
	% -----
		\label{1NEPHI-15}
		
		% ---
		\paragraph{1 Nephi 15:1} % *1NEPHI-15-1 
		% ---
			\label{1NEPHI-15-1}

			And it came to pass that
			after I Nephi had been carried away in the spirit and seen all these things,
			I returned to the tent of my father.

		% ---
		\paragraph{1 Nephi 15:2} % *1NEPHI-15-2 
		% ---
			\label{1NEPHI-15-2}

			And it came to pass that I beheld my brethren,
			and they were disputing one with another
			concerning the things which my father had spoken unto them.

		% ---
		\paragraph{1 Nephi 15:3} % *1NEPHI-15-3 
		% ---
			\label{1NEPHI-15-3}

			For he truly spake many great things unto them
			which was hard to be understood
			save a man should inquire of the Lord.
			And they being hard in their hearts,
			therefore they did not look unto the Lord as they had ought.

		% ---
		\paragraph{1 Nephi 15:4} % *1NEPHI-15-4 
		% ---
			\label{1NEPHI-15-4}

			And now I Nephi was grieved because of the hardness of their hearts,
			and also because of the things which I had seen,
			and knew they must unavoidably come to pass
			because of the great wickedness of the children of men.

		% ---
		\paragraph{1 Nephi 15:5} % *1NEPHI-15-5 
		% ---
			\label{1NEPHI-15-5}

			And it came to pass that I was overcome because of my afflictions,
			for I considered that mine afflictions was great above all
			because of the destruction of my people,
			for I had beheld their fall.

		% ---
		\paragraph{1 Nephi 15:6} % *1NEPHI-15-6 
		% ---
			\label{1NEPHI-15-6}

			And it came to pass that after I had received strength,
			I spake unto my brethren,
			desiring to know of them the cause of their disputations.

		% ---
		\paragraph{1 Nephi 15:7} % *1NEPHI-15-7 
		% ---
			\label{1NEPHI-15-7}

			And they said:
			Behold, we cannot understand the words which our father hath spoken
			concerning the natural branches of the olive tree
			and also concerning the Gentiles.

		% ---
		\paragraph{1 Nephi 15:8} % *1NEPHI-15-8 
		% ---
			\label{1NEPHI-15-8}

			And I said unto them:
			Have ye inquired of the Lord?

		% ---
		\paragraph{1 Nephi 15:9} % *1NEPHI-15-9 
		% ---
			\label{1NEPHI-15-9}

			And they said unto me:
			We have not,
			for the Lord maketh no such thing known unto us.

		% ---
		\paragraph{1 Nephi 15:10} % *1NEPHI-15-10 
		% ---
			\label{1NEPHI-15-10}

			Behold, I said unto them:
			How is it that ye do not keep the commandments of the Lord?
			How is it that ye will perish because of the hardness of your hearts?

		% ---
		\paragraph{1 Nephi 15:11} % *1NEPHI-15-11 
		% ---
			\label{1NEPHI-15-11}

			Do ye not remember the thing which the Lord hath said?
			---if ye will not harden your hearts and ask me in faith,
			believing that ye shall receive,
			with diligence in keeping my commandments,
			surely these things shall be made known unto you.

		% ---
		\paragraph{1 Nephi 15:12} % *1NEPHI-15-12 
		% ---
			\label{1NEPHI-15-12}

			Behold, I say unto you that the house of Israel was compared unto an olive tree
			by the Spirit of the Lord which was in our father.
			And behold, are we not broken off from the house of Israel?
			And are we not a branch of the house of Israel?

		% ---
		\paragraph{1 Nephi 15:13} % *1NEPHI-15-13 
		% ---
			\label{1NEPHI-15-13}

			And now the thing which our father meaneth
			concerning the grafting in of the natural branches
			through the fullness of the Gentiles
			is that in the latter days, when our seed shall have dwindled in unbelief,
			yea, for the space of many years and many generations
			after that the Messiah hath manifested himself in body unto the children of men,
			then shall the fullness of the gospel of the Messiah come unto the Gentiles,
			and from the Gentiles unto the remnant of our seed.

		% ---
		\paragraph{1 Nephi 15:14} % *1NEPHI-15-14 
		% ---
			\label{1NEPHI-15-14}

			And at that day shall the remnant of our seed know
			that they are of the house of Israel
			and that they are the covenant people of the Lord.
			And then shall they know and come to the knowledge of their forefathers
			and also to the knowledge of the gospel of their Redeemer,
			which was ministered unto their fathers by him.
			Wherefore they shall come to the knowledge of their Redeemer
			and the very points of his doctrine,
			that they may know how to come unto him and be saved.%
				\footnote{%
					\label{1NEPHI-15-14-NOTE-how}%
					Notice the relationship here between having knowledge of points of doctrine, and knowing how to do something. The purpose of the points of doctrine is practical application, or action: the doctrine tells us how to \textit{do} something. The purpose of doctrine is to guide us in how we live so that we can accomplish certain things, things that we wouldn't be able to accomplish if we didn't have the knowledge of the doctrine. See \hyperref[REPENTANCE-learning]{this study}.
					}

		% ---
		\paragraph{1 Nephi 15:15} % *1NEPHI-15-15 
		% ---
			\label{1NEPHI-15-15}

			And then at that day will they not rejoice
			and give praise unto their everlasting God, their rock and their salvation?
			Yea, at that day will they not receive strength and nourishment from the true vine?
			Yea, will they not come unto the true fold of God?

		% ---
		\paragraph{1 Nephi 15:16} % *1NEPHI-15-16 
		% ---
			\label{1NEPHI-15-16}

			Behold, I say unto you:
			Yea, they shall be numbered again among the house of Israel;
			they shall be grafted in,
			being a natural branch of the olive tree,
			into the true olive tree.

		% ---
		\paragraph{1 Nephi 15:17} % *1NEPHI-15-17 
		% ---
			\label{1NEPHI-15-17}

			And this is what our father meaneth.
			And he meaneth that it will not come to pass
			until after that they are scattered by the Gentiles.
			And he meaneth that it shall come by way of the Gentiles,
			that the Lord may shew his power unto the Gentiles,
			for the very cause that he shall be rejected of the Jews or of the house of Israel.

		% ---
		\paragraph{1 Nephi 15:18} % *1NEPHI-15-18 
		% ---
			\label{1NEPHI-15-18}

			Wherefore our father hath not spoken of our seed alone
			but also of all the house of Israel,
			pointing to the covenant which should be fulfilled in the latter days,
			which covenant the Lord made to our father Abraham, saying:
			In thy seed shall all the kindreds of the earth be blessed.

		% ---
		\paragraph{1 Nephi 15:19} % *1NEPHI-15-19 
		% ---
			\label{1NEPHI-15-19}

			And it came to pass that I Nephi spake much unto them concerning these things;
			yea, I spake unto them concerning the restoration of the Jews in the latter days.

		% ---
		\paragraph{1 Nephi 15:20} % *1NEPHI-15-20 
		% ---
			\label{1NEPHI-15-20}

			And I did rehearse unto them the words of Isaiah,
			which spake concerning the restoration of the Jews or of the house of Israel.
			And after that they were restored, they should no more be confounded,
			neither should they be scattered again.
			And it came to pass that I did speak so many words unto my brethren
			that they were pacified and did humble themselves before the Lord.

		% ---
		\paragraph{1 Nephi 15:21} % *1NEPHI-15-21 
		% ---
			\label{1NEPHI-15-21}

			And it came to pass that they did speak unto me again, saying:
			What meaneth the things which our father saw in a dream?
			What meaneth the tree which he saw?

		% ---
		\paragraph{1 Nephi 15:22} % *1NEPHI-15-22 
		% ---
			\label{1NEPHI-15-22}

			And I said unto them:
			It was a representation of the tree of life.

		% ---
		\paragraph{1 Nephi 15:23} % *1NEPHI-15-23 
		% ---
			\label{1NEPHI-15-23}

			And they said unto me:
			What meaneth the rod of iron which our father saw that led to the tree?

		% ---
		\paragraph{1 Nephi 15:24} % *1NEPHI-15-24 
		% ---
			\label{1NEPHI-15-24}

			And I said unto them that it was the word of God,
			and that whoso would hearken unto the word of God and would hold fast unto it,
			they would never perish,
			neither could the temptations and the fiery darts of the adversary
			overpower them unto blindness,
			to lead them away to destruction.

		% ---
		\paragraph{1 Nephi 15:25} % *1NEPHI-15-25 
		% ---
			\label{1NEPHI-15-25}

			Wherefore I Nephi did exhort them to give heed unto the word of the Lord.
			Yea, I did exhort them with all the energies of my soul
			and with all the faculty which I possessed
			that they would give heed to the word of God
			and remember to keep his commandments always in all things.

		% ---
		\paragraph{1 Nephi 15:26} % *1NEPHI-15-26 
		% ---
			\label{1NEPHI-15-26}

			And they said unto me:
			What meaneth the river of water which our father saw?

		% ---
		\paragraph{1 Nephi 15:27} % *1NEPHI-15-27 
		% ---
			\label{1NEPHI-15-27}

			And I said unto them that the water which my father saw was filthiness.
			And so much was his mind swallowed up in other things
			that he beheld not the filthiness of the water.

		% ---
		\paragraph{1 Nephi 15:28} % *1NEPHI-15-28 
		% ---
			\label{1NEPHI-15-28}

			And I said unto them that it was an awful gulf
			which separateth the wicked from the tree of life and also from the saints of God.

		% ---
		\paragraph{1 Nephi 15:29} % *1NEPHI-15-29 
		% ---
			\label{1NEPHI-15-29}

			And I said unto them that it was a representation of that awful hell
			which the angel said unto me was prepared for the wicked.

		% ---
		\paragraph{1 Nephi 15:30} % *1NEPHI-15-30 
		% ---
			\label{1NEPHI-15-30}

			And I said unto them that our father also saw
			that the justice of God did also divide the wicked from the righteous,
			and the brightness thereof was like unto the brightness of a flaming fire
			which ascendeth up unto God forever and ever and hath no end.

		% ---
		\paragraph{1 Nephi 15:31} % *1NEPHI-15-31 
		% ---
			\label{1NEPHI-15-31}

			And they said unto me:
			Doth this thing mean the torment of the body in the days of probation,
			or doth it mean the final state of the soul after the death of the temporal body,
			or doth it speak of the things which are temporal?

		% ---
		\paragraph{1 Nephi 15:32} % *1NEPHI-15-32 
		% ---
			\label{1NEPHI-15-32}

			And it came to pass that I said unto them
			that it was a representation of things both temporal and spiritual.
			For the day should come that they must be judged of their works,
			yea, even the works which were done by the temporal body
			in their days of probation.

		% ---
		\paragraph{1 Nephi 15:33} % *1NEPHI-15-33 
		% ---
			\label{1NEPHI-15-33}

			Wherefore if they should die in their wickedness,
			they must be cast off also as to the things which are spiritual,
			which are pertaining unto righteousness.
			Wherefore they must be brought to stand before God
			to be judged of their works.
			And if their works have been filthiness, they must needs be filthy.
			And if they be filthy, it must needs be
			that they cannot dwell in the kingdom of God;
			if so, the kingdom of God must be filthy also.

		% ---
		\paragraph{1 Nephi 15:34} % *1NEPHI-15-34 
		% ---
			\label{1NEPHI-15-34}

			But behold, I say unto you that the kingdom of God is not filthy,
			that there cannot any unclean thing enter into the kingdom of God.
			Wherefore there must needs be a place of filthiness
			prepared for that which is filthy.

		% ---
		\paragraph{1 Nephi 15:35} % *1NEPHI-15-35 
		% ---
			\label{1NEPHI-15-35}

			And there is a place prepared,
			yea, even that awful hell of which I have spoken,
			and the devil is the proprietor%
				\footnote{%
					\label{1NEPHI-15-35-NOTE-proprietor}%
					The current edition of the BOM has ``preparator'' here. The 1828 definition of ``proprietor'' is: ``An owner; the person who has the legal right or exclusive title to any thing whether in possession or not.'' This is actually an important difference, I think, because ``proprietor'' carries with it the idea of a legal right, and so you can see this place as the kingdom of a devil, governed by certain laws.
					}
			of it.
			Wherefore the final state of the soul of man
			is to dwell in the kingdom of God or to be cast out
			because of that justice of which I have spoken.

		% ---
		\paragraph{1 Nephi 15:36} % *1NEPHI-15-36 
		% ---
			\label{1NEPHI-15-36}

			Wherefore the wicked are separated from the righteous
			and also from that tree of life,
			whose fruit is most precious and most desirable of all other fruits;
			yea, and it is the greatest%
				\footnote{%
					% \label{}%
					Compare \hyperref[DC14-7]{DC 14:7}. The fruit is eternal life.
					}
			of all the gifts of God.
			And thus I spake unto my brethren.
			Amen.

	% -----
	\subsubsection{1 Nephi 16} % *1NEPHI-16 
	% -----
		\label{1NEPHI-16}
		
		% ---
		\paragraph{1 Nephi 16:1} % *1NEPHI-16-1 
		% ---
			\label{1NEPHI-16-1}

			And now it came to pass that
			after I Nephi had made an end of speaking to my brethren,
			behold, they said unto me:
			Thou hast declared unto us hard things,
			more than that which we are able to bear.

		% ---
		\paragraph{1 Nephi 16:2} % *1NEPHI-16-2 
		% ---
			\label{1NEPHI-16-2}

			And it came to pass that I said unto them
			that I knew that I had spoken hard things against the wicked
			according to the truth,
			and the righteous have I justified
			and testified that they should be lifted up at the last day.
			Wherefore the guilty taketh the truth to be hard,
			for it cutteth them to the very center.

		% ---
		\paragraph{1 Nephi 16:3} % *1NEPHI-16-3 
		% ---
			\label{1NEPHI-16-3}

			And now my brethren, if ye were righteous
			and were willing to hearken to the truth and give heed unto it,
			that ye might walk uprightly before God,
			then ye would not murmur because of the truth and say:
			Thou speakest hard things against us.

		% ---
		\paragraph{1 Nephi 16:4} % *1NEPHI-16-4 
		% ---
			\label{1NEPHI-16-4}

			And it came to pass that I Nephi did exhort my brethren with all diligence
			to keep the commandments of the Lord.

		% ---
		\paragraph{1 Nephi 16:5} % *1NEPHI-16-5 
		% ---
			\label{1NEPHI-16-5}

			And it came to pass that they did humble themselves before the Lord,
			insomuch that I had joy and great hopes of them,
			that they would walk in the paths of righteousness.

		% ---
		\paragraph{1 Nephi 16:6} % *1NEPHI-16-6 
		% ---
			\label{1NEPHI-16-6}

			Now all these things were said and done
			as my father dwelt in a tent in the valley which he called Lemuel.

		% ---
		\paragraph{1 Nephi 16:7} % *1NEPHI-16-7 
		% ---
			\label{1NEPHI-16-7}

			And it came to pass that I Nephi took one of the daughters of Ishmael to wife,
			and also my brethren took of the daughters of Ishmael to wife,
			and also Zoram took the eldest daughter of Ishmael to wife.%
				\footnote{%
					\label{1NEPHI-16-7-NOTE-daughters}%
					How were all these daughters still single? If they were all old enough to be married, how had none of them already gotten married?
					}

		% ---
		\paragraph{1 Nephi 16:8} % *1NEPHI-16-8 
		% ---
			\label{1NEPHI-16-8}

			And thus my father had fulfilled all the commandments of the Lord
			which had been given unto him.
			And also I Nephi had been blessed of the Lord exceedingly.

		% ---
		\paragraph{1 Nephi 16:9} % *1NEPHI-16-9 
		% ---
			\label{1NEPHI-16-9}

			And it came to pass that the voice of the Lord spake unto my father by night
			and commanded him that on the morrow
			he should take his journey into the wilderness.

		% ---
		\paragraph{1 Nephi 16:10} % *1NEPHI-16-10 
		% ---
			\label{1NEPHI-16-10}

			And it came to pass that as my father arose in the morning
			and went forth to the tent door,
			and to his great astonishment he beheld upon the ground
			a round ball of curious workmanship,
			and it was of fine brass.
			And within the ball was two spindles,%
				\footnote{%
					\label{1NEPHI-16-10-NOTE-spindles}%
					From the 1828 dictionary, for ``spindle'':
					\begin{compactenum}
						\item The pin used in spinning wheels for twisting the thread, and on which the thread when twisted, is wound.
						\item A slender pointed rod or pin on which any thing turn.
						\item The fusee of a watch.
						\item A long slender stalk.
						\item The lower end of a capstan, shod with iron; the pivot.
					\end{compactenum}
					Apart from definition \#4, all the definitions have in common the idea of a circular or conical rod around which things turn. Nephi says that the ball had two spindles and that one pointed the way to go. It's possible that one spindle rotated around the other and pointed the right way, or it's possible that the one spindle pointed the right way, maybe with an arrow painted onto it or something, and the other spindle did something else. \hyperref[ALMA-37-40]{Alma 37:40} says that the ``spindles should point the way they should go,'' so it seems like the spindles work together and don't have separate functions, but it's impossible to say for sure. See also \hyperref[1NEPHI-16-28]{verse 28}, where they are referred to as ``pointers.''
					}
			and the one pointed the way whither we should go into the wilderness.

		% ---
		\paragraph{1 Nephi 16:11} % *1NEPHI-16-11 
		% ---
			\label{1NEPHI-16-11}

			And it came to pass that we did gather together
			whatsoever things we should carry into the wilderness
			and all the remainder of our provisions which the Lord had given unto us.
			And we did take seed of every kind that we might carry into the wilderness.

		% ---
		\paragraph{1 Nephi 16:12} % *1NEPHI-16-12 
		% ---
			\label{1NEPHI-16-12}

			And it came to pass that we did take our tents
			and departed into the wilderness across the river Laman.

		% ---
		\paragraph{1 Nephi 16:13} % *1NEPHI-16-13 
		% ---
			\label{1NEPHI-16-13}

			And it came to pass that we traveled for the space of four days
			nearly a south-southeast direction.
			And we did pitch our tents again,
			and we did call the name of the place Shazer.%
				\footnote{%
					\label{1NEPHI-16-13-shazer}%
					From the BOM Onomasticon, copied 24 March 2021: ``The most likely suggestion for \textsc{SHAZER} is from Sidney Sperry, who took it from the \textsc{HEBREW} root \textit{\v{s}azar}, meaning `to twist, to intertwine,' but which only appears in the hophal participle in the \textsc{HEBREW} Bible....The understanding of \textsc{SHAZER} in \textsc{HEBREW} as `twisted' could be a reference to the acacia trees at this oasis where \textsc{LEHI}'s group camped and rested. According to James K. Hoffmeier in his book, \textit{Ancient Israel in Sinai: The Evidence for the Authenticity of the Wilderness Tradition}: `Most acacia are short and tend to be twisted.' Significantly, acacia trees tend to grow in wadis where there is moisture. It is cognate with Arabic \textit{\v{s}azara} `look askew at, suspiciously; twisted wrong way' as for a cord (from left). The vowels of \textsc{SHAZER} would be the vowels of the \textit{qal} participle form, and could mean something like, `twister.''' (Note that a \textit{wadi} is ``a valley, ravine, or channel that is dry except in the rainy season.'')
					}

		% ---
		\paragraph{1 Nephi 16:14} % *1NEPHI-16-14 
		% ---
			\label{1NEPHI-16-14}

			And it came to pass that we did take our bows and our arrows
			and go forth into the wilderness to slay food for our families.
			And after that we had slain food for our families,
			we did return again to our families in the wilderness to the place of Shazer.
			And we did go forth again in the wilderness, following the same direction,
			keeping in the most fertile parts of the wilderness
			which was in the borders near the Red Sea.

		% ---
		\paragraph{1 Nephi 16:15} % *1NEPHI-16-15 
		% ---
			\label{1NEPHI-16-15}

			And it came to pass that we did travel for the space of many days,
			slaying food by the way
			with our bows and our arrows and our stones and our slings.

		% ---
		\paragraph{1 Nephi 16:16} % *1NEPHI-16-16 
		% ---
			\label{1NEPHI-16-16}

			And we did follow the directions of the ball,
			which led us in the more fertile parts of the wilderness.

		% ---
		\paragraph{1 Nephi 16:17} % *1NEPHI-16-17 
		% ---
			\label{1NEPHI-16-17}

			And after that we had traveled for the space of many days,
			we did pitch our tents for the space of a time,
			that we might again rest ourselves and obtain food for our families.

		% ---
		\paragraph{1 Nephi 16:18} % *1NEPHI-16-18 
		% ---
			\label{1NEPHI-16-18}

			And it came to pass that as I Nephi went forth to slay food,
			behold, I did break my bow, which was made of fine steel.
			And after that I did break my bow,
			behold, my brethren were angry with me because of the loss of my bow,
			for we did obtain no food.

		% ---
		\paragraph{1 Nephi 16:19} % *1NEPHI-16-19 
		% ---
			\label{1NEPHI-16-19}

			And it came to pass that we did return without food to our families;
			and being much fatigued because of their journeying,
			they did suffer much for the want of food.

		% ---
		\paragraph{1 Nephi 16:20} % *1NEPHI-16-20 
		% ---
			\label{1NEPHI-16-20}

			And it came to pass that Laman and Lemuel and the sons of Ishmael
			did begin to murmur exceedingly
			because of their sufferings and afflictions in the wilderness.
			And also my father began to murmur against the Lord his God;
			yea, and they were all exceeding sorrowful,
			even that they did murmur against the Lord.

		% ---
		\paragraph{1 Nephi 16:21} % *1NEPHI-16-21 
		% ---
			\label{1NEPHI-16-21}

			Now it came to pass that
			I Nephi having been afflicted with my brethren because of the loss of my bow
			and their bows having lost their springs,
			it began to be exceeding difficult,
			yea, insomuch that we could obtain no food.

		% ---
		\paragraph{1 Nephi 16:22} % *1NEPHI-16-22 
		% ---
			\label{1NEPHI-16-22}

			And it came to pass that I Nephi did speak much unto my brethren
			because that they had hardened their hearts again,
			even unto complaining against the Lord their God.

		% ---
		\paragraph{1 Nephi 16:23} % *1NEPHI-16-23 
		% ---
			\label{1NEPHI-16-23}

			And it came to pass that I Nephi did make out of wood a bow
			and out of a straight stick an arrow;
			wherefore I did arm myself with a bow and an arrow,
			with a sling and with stones.
			And I said unto my father:
			Whither shall I go to obtain food?

		% ---
		\paragraph{1 Nephi 16:24} % *1NEPHI-16-24 
		% ---
			\label{1NEPHI-16-24}

			And it came to pass that he did inquire of the Lord,
			for they had humbled themselves because of my words;
			for I did say many things unto them in the energy of my soul.

		% ---
		\paragraph{1 Nephi 16:25} % *1NEPHI-16-25 
		% ---
			\label{1NEPHI-16-25}

			And it came to pass that the voice of the Lord came unto my father,
			and he was truly chastened because of his murmurings against the Lord,
			insomuch that he was brought down into the depths of sorrow.

		% ---
		\paragraph{1 Nephi 16:26} % *1NEPHI-16-26 
		% ---
			\label{1NEPHI-16-26}

			And it came to pass that the voice of the Lord said unto him:
			Look upon the ball and behold the things which are written.

		% ---
		\paragraph{1 Nephi 16:27} % *1NEPHI-16-27 
		% ---
			\label{1NEPHI-16-27}

			And it came to pass that when my father beheld the things
			which were written upon the ball,
			he did fear and tremble exceedingly,
			and also my brethren and the sons of Ishmael and our wives.

		% ---
		\paragraph{1 Nephi 16:28} % *1NEPHI-16-28 
		% ---
			\label{1NEPHI-16-28}

			And it came to pass that I Nephi beheld that the pointers which were in the ball
			that they did work according to the faith and diligence and heed
			which we did give unto them.

		% ---
		\paragraph{1 Nephi 16:29} % *1NEPHI-16-29 
		% ---
			\label{1NEPHI-16-29}

			And there was also written upon them a new writing which was plain to be read,
			which did give us understanding concerning the ways of the Lord.
			And it was written and changed from time to time
			according to the faith and diligence which we gave unto it.
			And thus we see that by small means the Lord can bring about great things.

		% ---
		\paragraph{1 Nephi 16:30} % *1NEPHI-16-30 
		% ---
			\label{1NEPHI-16-30}

			And it came to pass that I Nephi did go forth up into the top of the mountain
			according to the directions which was given upon the ball.

		% ---
		\paragraph{1 Nephi 16:31} % *1NEPHI-16-31 
		% ---
			\label{1NEPHI-16-31}

			And it came to pass that I did slay wild beasts,
			insomuch that I did obtain food for our families.%
				\footnote{%
					% \label{}%
					Note that it doesn't say he used his bow and arrow or the sling.
					}

		% ---
		\paragraph{1 Nephi 16:32} % *1NEPHI-16-32 
		% ---
			\label{1NEPHI-16-32}

			And it came to pass that I did return to our tents,
			bearing the beasts which I had slain.
			And now when they beheld that I had obtained food,
			how great was their joy!
			And it came to pass that they did humble themselves before the Lord
			and did give thanks unto him.

		% ---
		\paragraph{1 Nephi 16:33} % *1NEPHI-16-33 
		% ---
			\label{1NEPHI-16-33}

			And it came to pass that we did again take our journey,
			traveling nearly the same course as in the beginning.
			And after that we had traveled for the space of many days,
			we did pitch our tents again,
			that we might tarry for the space of a time.

		% ---
		\paragraph{1 Nephi 16:34} % *1NEPHI-16-34 
		% ---
			\label{1NEPHI-16-34}

			And it came to pass that Ishmael died
			and was buried in the place which was called Nahom.

		% ---
		\paragraph{1 Nephi 16:35} % *1NEPHI-16-35 
		% ---
			\label{1NEPHI-16-35}

			And it came to pass that the daughters of Ishmael did mourn exceedingly
			because of the loss of their father
			and because of their afflictions in the wilderness.
			And they did murmur against my father
			because that he had brought them out of the land of Jerusalem, saying:
			Our father is dead.
			Yea, and we have wandered much in the wilderness,
			and we have suffered much afflictions, hunger, thirst, and fatigue.
			And after all these sufferings we must perish in the wilderness with hunger.

		% ---
		\paragraph{1 Nephi 16:36} % *1NEPHI-16-36 
		% ---
			\label{1NEPHI-16-36}

			And thus they did murmur against my father and also against me.
			And they were desirous to return again to Jerusalem.

		% ---
		\paragraph{1 Nephi 16:37} % *1NEPHI-16-37 
		% ---
			\label{1NEPHI-16-37}

			And Laman saith unto Lemuel and also unto the sons of Ishmael:
			Behold, let us slay our father and also our brother Nephi,
			who hath taken it upon him to be our ruler and our teacher,
			who are his elder brethren.

		% ---
		\paragraph{1 Nephi 16:38} % *1NEPHI-16-38 
		% ---
			\label{1NEPHI-16-38}

			Now he saith that the Lord hath talked with him,
			and also that angels hath ministered unto him.
			But behold, we know that he lieth unto us.
			And he telleth us these things,
			and he worketh many things by his cunning arts
			that he may deceive our eyes,
			thinking perhaps that he may lead us away into some strange wilderness.
			And after that he hath led us away,
			he hath thought to make himself a king and a ruler over us,
			that he may do with us according to his will and pleasure.
			And after this manner did my brother Laman stir up their hearts to anger.%
				\footnote{%
					\label{1NEPHI-16-38-NOTE-lies}%
					Laman makes three important claims in this verse, and all three are things that we do to ourselves and with our own experiences at different times:
					\begin{compactdesc}
						\item[\#1: Openly deny past experiences]
						\item[\#2: Explain past experiences away]
						\item[\#3: Create a false narrative to reframe the present and future]
					\end{compactdesc}
					(Compare these attitudes with the remarks of Nephi's brothers at \hyperref[1NEPHI-17-17]{1 Nephi 17:17--22}.)
					}

		% ---
		\paragraph{1 Nephi 16:39} % *1NEPHI-16-39 
		% ---
			\label{1NEPHI-16-39}

			And it came to pass that the Lord was with us,
			yea, even the voice of the Lord came and did speak many words unto them
			and did chasten them exceedingly.
			And after that they were chastened by the voice of the Lord,
			they did turn away their anger and did repent of their sins,
			insomuch that the Lord did bless us again with food that we did not perish.

	% -----
	\subsubsection{1 Nephi 17} % *1NEPHI-17 
	% -----
		\label{1NEPHI-17}
		
		% ---
		\paragraph{1 Nephi 17:1} % *1NEPHI-17-1 
		% ---
			\label{1NEPHI-17-1}

			And it came to pass that we did again take our journey in the wilderness.
			And we did travel nearly eastward from that time forth.
			And we did travel and wade through much affliction in the wilderness,
			and our women bare children in the wilderness.

		% ---
		\paragraph{1 Nephi 17:2} % *1NEPHI-17-2 
		% ---
			\label{1NEPHI-17-2}

			And so great was the blessings of the Lord upon us
			that while we did live upon raw meat in the wilderness,
			our women did give plenty of suck for their children and were strong,
			yea, even like unto the men.
			And they began to bear their journeyings without murmuring.

		% ---
		\paragraph{1 Nephi 17:3} % *1NEPHI-17-3 
		% ---
			\label{1NEPHI-17-3}

			And thus we see that the commandments of God must be fulfilled.
			And if it so be that the children of men keep the commandments of God,
			he doth nourish them and strengthen them
			and provide ways and means whereby they can accomplish the thing
			which he hath commanded them.%
			\footnote{%
				\label{1NEPHI-17-3-NOTE-better}%
				This verse is a better version of the principle taught in \hyperref[1NEPHI-3-7]{1 Nephi 3:7}.
				} 
			Wherefore he did provide ways and means for us
			while we did sojourn in the wilderness.

		% ---
		\paragraph{1 Nephi 17:4} % *1NEPHI-17-4 
		% ---
			\label{1NEPHI-17-4}

			And we did sojourn for the space of many years,
			yea, even eight years in the wilderness.

		% ---
		\paragraph{1 Nephi 17:5} % *1NEPHI-17-5 
		% ---
			\label{1NEPHI-17-5}

			And we did come to the land which we called Bountiful
			because of its much fruit and also wild honey.
			And all these things were prepared of the Lord that we might not perish.
			And we beheld the sea, which we called Irreantum,
			which being interpreted is many waters.

		% ---
		\paragraph{1 Nephi 17:6} % *1NEPHI-17-6 
		% ---
			\label{1NEPHI-17-6}

			And it came to pass that we did pitch our tents by the seashore.
			And notwithstanding we had suffered many afflictions and much difficulty,
			yea, even so much that we cannot write them all,
			we was exceedingly rejoiced when we came to the seashore.
			And we called the place Bountiful because of its much fruit.

		% ---
		\paragraph{1 Nephi 17:7} % *1NEPHI-17-7 
		% ---
			\label{1NEPHI-17-7}

			And it came to pass that after I Nephi had been in the land Bountiful
			for the space of many days,
			the voice of the Lord came unto me, saying:
			Arise and get thee into the mountain.
			And it came to pass that I arose and went up into the mountain
			and cried unto the Lord.

		% ---
		\paragraph{1 Nephi 17:8} % *1NEPHI-17-8 
		% ---
			\label{1NEPHI-17-8}

			And it came to pass that the Lord spake unto me, saying:
			Thou shalt construct a ship after the manner which I shall shew thee
			that I may carry thy people across these waters.

		% ---
		\paragraph{1 Nephi 17:9} % *1NEPHI-17-9 
		% ---
			\label{1NEPHI-17-9}

			And I saith:
			Lord, whither shall I go that I may find ore to molten
			that I may make tools to construct the ship
			after the manner which thou hast shewn unto me?

		% ---
		\paragraph{1 Nephi 17:10} % *1NEPHI-17-10 
		% ---
			\label{1NEPHI-17-10}

			And it came to pass that the Lord told me
			whither I should go to find ore that I might make tools.

		% ---
		\paragraph{1 Nephi 17:11} % *1NEPHI-17-11 
		% ---
			\label{1NEPHI-17-11}

			And it came to pass that
			I Nephi did make bellowses wherewith to blow the fire
			of the skins of beasts.
			And after that I had made bellowses
			that I might have wherewith to blow the fire,
			I did smite two stones together that I might make fire.

		% ---
		\paragraph{1 Nephi 17:12} % *1NEPHI-17-12 
		% ---
			\label{1NEPHI-17-12}

			For the Lord had not hitherto suffered that we should make much fire
			as we journeyed in the wilderness,
			for he saith:
			I will make that thy food shall become sweet,
			that ye cook it not.

		% ---
		\paragraph{1 Nephi 17:13} % *1NEPHI-17-13 
		% ---
			\label{1NEPHI-17-13}

			And I will also be your light in the wilderness.
			And I will prepare the way before you
			if it so be that ye shall keep my commandments.
			Wherefore inasmuch as ye shall keep my commandments,
			ye shall be led towards the promised land.
			And ye shall know that it is by me that ye are led.

		% ---
		\paragraph{1 Nephi 17:14} % *1NEPHI-17-14 
		% ---
			\label{1NEPHI-17-14}

			Yea, and the Lord said also that
			after ye have arriven to the promised land,
			ye shall know that I the Lord am God
			and that I the Lord did deliver you from destruction,
			yea, that I did bring you out of the land of Jerusalem.

		% ---
		\paragraph{1 Nephi 17:15} % *1NEPHI-17-15 
		% ---
			\label{1NEPHI-17-15}

			Wherefore I Nephi did strive to keep the commandments of the Lord,
			and I did exhort my brethren to faithfulness and diligence.

		% ---
		\paragraph{1 Nephi 17:16} % *1NEPHI-17-16 
		% ---
			\label{1NEPHI-17-16}

			And it came to pass that I did make tools of the ore
			which I did molten out of the rock.

		% ---
		\paragraph{1 Nephi 17:17} % *1NEPHI-17-17 
		% ---
			\label{1NEPHI-17-17}

			And when my brethren saw that I was about to build a ship,
			they began to murmur against me, saying:
			Our brother is a fool,
			for he thinketh that he can build a ship.
			Yea, and he also thinketh that he can cross these great waters.

		% ---
		\paragraph{1 Nephi 17:18} % *1NEPHI-17-18 
		% ---
			\label{1NEPHI-17-18}

			And thus my brethren did complain against me
			and were desirous that they might not labor,
			for they did not believe that I could build a ship,
			neither would they believe that I were instructed of the Lord.

		% ---
		\paragraph{1 Nephi 17:19} % *1NEPHI-17-19 
		% ---
			\label{1NEPHI-17-19}

			And now it came to pass that I Nephi was exceeding sorrowful
			because of the hardness of their hearts.
			And now when they saw that I began to be sorrowful,
			they were glad in their hearts,
			insomuch that they did rejoice over me, saying:
			We knew that ye could not construct a ship,
			for we knew that ye were lacking in judgment;
			wherefore thou canst not accomplish so great a work.

		% ---
		\paragraph{1 Nephi 17:20} % *1NEPHI-17-20 
		% ---
			\label{1NEPHI-17-20}

			And thou art like unto our father,
			led away by the foolish imaginations of his heart.
			Yea, he hath led us out of the land of Jerusalem,
			and we have wandered in the wilderness for these many years.
			And our women have toiled being big with child,
			and they have borne children in the wilderness
			and suffered all things save it were death.
			And it would have been better that they had died
			before they came out of Jerusalem
			than to have suffered these afflictions.

		% ---
		\paragraph{1 Nephi 17:21} % *1NEPHI-17-21 
		% ---
			\label{1NEPHI-17-21}

			Behold, these many years we have suffered in the wilderness,
			which time we might have enjoyed our possessions
			and the land of our inheritance;
			yea, and we might have been happy.

		% ---
		\paragraph{1 Nephi 17:22} % *1NEPHI-17-22 
		% ---
			\label{1NEPHI-17-22}

			And we know that the people which were in the land of Jerusalem
			were a righteous people,
			for they keep the statutes and the judgments of the Lord
			and all his commandments according to the law of Moses;
			wherefore we know that they are a righteous people.
			And our father hath judged them
			and hath led us away because we would hearken unto his word;
			yea, and our brother is like unto him.
			And after this manner of language
			did my brethren murmur and complain against
				\footnote{%
					\label{1NEPHI-17-22-NOTE-us}%
					Who's the ``us''? Sam included? Nephi's wife?
					}%
			us.

		% ---
		\paragraph{1 Nephi 17:23} % *1NEPHI-17-23 
		% ---
			\label{1NEPHI-17-23}

			And it came to pass that I Nephi spake unto them, saying:
			Do ye believe that our fathers, which were the children of Israel,
			would have been led away out of the hands of the Egyptians
			if they had not hearkened unto the words of the Lord?

		% ---
		\paragraph{1 Nephi 17:24} % *1NEPHI-17-24 
		% ---
			\label{1NEPHI-17-24}

			Yea, do ye suppose that they would have been led out of bondage
			if the Lord had not commanded Moses
			that he should lead them out of bondage?

		% ---
		\paragraph{1 Nephi 17:25} % *1NEPHI-17-25 
		% ---
			\label{1NEPHI-17-25}

			Now ye know that the children of Israel were in bondage,
			and ye know that they were laden with tasks
			which were grievious to be borne.
			Wherefore ye know that it must needs be a good thing for them
			that they should be brought out of bondage.

		% ---
		\paragraph{1 Nephi 17:26} % *1NEPHI-17-26 
		% ---
			\label{1NEPHI-17-26}

			Now ye know that Moses was commanded of the Lord to do that great work.
			And ye know that by his word
			the waters of the Red Sea was divided hither and thither,
			and they passed through on dry ground.

		% ---
		\paragraph{1 Nephi 17:27} % *1NEPHI-17-27 
		% ---
			\label{1NEPHI-17-27}

			But ye know that the Egyptians were drowned in the Red Sea,
			which were the armies of Pharaoh.

		% ---
		\paragraph{1 Nephi 17:28} % *1NEPHI-17-28 
		% ---
			\label{1NEPHI-17-28}

			And ye also know that they were fed with manna in the wilderness.

		% ---
		\paragraph{1 Nephi 17:29} % *1NEPHI-17-29 
		% ---
			\label{1NEPHI-17-29}

			Yea, and ye also know that Moses by his word
			according to the power of God which was in him
			smote the rock and there came forth water,
			that the children of Israel might quench their thirst.

		% ---
		\paragraph{1 Nephi 17:30} % *1NEPHI-17-30 
		% ---
			\label{1NEPHI-17-30}

			And notwithstanding they being led,
			the Lord their God, their Redeemer, going before them,
			leading them by day and giving light unto them by night,
			and doing all things for them which was expedient for man to receive,
			they hardened their hearts and blinded their minds
			and reviled against Moses and against the true and living God.

		% ---
		\paragraph{1 Nephi 17:31} % *1NEPHI-17-31 
		% ---
			\label{1NEPHI-17-31}

			And it came to pass that according to his word he did destroy them,
			and according to his word he did lead them,
			and according to his word he did do all things for them.
			And there was not any thing done save it were by his word.

		% ---
		\paragraph{1 Nephi 17:32} % *1NEPHI-17-32 
		% ---
			\label{1NEPHI-17-32}

			And after they had crossed the river Jordan,
			he did make them mighty unto the driving out the children of the land,
			yea, unto the scattering them to destruction.

		% ---
		\paragraph{1 Nephi 17:33} % *1NEPHI-17-33 
		% ---
			\label{1NEPHI-17-33}

			And now do ye suppose that the children of this land,
			which were in the land of promise,
			which were driven out by our fathers,
			do ye suppose that they were righteous?
			Behold, I say unto you: Nay.

		% ---
		\paragraph{1 Nephi 17:34} % *1NEPHI-17-34 
		% ---
			\label{1NEPHI-17-34}

			Do ye suppose that our fathers would have been more choice than they
			if they had been righteous?
			I say unto you: Nay.

		% ---
		\paragraph{1 Nephi 17:35} % *1NEPHI-17-35 
		% ---
			\label{1NEPHI-17-35}

			Behold, the Lord \hyperref[ESTEEM]{esteemeth}%
				\footnote{%
					\label{1NEPHI-17-35-NOTE-esteem}%
					See \hyperref[ESTEEM-1828]{here} for the definitions of ``esteem.'' Meaning \#2 in that list is the one operative here? I'm not sure. At the very least the interpretation of this part of the verse, to me, is not that clear. The first three definitions of ``esteem'' could all work, and when you look at other uses of the word in the BOM---such as those in \hyperref[2NEPHI-33-2]{2 Nephi 33:2}, \hyperref[2NEPHI-33-3]{2 Nephi 33:3}, and \hyperref[MOSIAH-27-4]{Mosiah 27:4}, among others---it's pretty clear that the word is being used with definition \#3, where it just means ``consider'' or ``to think of as.'' But there are other uses, like \hyperref[MOSIAH-14-3]{Mosiah 14:3}, that look like definition \#2 or maybe \#1. So when you get a verse like this one above, 1 Nephi 17:35, without a lot of clues about the right interpretation, it's tough. And none of this analysis even takes into account the tricky later parts of the phrase, ``all flesh in one.''
					See also \hyperref[DC38-16]{DC 38:16}; \hyperref[DC38-24]{DC 38:24--27};
					}
			all flesh in one;
			he that is righteous is favored%
				\footnote{%
					% \label{}%
					See \hyperref[1NEPHI-1-1]{1 Nephi 1:1}, \hyperref[1NEPHI-3-6]{1 Nephi 3:6--8}, \hyperref[MOSIAH-10-13]{Mosiah 10:13}, \hyperref[MOSIAH-24-4]{Mosiah 24:4--7}, \hyperref[MOSIAH-26-1]{Mosiah 26:1--3},
					}
			of God.
			But behold, this people had rejected every word of God,
			and they were ripe in iniquity,
			and the fullness of the wrath of God was upon them.
			And the Lord did curse the land against them
			and bless it unto our fathers.
			Yea, he did curse it against them unto their destruction,
			and he did bless it unto our fathers
			unto their obtaining power over it.

		% ---
		\paragraph{1 Nephi 17:36} % *1NEPHI-17-36 
		% ---
			\label{1NEPHI-17-36}

			Behold, the Lord hath created the earth that it should be inhabited,
			and he hath created his children that they should possess it.%
				\footnote{%
					% \label{}%
					See \hyperref[DC88-20]{DC 88:20}.
					}

		% ---
		\paragraph{1 Nephi 17:37} % *1NEPHI-17-37 
		% ---
			\label{1NEPHI-17-37}

			And he raiseth up a righteous nation and destroyeth the nations of the wicked.

		% ---
		\paragraph{1 Nephi 17:38} % *1NEPHI-17-38 
		% ---
			\label{1NEPHI-17-38}

			And he leadeth away the righteous into precious lands,
			and the wicked he destroyeth and curseth the land unto them for their sakes.

		% ---
		\paragraph{1 Nephi 17:39} % *1NEPHI-17-39 
		% ---
			\label{1NEPHI-17-39}

			He ruleth high in the heavens,
			for it is his throne,
			and this earth is his footstool.

		% ---
		\paragraph{1 Nephi 17:40} % *1NEPHI-17-40 
		% ---
			\label{1NEPHI-17-40}

			And he loveth them which will have him to be their God.
			Behold, he loved our fathers;
			and he covenanted with them,
			yea, even Abraham and Isaac and Jacob,
			and he remembered the covenants which he had made;
			wherefore he did bring them out of the land of Egypt.

		% ---
		\paragraph{1 Nephi 17:41} % *1NEPHI-17-41 
		% ---
			\label{1NEPHI-17-41}

			And he did straiten them in the wilderness with his rod,
			for they hardened their hearts even as ye have.
			And the Lord straitened them because of their iniquity.
			He sent flying fiery serpents among them.
			And after they were bitten,
			he prepared a way that they might be healed.
			And the labor which they had to perform were to look.
			And because of the simpleness of the way or the easiness of it,
			there were many which perished.

		% ---
		\paragraph{1 Nephi 17:42} % *1NEPHI-17-42 
		% ---
			\label{1NEPHI-17-42}

			And they did harden their hearts from time to time,
			and they did revile against Moses and also against God.
			Nevertheless ye know that they were led forth
			by his matchless power into the land of promise.

		% ---
		\paragraph{1 Nephi 17:43} % *1NEPHI-17-43 
		% ---
			\label{1NEPHI-17-43}

			And now after all these things,
			the time has come that they have became wicked,
			yea, nearly unto ripeness.
			And I know not but they are at this day about to be destroyed,
			for I know that the day must surely come
			that they must be destroyed save a few only
			which shall be led away into captivity.

		% ---
		\paragraph{1 Nephi 17:44} % *1NEPHI-17-44 
		% ---
			\label{1NEPHI-17-44}

			Wherefore the Lord commanded my father
			that he should depart into the wilderness.
			And the Jews also sought to take away his life.
			Yea, and ye also have sought to take away his life.
			Wherefore ye are murderers in your hearts,
			and ye are like unto they.

		% ---
		\paragraph{1 Nephi 17:45} % *1NEPHI-17-45 
		% ---
			\label{1NEPHI-17-45}

			Ye are swift to do iniquity
			but slow to remember the Lord your God.
			Ye have seen an angel and he spake unto you.
			Yea, ye have heard his voice from time to time,
			and he hath spoken unto you in a still small voice;
			but ye were
				\footnote{%
					% \label{}%
					For ``past feeling,'' see \hyperref[EPHESIANS-4-19]{Ephesians 4:19} and \hyperref[MORONI-9-20]{Moroni 9:20}, which are the only other places that this phrase occurs.
					}%
			past feeling, that ye could not feel his words.
			Wherefore he hath spoken unto you like unto the voice of thunder,
			which did cause the earth to shake as if it were to divide asunder.

		% ---
		\paragraph{1 Nephi 17:46} % *1NEPHI-17-46 
		% ---
			\label{1NEPHI-17-46}

			And ye also know that by the power of his almighty word
			he can cause the earth that it shall pass away.
			Yea, and ye know that by his word he can cause
			that rough places be made smooth and smooth places shall be broken up.
			O then why is it that ye can be so hard in your hearts?

		% ---
		\paragraph{1 Nephi 17:47} % *1NEPHI-17-47 
		% ---
			\label{1NEPHI-17-47}

			Behold, my soul is rent with anguish because of you,
			and my heart is pained.
			I fear lest ye shall be cast off forever.
			Behold, I am full of the Spirit of God,
			insomuch as if my frame had no strength.

		% ---
		\paragraph{1 Nephi 17:48} % *1NEPHI-17-48 
		% ---
			\label{1NEPHI-17-48}

			And now it came to pass that when I had spoken these words,
			they were angry with me and were desirous
			to throw me into the depths of the sea.
			And as they came forth to lay their hands upon me,
			I spake unto them, saying:
			In the name of the Almighty God I command you that ye touch me not,
			for I am filled with the power of God,
			even unto the consuming of my flesh.%
				\footnote{%
					% \label{1NEPHI-17-48-consuming}%
					See \hyperref[2NEPHI-4-21]{2 Nephi 4:21}.
					}
			And whoso shall lay their hands upon me
			shall wither even as a dried reed,
			and he shall be as naught before the power of God,
			for God shall smite him.

		% ---
		\paragraph{1 Nephi 17:49} % *1NEPHI-17-49 
		% ---
			\label{1NEPHI-17-49}

			And it came to pass that I Nephi saith unto them
			that they should murmur no more against their father,
			neither should they withhold their labor from me,
			for God had commanded me that I should build a ship.

		% ---
		\paragraph{1 Nephi 17:50} % *1NEPHI-17-50 
		% ---
			\label{1NEPHI-17-50}

			And I saith unto them:
			If God had commanded me to do all things,
			I could do it.
			If he should command me
			that I should say unto this water:
			Be thou earth!
			---and it shall be earth.
			And if I should say it,
			it would be done.

		% ---
		\paragraph{1 Nephi 17:51} % *1NEPHI-17-51 
		% ---
			\label{1NEPHI-17-51}

			And now if the Lord hath such great power
			and hath wrought so many miracles among the children of men,
			how is it that he cannot instruct me that I should build a ship?

		% ---
		\paragraph{1 Nephi 17:52} % *1NEPHI-17-52 
		% ---
			\label{1NEPHI-17-52}

			And it came to pass that I Nephi said many things unto my brethren,
			insomuch that they were confounded and could not contend against me,
			neither durst they lay their hands upon me nor touch me with their fingers,
			even for the space of many days.
			Now they durst not do this lest they should wither before me,
			so powerful was the Spirit of God.
			And thus it had wrought upon them.

		% ---
		\paragraph{1 Nephi 17:53} % *1NEPHI-17-53 
		% ---
			\label{1NEPHI-17-53}

			And it came to pass that the Lord said unto me:
			Stretch forth thine hand again unto thy brethren.
			And they shall not wither before thee,
			but I will shake%
				\footnote{%
					\label{1NEPHI-17-53-NOTE-shake}%
					The current edition of the BOM has ``shock'' here. It is the only use of ``shock'' in the LDS scriptures where the word appears as a verb; it does appear twice in the OT, but there it's a noun referring to parts of corn. Here are the definitions from the 1828:
						\begin{compactitem}
							\item[1] To shake by the sudden collision of a body.
							\item[2] To meet with force; to encounter.
							\item[3] To strike, as with horror or disgust; to cause to recoil, as from something odious or horrible; to offend extremely; to disgust.
						\end{compactitem}
					The word ``shock'' could work here, but with the ``shake'' you get in the other verses, you can see how they could make this mistake.
					}
			them, saith the Lord.
			And this will I do that they may know that I am the Lord their God.

		% ---
		\paragraph{1 Nephi 17:54} % *1NEPHI-17-54 
		% ---
			\label{1NEPHI-17-54}

			And it came to pass that I stretched forth my hand unto my brethren.
			And they did not wither before me, but the Lord did shake them,
			even according to the word which he had spoken.

		% ---
		\paragraph{1 Nephi 17:55} % *1NEPHI-17-55 
		% ---
			\label{1NEPHI-17-55}

			And now they said:
			We know of a surety that the Lord is with thee,
			for we know that it is the power of the Lord that hath shaken us.
			And they fell down before me and were about to worship me,
			but I would not suffer them, saying:
			I am thy brother, yea, even thy younger brother.
			Wherefore worship the Lord thy God
			and honor thy father and thy mother,
			that thy days may be long in the land
			which the Lord thy God shall give thee.

	% -----
	\subsubsection{1 Nephi 18} % *1NEPHI-18 
	% -----
		\label{1NEPHI-18}
		
		% ---
		\paragraph{1 Nephi 18:1} % *1NEPHI-18-1 
		% ---
			\label{1NEPHI-18-1}

			And it came to pass that
			they did worship the Lord and did go forth with me,
			and we did work timbers of curious workmanship.
			And the Lord did shew me from time to time
			after what manner I should work the timbers of the ship.

		% ---
		\paragraph{1 Nephi 18:2} % *1NEPHI-18-2 
		% ---
			\label{1NEPHI-18-2}

			Now I Nephi did not work the timbers after the manner which was learned by men,
			neither did I build the ship after the manner of men,
			but I did build it after the manner which the Lord had shewn unto me;
			wherefore it was not after the manner of men.%
				\footnote{%
					\label{1NEPHI-18-2-NOTE-men}%
					As the Lord leads us into ``precious lands'' (\hyperref[1NEPHI-17-38]{1 Nephi 17:38}) and into our own promised land, we will need to do things differently from how other people have done them. We will become what we need to be through ways that we do not currently understand. It may not be after the manner that it's been done before, or after the manner in which people think it should be done, but the end result will be to land us in the land that is precious to us.
					}

		% ---
		\paragraph{1 Nephi 18:3} % *1NEPHI-18-3 
		% ---
			\label{1NEPHI-18-3}

			And I Nephi did go into the mount oft,
			and I did pray oft unto the Lord;
			wherefore the Lord shewed unto me great things.

		% ---
		\paragraph{1 Nephi 18:4} % *1NEPHI-18-4 
		% ---
			\label{1NEPHI-18-4}

			And it came to pass that after I had finished the ship
			according to the word of the Lord,
			my brethren beheld that it was good
			and that the workmanship thereof was exceeding fine;
			wherefore they did humble themselves again before the Lord.

		% ---
		\paragraph{1 Nephi 18:5} % *1NEPHI-18-5 
		% ---
			\label{1NEPHI-18-5}

			And it came to pass that the voice of the Lord came unto my father
			that we should arise and go down into the ship.

		% ---
		\paragraph{1 Nephi 18:6} % *1NEPHI-18-6 
		% ---
			\label{1NEPHI-18-6}

			And it came to pass that on the morrow,
			after that we had prepared all things,
			much fruits and meat from the wilderness
			and honey in abundance and provisions,
			according to that which the Lord had commanded us,
			we did go down into the ship with all our loading and our seeds
			and whatsoever things we had brought with us,
			every one according to his age;
			wherefore we did all go down into the ship with our wives and our children.

		% ---
		\paragraph{1 Nephi 18:7} % *1NEPHI-18-7 
		% ---
			\label{1NEPHI-18-7}

			And now my father had begat two sons in the wilderness;
				\footnote{%
					% \label{}%
					The names come from Lehi's discovery of his ancestors in \hyperref[1NEPHI-5-14]{1 Nephi 5:14}.
					}%
			the elder was called Jacob and the younger Joseph.

		% ---
		\paragraph{1 Nephi 18:8} % *1NEPHI-18-8 
		% ---
			\label{1NEPHI-18-8}

			And it came to pass that after we had all gone down into the ship
			and had taken with us our provisions and things which had been commanded us,
			we did put forth into the sea and were driven forth before the wind
			towards the promised land.

		% ---
		\paragraph{1 Nephi 18:9} % *1NEPHI-18-9 
		% ---
			\label{1NEPHI-18-9}

			And after that we had been driven forth before the wind
			for the space of many days,
			behold, my brethren and the sons of Ishmael and also their wives
			began to make themselves merry,
			insomuch that they began to dance and to sing
			and to speak with much rudeness,
			yea, even to that they did forget by what power they had been brought thither;
			yea, they were lifted up unto exceeding rudeness.%
				\footnote{%
					\label{1NEPHI-18-9-NOTE-rudeness}%
					``Rudeness'' only appears in the scriptures three times: twice here, and once at \hyperref[2NEPHI-2-1]{2 Nephi 2:1}, where Lehi is talking to Jacob. Here are the 1828 definitions:
						\begin{compactitem}
							\item[1] A rough broken state; unevenness; wildness; as the rudeness of a mountain, country or landscape.
							\item[2] Coarseness of manners; incivility; rusticity; vulgarity.
							\item[3] Ignorance; unskillfulness.
							\item[4] Artlessness; coarseness; inelegance; as the rudeness of a painting or piece of sculpture.
							\item[5] Violence; impetuosity; as the rudeness of an attack or shock.
							\item[6] Violence; storminess; as the rudeness of winds or of the season.
						\end{compactitem}
					It looks like the only definition that could be relevant here is \#2. Given Nephi's about propriety, I think it's hard to know exactly what they were doing on the ship that was the problem. As you see below in \hyperref[1NEPHI-18-12]{verse 12}, the compass only stops working after they tie Nephi up---it doesn't stop after earlier, when they are ``making merry.''
					}

		% ---
		\paragraph{1 Nephi 18:10} % *1NEPHI-18-10 
		% ---
			\label{1NEPHI-18-10}

			And I Nephi began to fear exceedingly lest the Lord should be angry with us
			and smite us because of our iniquity,
			that we should be swallowed up in the depths of the sea;
			wherefore I Nephi began to speak to them with much soberness.
			But behold, they were angry with me, saying:
			We will not that our younger brother shall be a ruler over us.

		% ---
		\paragraph{1 Nephi 18:11} % *1NEPHI-18-11 
		% ---
			\label{1NEPHI-18-11}

			And it came to pass that Laman and Lemuel did take me and bind me with cords,
			and they did treat me with much harshness.
			Nevertheless the Lord suffered it that he might shew forth his power
			unto the fulfilling of his word which he hath spoken concerning the wicked.

		% ---
		\paragraph{1 Nephi 18:12} % *1NEPHI-18-12 
		% ---
			\label{1NEPHI-18-12}

			And it came to pass that after they had bound me,
			insomuch that I could not move,
			the compass which had been prepared of the Lord did cease to work;%
				\footnote{%
					\label{1NEPHI-18-12-NOTE-cease}%
					I think it's really important to note that the compass only stops working\textit{after they tie Nephi up}, and \textit{not} as they are ``making merry'' or being ``lifted up unto exceeding rudeness.'' So it makes me wonder, if Nephi hadn't said anything, and had just left them alone, would the compass have stopped working? Would they have gotten lost? Would the storm have come up? I don't think we know, but the exact order of events is telling.
					}

		% ---
		\paragraph{1 Nephi 18:13} % *1NEPHI-18-13 
		% ---
			\label{1NEPHI-18-13}

			wherefore they knew not whither they should steer the ship,
			insomuch that there arose a great storm,
			yea, a great and terrible tempest,
			and we were driven back upon the waters for the space of three days.%
				\footnote{%
					\label{1NEPHI-18-13-NOTE-storm}%
					To me, this reads as though Laman and Lemuel and the others, without the benefit of a working compass, steer the ship directly into the storm. It's more like, the storm was already there, and had the compass been working, they would have avoided it, but since it wasn't working, they just didn't know which direction to go and headed into the storm. This also teaches us something interesting about the way the Lord is guiding them to the land of promise---it's not a direct straight line toward it, necessarily, because maybe going straight for it puts you in the middle of some storms. Maybe the Lord takes a circuitous route sometimes to avoid some storms.
					}
			And they began to be frightened exceedingly
			lest they should be drowned in the sea.
			Nevertheless they did loose me not.

		% ---
		\paragraph{1 Nephi 18:14} % *1NEPHI-18-14 
		% ---
			\label{1NEPHI-18-14}

			And on the fourth day which we had been driven back,
			the tempest began to be exceeding sore.

		% ---
		\paragraph{1 Nephi 18:15} % *1NEPHI-18-15 
		% ---
			\label{1NEPHI-18-15}

			And it came to pass that we were about to be swallowed up in the depths of the sea.
			And after that we had been driven back upon the waters for the space of four days,
			my brethren began to see that the judgments of God was upon them
			and that they must perish save that they should repent of their iniquities.
			Wherefore they came unto me and loosed the bands which was upon my wrists,
			and behold, they had swollen exceedingly;
			and also mine ankles were much swollen,
			and great was the soreness thereof.

		% ---
		\paragraph{1 Nephi 18:16} % *1NEPHI-18-16 
		% ---
			\label{1NEPHI-18-16}

			Nevertheless I did look unto my God
			and I did praise him all the day long,
			and I did not murmur against the Lord because of mine afflictions.

		% ---
		\paragraph{1 Nephi 18:17} % *1NEPHI-18-17 
		% ---
			\label{1NEPHI-18-17}

			Now my father Lehi had said many things unto them
			and also unto the sons of Ishmael,
			but behold, they did breathe out much threatenings
			against any one that should speak for me.
			And my parents being stricken in years
			and having suffered much grief because of their children,
			they were brought down, yea, even upon their sickbeds.

		% ---
		\paragraph{1 Nephi 18:18} % *1NEPHI-18-18 
		% ---
			\label{1NEPHI-18-18}

			Because of their grief and much sorrow and the iniquity of my brethren,
			they were brought near even to be carried out of this time to meet their God.
			Yea, their gray hairs were about to be brought down to lie low in the dust;
			yea, even they were near to be cast with sorrow into a watery grave.

		% ---
		\paragraph{1 Nephi 18:19} % *1NEPHI-18-19 
		% ---
			\label{1NEPHI-18-19}

			And Jacob and Joseph also being young,
			having need of much nourishment,
			were grieved because of the afflictions of their mother.
			And also my wife with her tears and prayers and also my children
			did not soften the hearts of my brethren
			that they would loose me.

		% ---
		\paragraph{1 Nephi 18:20} % *1NEPHI-18-20 
		% ---
			\label{1NEPHI-18-20}

			And there was nothing save it were the power of God
			which threatened them with destruction
			could soften their hearts.
			Wherefore when they saw that
			they were about to be swallowed up in the depths of the sea,
			they repented of the thing which they had done,
			insomuch that they loosed me.%
				\footnote{%
					\label{1NEPHI-18-20-NOTE-destruction}%
					There is something good about repenting here, even though it takes them a lot of time to get to it. In the life of every person there are times when we continue in a certain behavior until the time when we are ``threatened with destruction.'' Most of the time it's not as dramatic as this story here, and it instead has to do with the harmful consequences of individual behaviors, when we kind of bottom out with a certain activity and realize that we need to make a change. The trick is to learn to make those moments occur less and less often over time.
					}

		% ---
		\paragraph{1 Nephi 18:21} % *1NEPHI-18-21 
		% ---
			\label{1NEPHI-18-21}

			And it came to pass that after they had loosed me,
			behold, I took the compass and it did work whither I desired it.
			And it came to pass that I prayed unto the Lord;
			and after that I had prayed,
			the winds did cease and the storm did cease and there was a great calm.

		% ---
		\paragraph{1 Nephi 18:22} % *1NEPHI-18-22 
		% ---
			\label{1NEPHI-18-22}

			And it came to pass that I Nephi did guide the ship,
			that we sailed again towards the promised land.

		% ---
		\paragraph{1 Nephi 18:23} % *1NEPHI-18-23 
		% ---
			\label{1NEPHI-18-23}

			And it came to pass that after we had sailed for the space of many days,
			we did arrive to the promised land.
			And we went forth upon the land and did pitch our tents,
			and we did call it the promised land.

		% ---
		\paragraph{1 Nephi 18:24} % *1NEPHI-18-24 
		% ---
			\label{1NEPHI-18-24}

			And it came to pass that we did begin to till the earth,
			and we began to plant seeds;
			yea, we did put all our seeds into the earth
			which we had brought from the land of Jerusalem.
			And it came to pass that they did grow exceedingly;
			wherefore we were blessed in abundance.

		% ---
		\paragraph{1 Nephi 18:25} % *1NEPHI-18-25 
		% ---
			\label{1NEPHI-18-25}

			And it came to pass that we did find upon the land of promise
			as we journeyed in the wilderness
			that there was beasts in the forests of every kind,
			both the cow and the ox and the ass and the horse and the goat and the wild goat
			and all manner of wild animals which were for the use of man.
			And we did find all manner of ore, both of gold and of silver and of copper.

	% -----
	\subsubsection{1 Nephi 19} % *1NEPHI-19
	% -----
		\label{1NEPHI-19}
		
		% ---
		\paragraph{1 Nephi 19:1} % *1NEPHI-19-1 
		% ---
			\label{1NEPHI-19-1}

			And it came to pass that the Lord commanded me,
			wherefore I did make plates of ore
			that I might engraven upon them the record of my people.
			And upon the plates which I made I did engraven the record of my father
			and also our journeyings in the wilderness and the prophecies of my father.
			And also many of mine own prophecies have I engraven upon them.

		% ---
		\paragraph{1 Nephi 19:2} % *1NEPHI-19-2 
		% ---
			\label{1NEPHI-19-2}

			And I knew not at that time which I made them
			that I should be commanded of the Lord to make these plates.
			Wherefore the record of my father and the genealogy of his forefathers
			and the more part of all our proceedings in the wilderness
			are engraven upon those first plates of which I have spoken.
			Wherefore the things which transpired before that I made these plates
			are of a truth more particularly made mention upon the first plates.

		% ---
		\paragraph{1 Nephi 19:3} % *1NEPHI-19-3 
		% ---
			\label{1NEPHI-19-3}

			And after that I made these plates by way of commandment,
			I Nephi received a commandment
			that the ministry and the prophecies
			---the more plain and precious parts of them---
			should be written upon these plates,
			and that the things which were written
			should be kept for the instruction of my people,
			which should possess the land,
			and also for other wise purposes,
			which purposes are known unto the Lord.

		% ---
		\paragraph{1 Nephi 19:4} % *1NEPHI-19-4 
		% ---
			\label{1NEPHI-19-4}

			Wherefore I Nephi did make a record upon the other plates,
			which gives an account or which gives a greater account
			of the wars and contentions and destructions of my people.
			And now this have I done and commanded my people
			that they should do after that I was gone
			and that these plates should be handed down
			from one generation to another or from one prophet to another
			until further commandments of the Lord.

		% ---
		\paragraph{1 Nephi 19:5} % *1NEPHI-19-5 
		% ---
			\label{1NEPHI-19-5}

			And an account of my making these plates shall be given hereafter.
			And then behold, I proceed according to that which I have spoken;
			and this I do that the more sacred things may be kept
			for the knowledge of my people.

		% ---
		\paragraph{1 Nephi 19:6} % *1NEPHI-19-6 
		% ---
			\label{1NEPHI-19-6}

			Nevertheless I do not write any thing upon plates
			save it be that I think it be sacred.
			And now if I do err,
			even did they err of old---
			not that I would excuse myself because of other men,
			but because of the weakness which is in me according to the flesh,
			I would excuse myself.

		% ---
		\paragraph{1 Nephi 19:7} % *1NEPHI-19-7 
		% ---
			\label{1NEPHI-19-7}

			For the things which some men esteem to be of great worth,
			both to the body and soul,
			others set at naught and trample under their feet,
			yea, even the very God of Israel do men trample under their feet.
			I say trample under their feet,
			but I would speak in other words:
			they do set him at naught and hearken not to the voice of his counsels.

		% ---
		\paragraph{1 Nephi 19:8} % *1NEPHI-19-8 
		% ---
			\label{1NEPHI-19-8}

			And behold, he cometh according to the words of the angel
			in six hundred years from the time my father left Jerusalem.

		% ---
		\paragraph{1 Nephi 19:9} % *1NEPHI-19-9 
		% ---
			\label{1NEPHI-19-9}

			And the world because of their iniquity shall judge him to be a thing of naught.
			Wherefore they scourge him and he suffereth it;
			and they smite him and he suffereth it;
			yea, they spit upon him and he suffereth it
			because of his loving-kindness and his long-suffering towards the children of men.

		% ---
		\paragraph{1 Nephi 19:10} % *1NEPHI-19-10 
		% ---
			\label{1NEPHI-19-10}

			And the God of our fathers, which were led out of Egypt out of bondage,
			and also were preserved in the wilderness by him,
			yea, the God of Abraham and of Isaac and the God of Jacob yieldeth himself
			according to the words of the angel
			as a man into the hands of wicked men,
			to be lifted up according to the words of Zenoch,
			and to be crucified according to the words of Neum,
			and to be buried in a sepulchre according to the words of Zenos,
			which he spake concerning the three days of darkness
			which should be a sign given of his death unto them
			who should inhabit the isles of the sea,
			more especially given unto them which are of the house of Israel.

		% ---
		\paragraph{1 Nephi 19:11} % *1NEPHI-19-11 
		% ---
			\label{1NEPHI-19-11}

			For thus spake the prophet:
			The Lord God surely shall visit all the house of Israel at that day,
			some with his voice because of their righteousness,
			unto their great joy and salvation,
			and others with the thunderings and the lightnings of his power,
			by tempest, by fire, and by smoke and vapor of darkness
			and by the opening of the earth
			and by mountains which shall be carried up.

		% ---
		\paragraph{1 Nephi 19:12} % *1NEPHI-19-12 
		% ---
			\label{1NEPHI-19-12}

			And all these things must surely come, saith the prophet Zenos,
			and the rocks of the earth must rend.
			And because of the groanings of the earth,
			many of the kings of the isles of the sea
			shall be wrought upon by the Spirit of God to exclaim:
			The God of nature suffers.

		% ---
		\paragraph{1 Nephi 19:13} % *1NEPHI-19-13 
		% ---
			\label{1NEPHI-19-13}

			And as for they which are at Jerusalem, saith the prophet,
			they shall be scourged by all people, saith the prophet,
			because they crucified the God of Israel and turned their hearts aside,
			rejecting signs and wonders and power and glory of the God of Israel.

		% ---
		\paragraph{1 Nephi 19:14} % *1NEPHI-19-14 
		% ---
			\label{1NEPHI-19-14}

			And because they have turned their hearts aside, saith the prophet,
			and have despised the Holy One of Israel,
			they shall wander in the flesh and perish
			and become a hiss and a byword and be hated among all nations.

		% ---
		\paragraph{1 Nephi 19:15} % *1NEPHI-19-15 
		% ---
			\label{1NEPHI-19-15}

			Nevertheless when that day cometh, saith the prophet,
			that they no more turn aside their hearts against the Holy One of Israel,
			then will he remember the covenants which he made to their fathers.

		% ---
		\paragraph{1 Nephi 19:16} % *1NEPHI-19-16 
		% ---
			\label{1NEPHI-19-16}

			Yea, then will he remember the isles of the sea;
			yea, and all the people which are of the house of Israel
			will I gather in, saith the Lord,
			according to the words of the prophet Zenos,
			from the four quarters of the earth.

		% ---
		\paragraph{1 Nephi 19:17} % *1NEPHI-19-17 
		% ---
			\label{1NEPHI-19-17}

			Yea, and all the earth shall see the salvation of the Lord, saith the prophet;
			every nation, kindred, tongue, and people shall be blessed.

		% ---
		\paragraph{1 Nephi 19:18} % *1NEPHI-19-18 
		% ---
			\label{1NEPHI-19-18}

			And I Nephi have written these things unto my people
			that perhaps I might persuade them
			that they would remember the Lord their Redeemer.

		% ---
		\paragraph{1 Nephi 19:19} % *1NEPHI-19-19 
		% ---
			\label{1NEPHI-19-19}

			Wherefore I speak unto all the house of Israel,
			if it so be that they should obtain these things.

		% ---
		\paragraph{1 Nephi 19:20} % *1NEPHI-19-20 
		% ---
			\label{1NEPHI-19-20}

			For behold, I have workings in the spirit which doth weary me,
			even that all my joints are weak,
			for they which are at Jerusalem.
			For had not the Lord been merciful
			to shew unto me concerning them
			even as he had prophets of old---

		% ---
		\paragraph{1 Nephi 19:21} % *1NEPHI-19-21 
		% ---
			\label{1NEPHI-19-21}

			for he surely did shew unto prophets of old all things concerning them.
			And also he did shew unto many concerning us;
			wherefore it must needs be that we know concerning them,
			for they are written upon the plates of brass.

		% ---
		\paragraph{1 Nephi 19:22} % *1NEPHI-19-22 
		% ---
			\label{1NEPHI-19-22}

			Now it came to pass that I Nephi did teach my brethren these things.
			And it came to pass that I did read many things to them
			which were engraven upon the plates of brass,
			that they might know concerning the doings of the Lord
			in other lands among people of old.

		% ---
		\paragraph{1 Nephi 19:23} % *1NEPHI-19-23 
		% ---
			\label{1NEPHI-19-23}

			And I did read many things unto them which were in the books of Moses.
			But that I might more fully persuade them to believe in the Lord their Redeemer,
			wherefore I did read unto them that which was written by the prophet Isaiah;
			for I did liken%
				\footnote{%
					\label{1NEPHI-19-23-NOTE-liken}%
					See \hyperref[LIKEN]{Liken}, especially the \hyperref[LIKEN-1828]{1828} definition (there is only one) and the \hyperref[LIKEN-OED]{OED} definitions.
					}
			all scriptures unto us,
			that it%
				\footnote{%
					\label{1NEPHI-19-23-NOTE-it}%
					What is the referent of this pronoun? The teaching he does in this way? The likening of the scriptures? The learning from the scriptures? I'm not sure.
					}
			might be for our profit and learning.%
				\footnote{%
					% \label{}%
					This combination of words, ``profit and learning,'' occurs here and at \hyperref[2NEPHI-2-14]{2 Nephi 2:14}, \hyperref[2NEPHI-4-15]{2 Nephi 4:15}, \hyperref[2NEPHI-9-28]{2 Nephi 9:28}, and \hyperref[DC46-1]{DC 46:1}.
					}

		% ---
		\paragraph{1 Nephi 19:24} % *1NEPHI-19-24 
		% ---
			\label{1NEPHI-19-24}

			Wherefore I spake unto them, saying:
			Hear ye the words of the prophet,
			ye which are a remnant of the house of Israel,
			a branch which have been broken off.
			Hear ye the words of the prophet
			which was written unto all the house of Israel
			and liken it unto yourselves,
			that ye may have hope as well as your brethren
			from whom ye have been broken off.
			For after this manner hath the prophet written:

	% -----
	\subsubsection{1 Nephi 20} % *1NEPHI-20 
	% -----
		\label{1NEPHI-20} 
		
		% ---
		\paragraph{1 Nephi 20:1} % *1NEPHI-20-1 
		% ---
			\label{1NEPHI-20-1}

			Hearken and hear this, O house of Jacob,
			which are called by the name of Israel
			and are come forth out of the waters of Judah,
			which swear by the name of the Lord
			and make mention of the God of Israel;
			yet they swear not in truth nor in righteousness.

		% ---
		\paragraph{1 Nephi 20:2} % *1NEPHI-20-2 
		% ---
			\label{1NEPHI-20-2}

			Nevertheless they call themselves of the holy city,
			but they do not stay themselves upon the God of Israel,
			which is the Lord of Hosts;
			yea, the Lord of Hosts is his name.

		% ---
		\paragraph{1 Nephi 20:3} % *1NEPHI-20-3 
		% ---
			\label{1NEPHI-20-3}

			Behold, I have declared the former things from the beginning.
			And they went forth out of my mouth and I shewed them.
			I did shew them suddenly.

		% ---
		\paragraph{1 Nephi 20:4} % *1NEPHI-20-4 
		% ---
			\label{1NEPHI-20-4}

			And I did it because I knew that thou art obstinate,
			and thy neck was an iron sinew and thy brow brass.

		% ---
		\paragraph{1 Nephi 20:5} % *1NEPHI-20-5 
		% ---
			\label{1NEPHI-20-5}

			And I have even from the beginning declared to thee;
			before it came to pass, I shewed them thee.
			And I shewed them for fear lest thou shouldest say:
			Mine idol hath done them,
			and my graven image and my molten image hath commanded them.

		% ---
		\paragraph{1 Nephi 20:6} % *1NEPHI-20-6 
		% ---
			\label{1NEPHI-20-6}

			Thou hast heard and seen all this,
			and will ye not declare them?
			And that I have shewed thee new things from this time,
			even hidden things,
			and thou didst not know them.

		% ---
		\paragraph{1 Nephi 20:7} % *1NEPHI-20-7 
		% ---
			\label{1NEPHI-20-7}

			They are created now and not from the beginning.
			Even before the day when thou heardest them not,
			they were declared unto thee,
			lest thou shouldst say:
			Behold, I knew them.

		% ---
		\paragraph{1 Nephi 20:8} % *1NEPHI-20-8 
		% ---
			\label{1NEPHI-20-8}

			Yea, and thou heardest not,
			yea, thou knewest not;
			yea, from that time thine ear was not opened.
			For I knew that thou wouldst deal very treacherously
			and wast called a transgressor from the womb.

		% ---
		\paragraph{1 Nephi 20:9} % *1NEPHI-20-9 
		% ---
			\label{1NEPHI-20-9}

			Nevertheless for my name’s sake will I defer mine anger.
			And for my praise will I refrain from thee,
			that I cut thee not off.

		% ---
		\paragraph{1 Nephi 20:10} % *1NEPHI-20-10 
		% ---
			\label{1NEPHI-20-10}

			For behold, I have refined thee;
			I have chosen thee in the furnace of affliction.

		% ---
		\paragraph{1 Nephi 20:11} % *1NEPHI-20-11 
		% ---
			\label{1NEPHI-20-11}

			For mine own sake—yea, for mine own sake—will I do this.
			For how should I suffer my name to be polluted?
			And I will not give my glory unto another.

		% ---
		\paragraph{1 Nephi 20:12} % *1NEPHI-20-12 
		% ---
			\label{1NEPHI-20-12}

			Hearken unto me, O Jacob and Israel my called.
			For I am he, and I am the first and I am also the last.

		% ---
		\paragraph{1 Nephi 20:13} % *1NEPHI-20-13 
		% ---
			\label{1NEPHI-20-13}

			Mine hand hath also laid the foundation of the earth,
			and my right hand hath spanned the heavens.
			And I called unto them and they stand up together.

		% ---
		\paragraph{1 Nephi 20:14} % *1NEPHI-20-14 
		% ---
			\label{1NEPHI-20-14}

			All ye, assemble yourselves and hear.
			Which among them hath declared these things unto them?
			The Lord hath loved him.
			Yea, and he will fulfill his word which he hath declared by them.
			And he will do his pleasure on Babylon,
			and his arm shall come upon the Chaldeans.

		% ---
		\paragraph{1 Nephi 20:15} % *1NEPHI-20-15 
		% ---
			\label{1NEPHI-20-15}

			Also saith the Lord:
			I the Lord, yea, I have spoken.
			Yea, I have called him to declare;
			I have brought him,
			and he shall make his way prosperous.

		% ---
		\paragraph{1 Nephi 20:16} % *1NEPHI-20-16 
		% ---
			\label{1NEPHI-20-16}

			Come ye near unto me.
			I have not spoken in secret from the beginning;
			from the time that it was declared have I spoken.
			And the Lord God and his Spirit hath sent me.

		% ---
		\paragraph{1 Nephi 20:17} % *1NEPHI-20-17 
		% ---
			\label{1NEPHI-20-17}

			And thus saith the Lord thy Redeemer, the Holy One of Israel:
			I have sent him.
			The Lord thy God,
			which teacheth thee to profit,
			which leadeth thee by the way thou shouldst go,
			hath done it.

		% ---
		\paragraph{1 Nephi 20:18} % *1NEPHI-20-18 
		% ---
			\label{1NEPHI-20-18}

			O that thou hadst hearkened to my commandments!
			Then had thy peace been as a river
			and thy righteousness as the waves of the sea.

		% ---
		\paragraph{1 Nephi 20:19} % *1NEPHI-20-19 
		% ---
			\label{1NEPHI-20-19}

			Thy seed also had been as the sand,
			the offspring of thy bowels like the gravel thereof.
			His name should not have been cut off,
			nor destroyed from before me.

		% ---
		\paragraph{1 Nephi 20:20} % *1NEPHI-20-20 
		% ---
			\label{1NEPHI-20-20}

			Go ye forth of Babylon;
			flee ye from the Chaldeans.
			With a voice of singing declare ye, tell this;
			utter to the end of the earth, say ye:
			The Lord hath redeemed his servant Jacob.

		% ---
		\paragraph{1 Nephi 20:21} % *1NEPHI-20-21 
		% ---
			\label{1NEPHI-20-21}

			And they thirsted not.
			He led them through the deserts.
			He caused the waters to flow out of the rock for them.
			He clave the rock also and the waters gushed out.

		% ---
		\paragraph{1 Nephi 20:22} % *1NEPHI-20-22 
		% ---
			\label{1NEPHI-20-22}

			And notwithstanding he hath done all this and greater also,
			there is no peace, saith the Lord, unto the wicked.

	% -----
	\subsubsection{1 Nephi 21} % *1NEPHI-21 
	% -----
		\label{1NEPHI-21}
		
		% ---
		\paragraph{1 Nephi 21:1} % *1NEPHI-21-1 
		% ---
			\label{1NEPHI-21-1}

			And again, hearken, O ye house of Israel,
			all ye that are broken off and are driven out
			because of the wickedness of the pastors of my people,
			yea, all ye that are broken off,
			that are scattered abroad,
			which are of my people, O house of Israel.
			Listen, O isles, unto me,
			and hearken, ye people, from far.
			The Lord hath called me from the womb;
			from the bowels of my mother hath he made mention of my name.

		% ---
		\paragraph{1 Nephi 21:2} % *1NEPHI-21-2 
		% ---
			\label{1NEPHI-21-2}

			And he hath made my mouth like a sharp sword.
			In the shadow of his hand hath he hid me
			and made me a polished shaft;
			in his quiver hath he hid me—

		% ---
		\paragraph{1 Nephi 21:3} % *1NEPHI-21-3 
		% ---
			\label{1NEPHI-21-3}

			and said unto me:
			Thou art my servant, O Israel,
			in whom I will be glorified.

		% ---
		\paragraph{1 Nephi 21:4} % *1NEPHI-21-4 
		% ---
			\label{1NEPHI-21-4}

			Then I said:
			I have labored in vain.
			I have spent my strength for naught and in vain.
			Surely my judgment is with the Lord
			and my work with my God.

		% ---
		\paragraph{1 Nephi 21:5} % *1NEPHI-21-5 
		% ---
			\label{1NEPHI-21-5}

			And now saith the Lord that formed me from the womb
			that I should be his servant to bring Jacob again to him.
			Though Israel be not gathered,
			yet shall I be glorious in the eyes of the Lord,
			and my God shall be my strength.

		% ---
		\paragraph{1 Nephi 21:6} % *1NEPHI-21-6 
		% ---
			\label{1NEPHI-21-6}

			And he said:
			It is a light thing
			that thou shouldst be my servant,
			to raise up the tribes of Jacob
			and to restore the preserved of Israel.
			I will also give thee for a light to the Gentiles,
			that thou mayest be my salvation unto the ends of the earth.

		% ---
		\paragraph{1 Nephi 21:7} % *1NEPHI-21-7 
		% ---
			\label{1NEPHI-21-7}

			Thus saith the Lord, the Redeemer of Israel, his Holy One,
			to him whom man despiseth,
			to him whom the nation abhorreth,
			to a servant of rulers:
			Kings shall see and arise,
			princes also shall worship,
			because of the Lord that is faithful.

		% ---
		\paragraph{1 Nephi 21:8} % *1NEPHI-21-8 
		% ---
			\label{1NEPHI-21-8}

			Thus saith the Lord:
			In an acceptable time have I heard thee, O isles of the sea;
			and in a day of salvation have I helped thee.
			And I will preserve thee and give thee my servant
			for a covenant of the people,
			to establish the earth,
			to cause to inherit the desolate heritages,

		% ---
		\paragraph{1 Nephi 21:9} % *1NEPHI-21-9 
		% ---
			\label{1NEPHI-21-9}

			that thou mayest say to the prisoners:
			Go forth;
			to them that sit in darkness:
			Shew yourselves.
			They shall feed in the ways,
			and their pastures shall be in all high places.

		% ---
		\paragraph{1 Nephi 21:10} % *1NEPHI-21-10 
		% ---
			\label{1NEPHI-21-10}

			They shall not hunger nor thirst,
			neither shall the heat nor the sun smite them.
			For he that hath mercy on them shall lead them;
			even by the springs of water shall he guide them.

		% ---
		\paragraph{1 Nephi 21:11} % *1NEPHI-21-11 
		% ---
			\label{1NEPHI-21-11}

			And I will make all my mountains a way,
			and my highways shall be exalted.

		% ---
		\paragraph{1 Nephi 21:12} % *1NEPHI-21-12 
		% ---
			\label{1NEPHI-21-12}

			And then, O house of Israel,
			behold, these shall come from far;
			and lo, these from the north and from the west,
			and these from the land of Sinim.

		% ---
		\paragraph{1 Nephi 21:13} % *1NEPHI-21-13 
		% ---
			\label{1NEPHI-21-13}

			Sing, O heavens, and be joyful, O earth,
			for the feet of them which are in the east shall be established.
			And break forth into singing, O mountains,
			for they shall be smitten no more.
			For the Lord hath comforted his people
			and will have mercy upon his afflicted.

		% ---
		\paragraph{1 Nephi 21:14} % *1NEPHI-21-14 
		% ---
			\label{1NEPHI-21-14}

			But behold, Zion hath said:
			The Lord hath forsaken me,
			and my Lord hath forgotten me.
			But he will shew that he hath not.

		% ---
		\paragraph{1 Nephi 21:15} % *1NEPHI-21-15 
		% ---
			\label{1NEPHI-21-15}

			For can a woman forget her sucking child,
			that she should not have compassion on the son of her womb?
			Yea, they may forget,
			yet will I not forget thee, O house of Israel.

		% ---
		\paragraph{1 Nephi 21:16} % *1NEPHI-21-16 
		% ---
			\label{1NEPHI-21-16}

			Behold, I have graven thee upon the palms of my hands.
			Thy walls are continually before me.

		% ---
		\paragraph{1 Nephi 21:17} % *1NEPHI-21-17 
		% ---
			\label{1NEPHI-21-17}

			Thy children shall make haste against thy destroyers,
			and they that made thee waste shall go forth of thee.

		% ---
		\paragraph{1 Nephi 21:18} % *1NEPHI-21-18 
		% ---
			\label{1NEPHI-21-18}

			Lift up thine eyes round about and behold,
			all these gather themselves together
			and they shall come to thee.
			And as I live, saith the Lord,
			thou shalt surely clothe thee with them all as with an ornament
			and bind them on, even as a bride.

		% ---
		\paragraph{1 Nephi 21:19} % *1NEPHI-21-19 
		% ---
			\label{1NEPHI-21-19}

			For thy waste and thy desolate places and the land of thy destruction
			shall even now be too narrow by reason of the inhabitants.
			And they that swallowed thee up shall be far away.

		% ---
		\paragraph{1 Nephi 21:20} % *1NEPHI-21-20 
		% ---
			\label{1NEPHI-21-20}

			The children which thou shalt have after thou hast lost the other
			shall say again in thine ears:
			The place is too strait for me;
			give place to me that I may dwell.

		% ---
		\paragraph{1 Nephi 21:21} % *1NEPHI-21-21 
		% ---
			\label{1NEPHI-21-21}

			Then shalt thou say in thine heart:
			Who hath begotten me these,
			seeing I have lost my children
			and am desolate, a captive,
			and removing to and fro?
			And who hath brought up these?
			Behold, I was left alone.
			These, where have they been?

		% ---
		\paragraph{1 Nephi 21:22} % *1NEPHI-21-22 
		% ---
			\label{1NEPHI-21-22}

			Thus saith the Lord God:
			Behold, I will lift up mine hand to the Gentiles
			and set up my standard to the people.
			And they shall bring thy sons in their arms
			and thy daughters shall be carried upon their shoulders.

		% ---
		\paragraph{1 Nephi 21:23} % *1NEPHI-21-23 
		% ---
			\label{1NEPHI-21-23}

			And kings shall be thy nursing fathers
			and their queens thy nursing mothers.
			They shall bow down to thee with their face towards the earth
			and lick up the dust of thy feet.
			And thou shalt know that I am the Lord,
			for they shall not be ashamed that wait for me.

		% ---
		\paragraph{1 Nephi 21:24} % *1NEPHI-21-24 
		% ---
			\label{1NEPHI-21-24}

			For shall the prey be taken from the mighty
			or the lawful captive delivered?

		% ---
		\paragraph{1 Nephi 21:25} % *1NEPHI-21-25 
		% ---
			\label{1NEPHI-21-25}

			But thus saith the Lord:
			Even the captive of the mighty shall be taken away,
			and the prey of the terrible shall be delivered.
			For I will contend with him that contendeth with thee,
			and I will save thy children.

		% ---
		\paragraph{1 Nephi 21:26} % *1NEPHI-21-26 
		% ---
			\label{1NEPHI-21-26}

			And I will feed them that oppress thee with their own flesh.
			They shall be drunken with their own blood as with sweet wine.
			And all flesh shall know that
			I the Lord am thy Savior and thy Redeemer, the Mighty One of Jacob.

	% -----
	\subsubsection{1 Nephi 22} % *1NEPHI-22 
	% -----
		\label{1NEPHI-22}
		
		% ---
		\paragraph{1 Nephi 22:1} % *1NEPHI-22-1 
		% ---
			\label{1 Nephi-22-1}

			And now it came to pass that after I Nephi—
			after that I had read these things which were engraven upon the plates of brass,
			my brethren came unto me and said unto me:
			What mean these things which ye have read?
			Behold, are they to be understood according to things which are spiritual
			which shall come to pass according to the Spirit and not the flesh?

		% ---
		\paragraph{1 Nephi 22:2} % *1NEPHI-22-2 
		% ---
			\label{1 Nephi-22-2}

			And I Nephi saith unto them:
			Behold, they were made manifest unto the prophet by the voice of the Spirit,
			for by the Spirit are all things made known unto the prophets
			which shall come upon the children of men according to the flesh.

		% ---
		\paragraph{1 Nephi 22:3} % *1NEPHI-22-3 
		% ---
			\label{1 Nephi-22-3}

			Wherefore the things of which I have read
			are things pertaining to things both temporal and spiritual.
			For it appears that the house of Israel sooner or later will be scattered
			upon all the face of the earth and also among all nations.

		% ---
		\paragraph{1 Nephi 22:4} % *1NEPHI-22-4 
		% ---
			\label{1 Nephi-22-4}

			And behold, there are many which are already lost
			from the knowledge of they which are at Jerusalem;
			yea, the more part of all the tribes have been led away,
			and they are scattered to and fro upon the isles of the sea.
			And whither they are none of us knoweth,
			save that we know that they have been led away.

		% ---
		\paragraph{1 Nephi 22:5} % *1NEPHI-22-5 
		% ---
			\label{1 Nephi-22-5}

			And since that they have been led away,
			these things have been prophesied concerning them,
			and also concerning all they
			which shall hereafter be scattered and be confounded
			because of the Holy One of Israel,
			for against him will they harden their hearts.
			Wherefore they shall be scattered among all nations
			and shall be hated by all men.

		% ---
		\paragraph{1 Nephi 22:6} % *1NEPHI-22-6 
		% ---
			\label{1 Nephi-22-6}

			Nevertheless, after that they have been nursed by the Gentiles,
			and the Lord hath lifted up his hand upon the Gentiles
			and set them up for a standard,
			and their children shall be carried in their arms
			and their daughters shall be carried upon their shoulders—
			behold, these things of which are spoken are temporal,
			for thus is the covenants of the Lord with our fathers.
			And it meaneth us in the days to come
			and also all our brethren which are of the house of Israel.

		% ---
		\paragraph{1 Nephi 22:7} % *1NEPHI-22-7 
		% ---
			\label{1 Nephi-22-7}

			And it meaneth that the time cometh that
			after all the house of Israel have been scattered and confounded
			that the Lord God will raise up a mighty nation among the Gentiles,
			yea, even upon the face of this land,
			and by them shall our seed be scattered.

		% ---
		\paragraph{1 Nephi 22:8} % *1NEPHI-22-8 
		% ---
			\label{1 Nephi-22-8}

			And after that our seed is scattered,
			the Lord God will proceed to do a marvelous work among the Gentiles
			which shall be of great worth unto our seed.
			Wherefore it is likened unto the being nursed by the Gentiles
			and being carried in their arms and upon their shoulders.

		% ---
		\paragraph{1 Nephi 22:9} % *1NEPHI-22-9 
		% ---
			\label{1 Nephi-22-9}

			And it shall also be of worth unto the Gentiles
			—and not only unto the Gentiles but unto all the house of Israel—
			unto the making known of the covenants of the Father of heaven
			unto Abraham, saying:
			In thy seed shall all the kindreds of the earth be blessed.

		% ---
		\paragraph{1 Nephi 22:10} % *1NEPHI-22-10 
		% ---
			\label{1 Nephi-22-10}

			And I would, my brethren, that ye should know
			that all the kindreds of the earth cannot be blessed
			unless he shall make bare his arm in the eyes of the nations.

		% ---
		\paragraph{1 Nephi 22:11} % *1NEPHI-22-11 
		% ---
			\label{1 Nephi-22-11}

			Wherefore the Lord God will proceed to make bare his arm
			in the eyes of all the nations,
			in bringing about his covenants and his gospel unto they
			which are of the house of Israel.

		% ---
		\paragraph{1 Nephi 22:12} % *1NEPHI-22-12 
		% ---
			\label{1 Nephi-22-12}

			Wherefore he will bring them again out of captivity,
			and they shall be gathered together to the lands of their first inheritance.
			And they shall be brought out of obscurity and out of darkness,
			and they shall know that
			the Lord is their Savior and their Redeemer, the Mighty One of Israel.

		% ---
		\paragraph{1 Nephi 22:13} % *1NEPHI-22-13 
		% ---
			\label{1 Nephi-22-13}

			And the blood of that great and abominable church,
			which is the whore of all the earth,
			shall turn upon their own heads,
			for they shall war among themselves.
			And the sword of their own hands shall fall upon their own heads,
			and they shall be drunken with their own blood.

		% ---
		\paragraph{1 Nephi 22:14} % *1NEPHI-22-14 
		% ---
			\label{1 Nephi-22-14}

			And every nation which shall war against thee, O house of Israel,
			shall be turned one against another.
			And they shall fall into the pit
			which they digged to ensnare the people of the Lord.
			And all they which fight against Zion shall be destroyed.
			And that great whore which hath perverted the right ways of the Lord,
			yea, that great and abominable church,
			shall tumble to the dust,
			and great shall be the fall of it.

		% ---
		\paragraph{1 Nephi 22:15} % *1NEPHI-22-15 
		% ---
			\label{1 Nephi-22-15}

			For behold, saith the prophet, that the time cometh speedily
			that Satan shall have no more power over the hearts of the children of men.
			For the day soon cometh
			that all the proud and they which do wickedly shall be as stubble.
			And the day cometh that they must be burned.

		% ---
		\paragraph{1 Nephi 22:16} % *1NEPHI-22-16 
		% ---
			\label{1 Nephi-22-16}

			For the time soon cometh that
			the fullness of the wrath of God shall be poured out upon all the children of men,
			for he will not suffer that the wicked shall destroy the righteous.

		% ---
		\paragraph{1 Nephi 22:17} % *1NEPHI-22-17 
		% ---
			\label{1 Nephi-22-17}

			Wherefore he will preserve the righteous by his power,
			even if it so be that the fullness of his wrath must come
			and the righteous be preserved,
			even unto the destruction of their enemies by fire.
			Wherefore the righteous need not fear,
			for thus, saith the prophet, they shall be saved,
			even if it so be as by fire.

		% ---
		\paragraph{1 Nephi 22:18} % *1NEPHI-22-18 
		% ---
			\label{1 Nephi-22-18}

			Behold, my brethren, I say unto you
			that these things must shortly come;
			yea, even blood and fire and vapor of smoke must come,
			and it must needs be upon the face of this earth.
			And it cometh unto men according to the flesh
			if it so be that they will harden their hearts against the Holy One of Israel.

		% ---
		\paragraph{1 Nephi 22:19} % *1NEPHI-22-19 
		% ---
			\label{1 Nephi-22-19}

			For behold, the righteous shall not perish,
			for the time surely must come
			that all they which fight against Zion shall be cut off.

		% ---
		\paragraph{1 Nephi 22:20} % *1NEPHI-22-20 
		% ---
			\label{1 Nephi-22-20}

			And the Lord will surely prepare a way for his people
			unto the fulfilling of the words of Moses,
			which he spake, saying:
			A prophet shall the Lord your God raise up unto you like unto me;
			him shall ye hear in all things whatsoever he shall say unto you.
			And it shall come to pass that all they which will not hear that prophet
			shall be cut off from among the people.

		% ---
		\paragraph{1 Nephi 22:21} % *1NEPHI-22-21 
		% ---
			\label{1 Nephi-22-21}

			And now I Nephi declare unto you that
			this prophet of whom Moses spake was the Holy One of Israel.
			Wherefore he shall execute judgment in righteousness.

		% ---
		\paragraph{1 Nephi 22:22} % *1NEPHI-22-22 
		% ---
			\label{1 Nephi-22-22}

			And the righteous need not fear.
			For it is they which shall not be confounded,
			but it is the kingdom of the devil
			which shall be built up among the children of men,
			which kingdom is established among them which are in the flesh.

		% ---
		\paragraph{1 Nephi 22:23} % *1NEPHI-22-23 
		% ---
			\label{1 Nephi-22-23}

			For the time speedily shall come that all churches
			which are built up to get gain
			and all they which are built up to get power over the flesh
			and they which are built up to become popular in the eyes of the world
			and they which seek the lusts of the flesh and the things of the world
			and to do all manner of iniquity,
			yea, in fine, all they which belong to the kingdom of the devil,
			it is they which need fear and tremble and quake;
			it is they which must be brought low in the dust;
			it is they which must be consumed as stubble.
			And this is according to the words of the prophet.

		% ---
		\paragraph{1 Nephi 22:24} % *1NEPHI-22-24 
		% ---
			\label{1 Nephi-22-24}

			And the time cometh speedily
			that the righteous must be led up as calves of the stall,
			and the Holy One of Israel must reign
			in dominion and might and power and great glory.

		% ---
		\paragraph{1 Nephi 22:25} % *1NEPHI-22-25 
		% ---
			\label{1 Nephi-22-25}

			And he gathereth his children from the four quarters of the earth,
			and he numbereth his sheep and they know him.
			And there shall be one fold and one shepherd;
			and he shall feed his sheep
			and in him they shall find pasture.

		% ---
		\paragraph{1 Nephi 22:26} % *1NEPHI-22-26 
		% ---
			\label{1 Nephi-22-26}

			And because of the righteousness of his people, Satan hath no power;
			wherefore he cannot be loosed for the space of many years,
			for he hath no power over the hearts of the people,
			for they dwell in righteousness
			and the Holy One of Israel reigneth.

		% ---
		\paragraph{1 Nephi 22:27} % *1NEPHI-22-27 
		% ---
			\label{1 Nephi-22-27}

			And now behold, I Nephi say unto you
			that all these things must come according to the flesh.

		% ---
		\paragraph{1 Nephi 22:28} % *1NEPHI-22-28 
		% ---
			\label{1 Nephi-22-28}

			But behold, all nations, kindreds, tongues, and people
			shall dwell safely in the Holy One of Israel
			if it so be that they will repent.

		% ---
		\paragraph{1 Nephi 22:29} % *1NEPHI-22-29 
		% ---
			\label{1 Nephi-22-29}

			And now I Nephi make an end,
			for I durst not speak further as yet concerning these things.

		% ---
		\paragraph{1 Nephi 22:30} % *1NEPHI-22-30 
		% ---
			\label{1 Nephi-22-30}

			Wherefore, my brethren, I would that ye should consider
			that the things which have been written upon the plates of brass are true,
			and they testify that a man must be obedient to the commandments of God.

		% ---
		\paragraph{1 Nephi 22:31} % *1NEPHI-22-31 
		% ---
			\label{1 Nephi-22-31}

			Wherefore ye need not suppose that I and my father are the only ones
			which have testified and also taught them.
			Wherefore if ye shall be obedient to the commandments and endure to the end,
			ye shall be saved at the last day.
			And thus it is.
			Amen.